\printleanrefs
\vspace{2em}
\begin{minted}{lean}

-- MODULE: CGD.AntiSelfDualSector.Decoupling
--------------------------------------------------------------------
/--
This theorem asserts that the topological vacuum action is completely
insensitive to the Anti-Self-Dual gauge field connection, securing
the chiral asymmetry of the geometry.
-/
@[litlib_track "Anti-Self-Dual Matter Decoupling"]
theorem §\phantomsection\label{lean:CGD.AntiSelfDualSector.algebraicAntiSelfDualSectorDecoupling}§CGD.AntiSelfDualSector.algebraicAntiSelfDualSectorDecoupling (pu : §\hyperref[lean:CGD.Axioms.PhysicalUniverse]{\detokenize{PhysicalUniverse}}§)
  (A_R_alt : §\hyperref[lean:CGD.Axioms.Sl2cGaugeField]{\detokenize{Sl2cGaugeField}}§) (x : §\hyperref[lean:CGD.Foundations.SpacetimePoint]{\detokenize{SpacetimePoint}}§) :
  §\hyperref[lean:CGD.Foundations.actionVacuum]{\detokenize{actionVacuum}}§ (fun mu nu => §\hyperref[lean:CGD.Foundations.curvature]{\detokenize{curvature}}§ (fun m p => pu.toUniverse.spin4c_connection m p) mu nu x) =
  §\hyperref[lean:CGD.Foundations.actionVacuum]{\detokenize{actionVacuum}}§ (fun mu nu => §\hyperref[lean:CGD.Foundations.curvature]{\detokenize{curvature}}§ (fun m p => §\hyperref[lean:CGD.Foundations.embedSelfDual]{\detokenize{embedSelfDual}}§ (pu.toUniverse.sd_sector m p) + §\hyperref[lean:CGD.Foundations.embedAntiSelfDual]{\detokenize{embedAntiSelfDual}}§ (A_R_alt m p)) mu nu x)


-- MODULE: CGD.AntiSelfDualSector.SelfInteracting
--------------------------------------------------------------------
/--
This theorem proves that non-commuting SU(2) fields (which correspond to non-trivial matter) natively expand into non-zero topological density traces without needing a background metric.
-/
@[litlib_track "Anti-Self-Dual Matter Trace Constraint"]
theorem §\phantomsection\label{lean:CGD.AntiSelfDualSector.kinematicSIDMTrace}§CGD.AntiSelfDualSector.kinematicSIDMTrace (pu : §\hyperref[lean:CGD.Axioms.PhysicalUniverse]{\detokenize{PhysicalUniverse}}§)
  (x : §\hyperref[lean:CGD.Foundations.SpacetimePoint]{\detokenize{SpacetimePoint}}§) (μ ν : Fin 4)
  (h_anti_self_dual_su2 : ∀ m p, §\hyperref[lean:CGD.Foundations.isSu2]{\detokenize{isSu2}}§ (pu.toUniverse.asd_sector m p).val)
  (h_comm : ((pu.toUniverse.asd_sector μ x).val * (pu.toUniverse.asd_sector ν x).val - (pu.toUniverse.asd_sector ν x).val * (pu.toUniverse.asd_sector μ x).val) ≠ 0) :
  Matrix.trace (((pu.toUniverse.asd_sector μ x).val * (pu.toUniverse.asd_sector ν x).val - (pu.toUniverse.asd_sector ν x).val * (pu.toUniverse.asd_sector μ x).val) *
                ((pu.toUniverse.asd_sector μ x).val * (pu.toUniverse.asd_sector ν x).val - (pu.toUniverse.asd_sector ν x).val * (pu.toUniverse.asd_sector μ x).val)) ≠ 0


-- MODULE: CGD.AntiSelfDualSector.VacuumDegeneracy
--------------------------------------------------------------------
/--
Demonstrates the topological collapse of the Anti-Self-Dual vacuum.
If the ASD sector of the universe remains in the trivial vacuum state (zero gauge field),
its resulting emergent Urbantke metric mathematically degenerates to a zero determinant.
This strictly enforces that macroscopic spacetime can only emerge from non-trivial,
non-Abelian gauge condensates, cementing the geometric chirality of the universe.
-/
@[litlib_track "Topological Collapse of the Anti-Self-Dual Vacuum"]
theorem §\phantomsection\label{lean:CGD.AntiSelfDualSector.kinematicAsdVacuumDegeneracy}§CGD.AntiSelfDualSector.kinematicAsdVacuumDegeneracy (pu : §\hyperref[lean:CGD.Axioms.PhysicalUniverse]{\detokenize{PhysicalUniverse}}§) :
  pu.toUniverse.asd_sector.val = (fun _ _ => (0 : §\hyperref[lean:CGD.Foundations.SL2C]{\detokenize{SL2C}}§)) →
  ∀ x, (§\hyperref[lean:CGD.Gravity.urbantkeMetric]{\detokenize{urbantkeMetric}}§ (fun m n => §\hyperref[lean:CGD.Foundations.curvatureSl2c]{\detokenize{curvatureSl2c}}§ pu.toUniverse.asd_sector m n x)).det = 0


-- MODULE: CGD.Axioms.MacroscopicVolume
--------------------------------------------------------------------
/--
Axiom II: Macroscopic Existence (The "Non-Empty" Condition)
We constrain the physics by demanding that the emergent metric is
non-degenerate over an open topological bulk. By defining the bulk as
an `IsOpen` set, we safely allow the volume to mathematically crash to zero
on its boundaries (the Big Bang) or at topological punctures (derived matter defects),
without destroying the extended macroscopic spacetime.
-/
class §\phantomsection\label{lean:CGD.Axioms.MacroscopicVolume}§CGD.Axioms.MacroscopicVolume (u : §\hyperref[lean:CGD.Axioms.Universe]{\detokenize{Universe}}§) (bulk : Set §\hyperref[lean:CGD.Foundations.SpacetimePoint]{\detokenize{SpacetimePoint}}§) : Prop where
  h_bulk_open : IsOpen bulk
  h_bulk_nonempty : Set.Nonempty bulk
  volume_exists : ∀ x ∈ bulk,
    (§\hyperref[lean:CGD.Gravity.urbantkeMetric]{\detokenize{urbantkeMetric}}§ (fun μ ν => §\hyperref[lean:CGD.Foundations.curvatureSl2c]{\detokenize{curvatureSl2c}}§ u.sd_sector μ ν x)).det ≠ 0


-- MODULE: CGD.Axioms.Ontology
--------------------------------------------------------------------
/--
A smooth SL(2, C) gauge field.
As an abstract Lie algebra, it possesses no native spacetime chirality.
-/
structure §\phantomsection\label{lean:CGD.Axioms.Sl2cGaugeField}§CGD.Axioms.Sl2cGaugeField where
  val : Fin 4 → §\hyperref[lean:CGD.Foundations.SpacetimePoint]{\detokenize{CGD.Foundations.SpacetimePoint}}§ → §\hyperref[lean:CGD.Foundations.SL2C]{\detokenize{SL2C}}§
  is_smooth : ∀ mu i j, ContDiff ℝ ⊤ (fun x => (val mu x).val i j)

/--
The Core Ontology: The Universe is a macroscopic Spin(4, C) connection.
We define it strictly as a 4x4 continuous gauge field constrained to the
6-dimensional Spin(4, C) Lie algebra.
-/
structure §\phantomsection\label{lean:CGD.Axioms.Universe}§CGD.Axioms.Universe where
  /-- The unified 4x4 Chiral Dirac field -/
  val : Fin 4 → §\hyperref[lean:CGD.Foundations.SpacetimePoint]{\detokenize{CGD.Foundations.SpacetimePoint}}§ → §\hyperref[lean:CGD.Foundations.ChiralM]{\detokenize{ChiralM}}§

  /-- The geometric constraint forcing the 16D GL(4,C) field into the 6D Spin(4,C) algebra -/
  is_spin4c : ∀ mu x, val mu x =
    §\hyperref[lean:CGD.Foundations.embedSelfDual]{\detokenize{embedSelfDual}}§ (§\hyperref[lean:CGD.Foundations.chiralProject]{\detokenize{chiralProject}}§ (val mu x)).self_dual +
    §\hyperref[lean:CGD.Foundations.embedAntiSelfDual]{\detokenize{embedAntiSelfDual}}§ (§\hyperref[lean:CGD.Foundations.chiralProject]{\detokenize{chiralProject}}§ (val mu x)).anti_self_dual

  /-- Smoothness conditions on the topological projections -/
  sd_is_smooth : ∀ mu i j, ContDiff ℝ ⊤ (fun x => (§\hyperref[lean:CGD.Foundations.chiralProject]{\detokenize{chiralProject}}§ (val mu x)).self_dual.val i j)
  asd_is_smooth : ∀ mu i j, ContDiff ℝ ⊤ (fun x => (§\hyperref[lean:CGD.Foundations.chiralProject]{\detokenize{chiralProject}}§ (val mu x)).anti_self_dual.val i j)

§\phantomsection\label{lean:CGD.Axioms.Universe.asd_sector}§/-- Derives the anti-self-dual SL(2,C) sector from the unified Spin(4,C) field. -/
noncomputable def Universe.asd_sector (u : §\hyperref[lean:CGD.Axioms.Universe]{\detokenize{Universe}}§) : §\hyperref[lean:CGD.Axioms.Sl2cGaugeField]{\detokenize{Sl2cGaugeField}}§ :=
  ⟨fun mu x => (§\hyperref[lean:CGD.Foundations.chiralProject]{\detokenize{chiralProject}}§ (u.val mu x)).anti_self_dual, u.asd_is_smooth⟩

§\phantomsection\label{lean:CGD.Axioms.Universe.sd_sector}§/-- Derives the self-dual SL(2,C) sector from the unified Spin(4,C) field. -/
noncomputable def Universe.sd_sector (u : §\hyperref[lean:CGD.Axioms.Universe]{\detokenize{Universe}}§) : §\hyperref[lean:CGD.Axioms.Sl2cGaugeField]{\detokenize{Sl2cGaugeField}}§ :=
  ⟨fun mu x => (§\hyperref[lean:CGD.Foundations.chiralProject]{\detokenize{chiralProject}}§ (u.val mu x)).self_dual, u.sd_is_smooth⟩

§\phantomsection\label{lean:CGD.Axioms.Universe.spin4c_connection}§/--
The Spin(4, C) 4x4 Dirac Representation.
This definition evaluates to the raw 4x4 field, but its underlying
algebra guarantees it matches the spacetime chiral sum exactly.
-/
noncomputable def Universe.spin4c_connection (u : §\hyperref[lean:CGD.Axioms.Universe]{\detokenize{Universe}}§) (mu : Fin 4) (x : §\hyperref[lean:CGD.Foundations.SpacetimePoint]{\detokenize{CGD.Foundations.SpacetimePoint}}§) : §\hyperref[lean:CGD.Foundations.ChiralM]{\detokenize{ChiralM}}§ :=
  u.val mu x


-- MODULE: CGD.Axioms.PhysicalUniverse
--------------------------------------------------------------------
/--
The Comprehensive Physical Reality.
Bundles the foundational ontology with the topological bulk and its required physical mixins.
Matter is never inputted as an assumption; it exists entirely as the topological complement
of the macroscopic bulk.
-/
structure §\phantomsection\label{lean:CGD.Axioms.PhysicalUniverse}§CGD.Axioms.PhysicalUniverse where
  toUniverse : §\hyperref[lean:CGD.Axioms.Universe]{\detokenize{Universe}}§
  bulk : Set §\hyperref[lean:CGD.Foundations.SpacetimePoint]{\detokenize{SpacetimePoint}}§
  has_volume : §\hyperref[lean:CGD.Axioms.MacroscopicVolume]{\detokenize{MacroscopicVolume}}§ toUniverse bulk
  has_vacuum : §\hyperref[lean:CGD.Axioms.UnimodularVacuum]{\detokenize{UnimodularVacuum}}§ toUniverse bulk


-- MODULE: CGD.Axioms.UnimodularVacuum
--------------------------------------------------------------------
/--
Axiom III: Unimodular Vacuum Isotropy (Unbroken Gauge Symmetry)
Because the `bulk` is defined as the region where det g ≠ 0, it natively
excludes all topological defects. Therefore, the entire `bulk` is the pure vacuum.

This constraint mathematically enforces that the pure vacuum spacetime must not
spontaneously break its internal SU(2) unimodular symmetry. The only symmetric
tensor invariant under local SO(3) rotations is a multiple of the identity matrix.
This natively bounds the field to generate an Einstein manifold.
-/
class §\phantomsection\label{lean:CGD.Axioms.UnimodularVacuum}§CGD.Axioms.UnimodularVacuum (u : §\hyperref[lean:CGD.Axioms.Universe]{\detokenize{Universe}}§) (bulk : Set §\hyperref[lean:CGD.Foundations.SpacetimePoint]{\detokenize{SpacetimePoint}}§) : Prop where
  vacuum_equipartition : ∀ x ∈ bulk,
    let F_adj := fun m n => §\hyperref[lean:CGD.Gravity.cgdAdjointCurvature]{\detokenize{cgdAdjointCurvature}}§ u m n x;
    let Sigma := ∑ μ : Fin 4, ∑ ν : Fin 4, ∑ ρ : Fin 4, ∑ σ : Fin 4,
      §\hyperref[lean:CGD.Gravity.epsilon4]{\detokenize{epsilon4}}§ μ ν ρ σ • (F_adj μ ν * F_adj ρ σ);
    Sigma = (Matrix.trace Sigma / 3) • (1 : Matrix (Fin 3) (Fin 3) ℂ)


-- MODULE: CGD.Cosmology.BigBang
--------------------------------------------------------------------
/--
A kinematic boundary evaluation. This demonstrates that if a phase region is 
artificially constrained to a fully 4D symmetric SO(4) state, the resulting 
Urbantke metric takes a purely Euclidean form (c • 1).

Rather than representing a dynamic cosmological evolution, this acts as a 
mathematical sanity check: it shows that highly symmetric ground-state toy models 
evaluate to non-singular Euclidean spaces rather than generating mathematical 
singularities in the metric.
-/
@[litlib_track "Big Bang Singularity Resolution (Euclidean Bounce)"]
theorem §\phantomsection\label{lean:CGD.Cosmology.kinematicBigBang}§CGD.Cosmology.kinematicBigBang (pu : §\hyperref[lean:CGD.Axioms.PhysicalUniverse]{\detokenize{PhysicalUniverse}}§) (phaseRegion : Set §\hyperref[lean:CGD.Foundations.SpacetimePoint]{\detokenize{SpacetimePoint}}§)
  (h_subset : phaseRegion ⊆ pu.bulk)
  (h_symm : ∀ x ∈ phaseRegion, §\hyperref[lean:CGD.Cosmology.isFully4DSymmetric]{\detokenize{isFully4DSymmetric}}§ (fun mu nu => §\hyperref[lean:CGD.Foundations.curvatureSl2c]{\detokenize{curvatureSl2c}}§ pu.toUniverse.sd_sector mu nu x)) :
  ∀ x ∈ phaseRegion, ∃ c : Complex, c ≠ 0 ∧ §\hyperref[lean:CGD.Gravity.urbantkeMetric]{\detokenize{urbantkeMetric}}§ (fun mu nu => §\hyperref[lean:CGD.Foundations.curvatureSl2c]{\detokenize{curvatureSl2c}}§ pu.toUniverse.sd_sector mu nu x) = c • 1

/--
In pure connection gravity, 4D spacetime volume geometrically requires time-evolution. A completely static universe (where electric/temporal components of the field strength are zero) cannot sustain a macroscopic 4D metric and topologically collapses, forcing the metric determinant to zero.
-/
@[litlib_track "Static Universe Degeneracy"]
theorem §\phantomsection\label{lean:CGD.Cosmology.kinematicStaticUniverseDegeneracy}§CGD.Cosmology.kinematicStaticUniverseDegeneracy (pu : §\hyperref[lean:CGD.Axioms.PhysicalUniverse]{\detokenize{PhysicalUniverse}}§) :
  §\hyperref[lean:CGD.Cosmology.isStaticUniverse]{\detokenize{isStaticUniverse}}§ pu.toUniverse →
  ∀ x, (§\hyperref[lean:CGD.Gravity.urbantkeMetric]{\detokenize{urbantkeMetric}}§ (fun mu nu => §\hyperref[lean:CGD.Foundations.curvatureSl2c]{\detokenize{curvatureSl2c}}§ pu.toUniverse.sd_sector mu nu x)).det = 0


-- MODULE: CGD.Cosmology.Definitions
--------------------------------------------------------------------
/--
Time Emergence: A mathematically rigorous definition of pre-time SO(4) Euclidean symmetry.
Instead of setting time components to zero, we define the symmetric state as a strictly
self-dual configuration.
-/
def §\phantomsection\label{lean:CGD.Cosmology.isFully4DSymmetric}§CGD.Cosmology.isFully4DSymmetric (F : Fin 4 → Fin 4 → §\hyperref[lean:CGD.Foundations.SL2C]{\detokenize{SL2C}}§) : Prop :=
  ∀ mu nu, F mu nu = ∑ rho : Fin 4, ∑ sigma : Fin 4,
    ((1/2 : ℂ) * §\hyperref[lean:CGD.Gravity.epsilon4]{\detokenize{epsilon4}}§ mu nu rho sigma) • F rho sigma

/-- Defines the geometric relationship of a parity-inverted curvature tensor. -/
def §\phantomsection\label{lean:CGD.Cosmology.isParityInvertedTensor}§CGD.Cosmology.isParityInvertedTensor (F P_F : Fin 4 → Fin 4 → §\hyperref[lean:CGD.Foundations.SL2C]{\detokenize{SL2C}}§) (_ : §\hyperref[lean:CGD.Foundations.SpacetimePoint]{\detokenize{SpacetimePoint}}§) : Prop :=
  (∀ i : Fin 4, i ≠ 0 → (P_F 0 i).val = -(F 0 i).val ∧ (P_F i 0).val = -(F i 0).val) ∧
  (∀ i j : Fin 4, i ≠ 0 → j ≠ 0 → (P_F i j).val = (F i j).val) ∧
  ((P_F 0 0).val = (F 0 0).val)

/--
Identifies a static universe where all electric (temporal) curvature components are identically zero.
-/
def §\phantomsection\label{lean:CGD.Cosmology.isStaticUniverse}§CGD.Cosmology.isStaticUniverse (u : §\hyperref[lean:CGD.Axioms.Universe]{\detokenize{Universe}}§) : Prop :=
  ∀ x j, §\hyperref[lean:CGD.Foundations.curvatureSl2c]{\detokenize{curvatureSl2c}}§ u.sd_sector 0 j x = 0


-- MODULE: CGD.Cosmology.ParityInversion
--------------------------------------------------------------------
/--
Demonstrates that the fully antisymmetric topological density (Pontryagin action) natively obeys geometric CPT symmetry. A passive time-reversal transformation (negating the temporal components of the field strength tensor) strictly inverts the sign of the topological charge.
-/
@[litlib_track "Geometric Parity Inversion"]
theorem §\phantomsection\label{lean:CGD.Cosmology.kinematicParityInversion}§CGD.Cosmology.kinematicParityInversion (pu : §\hyperref[lean:CGD.Axioms.PhysicalUniverse]{\detokenize{PhysicalUniverse}}§) :
  ∀ (x : §\hyperref[lean:CGD.Foundations.SpacetimePoint]{\detokenize{SpacetimePoint}}§) (P_F : Fin 4 → Fin 4 → §\hyperref[lean:CGD.Foundations.SL2C]{\detokenize{SL2C}}§),
    §\hyperref[lean:CGD.Cosmology.isParityInvertedTensor]{\detokenize{isParityInvertedTensor}}§ (fun m n => §\hyperref[lean:CGD.Foundations.curvatureSl2c]{\detokenize{curvatureSl2c}}§ pu.toUniverse.sd_sector m n x) P_F x →
    §\hyperref[lean:CGD.Cosmology.pontryaginDensity]{\detokenize{pontryaginDensity}}§ P_F = - §\hyperref[lean:CGD.Cosmology.pontryaginDensity]{\detokenize{pontryaginDensity}}§ (fun mu nu => §\hyperref[lean:CGD.Foundations.curvatureSl2c]{\detokenize{curvatureSl2c}}§ pu.toUniverse.sd_sector mu nu x)

noncomputable def §\phantomsection\label{lean:CGD.Cosmology.pontryaginDensity}§CGD.Cosmology.pontryaginDensity (F : Fin 4 → Fin 4 → §\hyperref[lean:CGD.Foundations.SL2C]{\detokenize{SL2C}}§) : Complex :=
  ∑ μ, ∑ ν, ∑ ρ, ∑ σ, §\hyperref[lean:CGD.Gravity.epsilon4]{\detokenize{epsilon4}}§ μ ν ρ σ * Matrix.trace ((F μ ν).val * (F ρ σ).val)


-- MODULE: CGD.Cosmology.ScaleBreaking
--------------------------------------------------------------------
/--
The conformal scaling property of the Urbantke metric. A scale transformation of the field strength tensor (F → λ² F) results in a strict λ²⁴ scaling of the emergent metric determinant, demonstrating that the macroscopic metric possesses a conformal weight of 6.
-/
@[litlib_track "Kinematic Classical Scale Breaking"]
theorem §\phantomsection\label{lean:CGD.Cosmology.kinematicClassicalScaleBreaking}§CGD.Cosmology.kinematicClassicalScaleBreaking (F : Fin 4 → Fin 4 → §\hyperref[lean:CGD.Foundations.SL2C]{\detokenize{SL2C}}§) (lambda_scale : ℂ) :
  let F_scaled := fun μ ν => §\hyperref[lean:CGD.Foundations.toSl2c]{\detokenize{toSl2c}}§ (lambda_scale^2 • (F μ ν).val);
  (∀ μ ν, §\hyperref[lean:CGD.Gravity.urbantkeMetric]{\detokenize{urbantkeMetric}}§ F_scaled μ ν = lambda_scale^6 * §\hyperref[lean:CGD.Gravity.urbantkeMetric]{\detokenize{urbantkeMetric}}§ F μ ν) ∧
  (§\hyperref[lean:CGD.Gravity.urbantkeMetric]{\detokenize{urbantkeMetric}}§ F_scaled).det = lambda_scale^24 * (§\hyperref[lean:CGD.Gravity.urbantkeMetric]{\detokenize{urbantkeMetric}}§ F).det


-- MODULE: CGD.Foundations.Action
--------------------------------------------------------------------
noncomputable def §\phantomsection\label{lean:CGD.Foundations.actionAntiSelfDual}§CGD.Foundations.actionAntiSelfDual (F : Fin 4 -> Fin 4 -> §\hyperref[lean:CGD.Foundations.ChiralM]{\detokenize{ChiralM}}§) : Complex :=
  let F_R := fun mu nu => (§\hyperref[lean:CGD.Foundations.chiralProject]{\detokenize{chiralProject}}§ (F mu nu)).anti_self_dual
  (-0.5 : Complex) * (
    ∑ mu : Fin 4, ∑ nu : Fin 4, ∑ rho : Fin 4, ∑ sigma : Fin 4,
      §\hyperref[lean:CGD.Gravity.epsilon4]{\detokenize{CGD.Gravity.epsilon4}}§ mu nu rho sigma * Matrix.trace ((F_R mu nu).val * (F_R rho sigma).val)
  )

noncomputable def §\phantomsection\label{lean:CGD.Foundations.actionVacuum}§CGD.Foundations.actionVacuum (F : Fin 4 -> Fin 4 -> §\hyperref[lean:CGD.Foundations.ChiralM]{\detokenize{ChiralM}}§) : Complex :=
  let F_L := fun mu nu => (§\hyperref[lean:CGD.Foundations.chiralProject]{\detokenize{chiralProject}}§ (F mu nu)).self_dual
  (-0.5 : Complex) * (
    ∑ mu : Fin 4, ∑ nu : Fin 4, ∑ rho : Fin 4, ∑ sigma : Fin 4,
      §\hyperref[lean:CGD.Gravity.epsilon4]{\detokenize{CGD.Gravity.epsilon4}}§ mu nu rho sigma * Matrix.trace ((F_L mu nu).val * (F_L rho sigma).val)
  )

/--
A physically valid universe variation is mathematically constrained:
1. It must be a smooth path through configuration space.
2. It must have uniform compact support (vanish at spatial infinity consistently across the perturbation).
3. The underlying physical Lagrangian density must be Lebesgue Integrable.
4. It must map `ℝ → PhysicalUniverse` to ensure the macroscopic volume and vacuum
   constraints are not violated during the perturbation.
-/
def §\phantomsection\label{lean:CGD.Foundations.isValidPhysicalVariation}§CGD.Foundations.isValidPhysicalVariation (v : ℝ → §\hyperref[lean:CGD.Axioms.PhysicalUniverse]{\detokenize{PhysicalUniverse}}§) : Prop :=
  (∀ mu i j, ContDiff ℝ ⊤ (fun (tx : ℝ × §\hyperref[lean:CGD.Foundations.SpacetimePoint]{\detokenize{CGD.Foundations.SpacetimePoint}}§) => ((v tx.1).toUniverse.sd_sector mu tx.2).val i j)) ∧
  (∀ mu i j, ContDiff ℝ ⊤ (fun (tx : ℝ × §\hyperref[lean:CGD.Foundations.SpacetimePoint]{\detokenize{CGD.Foundations.SpacetimePoint}}§) => ((v tx.1).toUniverse.asd_sector mu tx.2).val i j)) ∧
  (∃ R > 0, ∀ t, ∀ x, (x 0)^2 + (x 1)^2 + (x 2)^2 + (x 3)^2 > R^2 →
    (∀ mu, (v t).toUniverse.sd_sector mu x = (v 0).toUniverse.sd_sector mu x) ∧
    (∀ mu, (v t).toUniverse.asd_sector mu x = (v 0).toUniverse.asd_sector mu x)) ∧
  (∀ t, MeasureTheory.Integrable (fun p => §\hyperref[lean:CGD.Foundations.lagrangianDensity]{\detokenize{lagrangianDensity}}§ (fun mu nu => §\hyperref[lean:CGD.Foundations.curvature]{\detokenize{curvature}}§ (fun m x => (v t).toUniverse.spin4c_connection m x) mu nu p)))

noncomputable def §\phantomsection\label{lean:CGD.Foundations.lagrangianDensity}§CGD.Foundations.lagrangianDensity (F : Fin 4 -> Fin 4 -> §\hyperref[lean:CGD.Foundations.ChiralM]{\detokenize{ChiralM}}§) : Complex :=
  (-0.5 : Complex) * (
    ∑ mu : Fin 4, ∑ nu : Fin 4, ∑ rho : Fin 4, ∑ sigma : Fin 4,
      §\hyperref[lean:CGD.Gravity.epsilon4]{\detokenize{CGD.Gravity.epsilon4}}§ mu nu rho sigma * Matrix.trace (F mu nu * F rho sigma)
  )


-- MODULE: CGD.Foundations.Bianchi
--------------------------------------------------------------------
/--
The Differential Bianchi Identity.
Derived strictly from fundamental Lie algebra properties and Clairaut's theorem,
simultaneously performing the $d^2=0$ and Jacobi cancellations algebraically.
-/
@[litlib_track "Kinematic Bianchi Identity"]
theorem §\phantomsection\label{lean:CGD.Foundations.kinematicBianchiIdentity}§CGD.Foundations.kinematicBianchiIdentity
  [clairaut : §\hyperref[lean:Litlib.Y1976.rudin1976principles.ClairautTheoremNDimensional]{\detokenize{Litlib.Y1976.rudin1976principles.ClairautTheoremNDimensional}}§]
  (A : §\hyperref[lean:CGD.Axioms.Sl2cGaugeField]{\detokenize{CGD.Axioms.Sl2cGaugeField}}§) (ρ μ ν : Fin 4) (x : §\hyperref[lean:CGD.Foundations.SpacetimePoint]{\detokenize{SpacetimePoint}}§) :
  §\hyperref[lean:CGD.Foundations.covariantDeriv]{\detokenize{covariantDeriv}}§ A ρ μ ν x +
  §\hyperref[lean:CGD.Foundations.covariantDeriv]{\detokenize{covariantDeriv}}§ A μ ν ρ x +
  §\hyperref[lean:CGD.Foundations.covariantDeriv]{\detokenize{covariantDeriv}}§ A ν ρ μ x = 0


-- MODULE: CGD.Foundations.Calculus
--------------------------------------------------------------------
noncomputable def §\phantomsection\label{lean:CGD.Foundations.curvature}§CGD.Foundations.curvature (A : Fin 4 → §\hyperref[lean:CGD.Foundations.SpacetimePoint]{\detokenize{SpacetimePoint}}§ → §\hyperref[lean:CGD.Foundations.ChiralM]{\detokenize{ChiralM}}§) (mu nu : Fin 4) (x : §\hyperref[lean:CGD.Foundations.SpacetimePoint]{\detokenize{SpacetimePoint}}§) : §\hyperref[lean:CGD.Foundations.ChiralM]{\detokenize{ChiralM}}§ :=
  let dA_nu := partialDerivChiral mu (fun p => A nu p) x
  let dA_mu := partialDerivChiral nu (fun p => A mu p) x
  let raw_comm := bracket (A mu x) (A nu x)
  let proj_comm := §\hyperref[lean:CGD.Foundations.embedSelfDual]{\detokenize{embedSelfDual}}§ (§\hyperref[lean:CGD.Foundations.chiralProject]{\detokenize{chiralProject}}§ raw_comm).self_dual + §\hyperref[lean:CGD.Foundations.embedAntiSelfDual]{\detokenize{embedAntiSelfDual}}§ (§\hyperref[lean:CGD.Foundations.chiralProject]{\detokenize{chiralProject}}§ raw_comm).anti_self_dual
  dA_nu - dA_mu + proj_comm

noncomputable def §\phantomsection\label{lean:CGD.Foundations.curvatureSl2c}§CGD.Foundations.curvatureSl2c (A : Fin 4 → §\hyperref[lean:CGD.Foundations.SpacetimePoint]{\detokenize{SpacetimePoint}}§ → §\hyperref[lean:CGD.Foundations.SL2C]{\detokenize{SL2C}}§) (mu nu : Fin 4) (x : §\hyperref[lean:CGD.Foundations.SpacetimePoint]{\detokenize{SpacetimePoint}}§) : §\hyperref[lean:CGD.Foundations.SL2C]{\detokenize{SL2C}}§ :=
  let dA_nu := partialDerivSl2c mu (A nu) x
  let dA_mu := partialDerivSl2c nu (A mu) x
  let comm  := ⁅A mu x, A nu x⁆
  dA_nu - dA_mu + comm


-- MODULE: CGD.Foundations.ChiralDecomposition
--------------------------------------------------------------------
/--
The topological Pontryagin action strictly preserves the chiral split. Because the cross terms vanish orthogonally, the 4D spacetime topology cleanly factorizes into a self-dual topological charge and an anti-self-dual topological charge.
-/
@[litlib_track "Topological Chiral Decomposition"]
theorem §\phantomsection\label{lean:CGD.Foundations.algebraicChiralDecomposition}§CGD.Foundations.algebraicChiralDecomposition (pu : §\hyperref[lean:CGD.Axioms.PhysicalUniverse]{\detokenize{PhysicalUniverse}}§) (x : §\hyperref[lean:CGD.Foundations.SpacetimePoint]{\detokenize{SpacetimePoint}}§) :
  §\hyperref[lean:CGD.Foundations.lagrangianDensity]{\detokenize{lagrangianDensity}}§ (fun mu nu => §\hyperref[lean:CGD.Foundations.curvature]{\detokenize{curvature}}§ (fun m p => pu.toUniverse.spin4c_connection m p) mu nu x) =
  §\hyperref[lean:CGD.Foundations.actionVacuum]{\detokenize{actionVacuum}}§ (fun mu nu => §\hyperref[lean:CGD.Foundations.curvature]{\detokenize{curvature}}§ (fun m p => pu.toUniverse.spin4c_connection m p) mu nu x) +
  §\hyperref[lean:CGD.Foundations.actionAntiSelfDual]{\detokenize{actionAntiSelfDual}}§ (fun mu nu => §\hyperref[lean:CGD.Foundations.curvature]{\detokenize{curvature}}§ (fun m p => pu.toUniverse.spin4c_connection m p) mu nu x)


-- MODULE: CGD.Foundations.GaugeGroup
--------------------------------------------------------------------
structure §\phantomsection\label{lean:CGD.Foundations.ChiralComponents}§CGD.Foundations.ChiralComponents where
  self_dual : §\hyperref[lean:CGD.Foundations.SL2C]{\detokenize{SL2C}}§
  anti_self_dual  : §\hyperref[lean:CGD.Foundations.SL2C]{\detokenize{SL2C}}§

abbrev §\phantomsection\label{lean:CGD.Foundations.ChiralM}§CGD.Foundations.ChiralM := Matrix (Fin 4) (Fin 4) Complex

abbrev §\phantomsection\label{lean:CGD.Foundations.SL2C}§CGD.Foundations.SL2C := ↥(LieAlgebra.SpecialLinear.sl (Fin 2) Complex)

def §\phantomsection\label{lean:CGD.Foundations.SU2Group}§CGD.Foundations.SU2Group := { M : Matrix (Fin 2) (Fin 2) Complex // M * M.conjTranspose = 1 ∧ M.det = 1 }

noncomputable def §\phantomsection\label{lean:CGD.Foundations.chiralProject}§CGD.Foundations.chiralProject (M : §\hyperref[lean:CGD.Foundations.ChiralM]{\detokenize{ChiralM}}§) : §\hyperref[lean:CGD.Foundations.ChiralComponents]{\detokenize{ChiralComponents}}§ :=
  let b := M.submatrix chiralIso chiralIso
  { self_dual := §\hyperref[lean:CGD.Foundations.toSl2c]{\detokenize{toSl2c}}§ (fun i j => b (Sum.inl i) (Sum.inl j)), anti_self_dual := §\hyperref[lean:CGD.Foundations.toSl2c]{\detokenize{toSl2c}}§ (fun i j => b (Sum.inr i) (Sum.inr j)) }

noncomputable def §\phantomsection\label{lean:CGD.Foundations.embedAntiSelfDual}§CGD.Foundations.embedAntiSelfDual (m : §\hyperref[lean:CGD.Foundations.SL2C]{\detokenize{SL2C}}§) : §\hyperref[lean:CGD.Foundations.ChiralM]{\detokenize{ChiralM}}§ := Matrix.of (fun i j => match (chiralIso.symm i), (chiralIso.symm j) with | Sum.inr i', Sum.inr j' => m.val i' j' | _, _ => 0)

noncomputable def §\phantomsection\label{lean:CGD.Foundations.embedSelfDual}§CGD.Foundations.embedSelfDual (m : §\hyperref[lean:CGD.Foundations.SL2C]{\detokenize{SL2C}}§) : §\hyperref[lean:CGD.Foundations.ChiralM]{\detokenize{ChiralM}}§ := Matrix.of (fun i j => match (chiralIso.symm i), (chiralIso.symm j) with | Sum.inl i', Sum.inl j' => m.val i' j' | _, _ => 0)

§\phantomsection\label{lean:CGD.Foundations.instOneSU2Group}§instance : One §\hyperref[lean:CGD.Foundations.SU2Group]{\detokenize{SU2Group}}§

def §\phantomsection\label{lean:CGD.Foundations.isSu2}§CGD.Foundations.isSu2 (M : Matrix (Fin 2) (Fin 2) Complex) : Prop :=
  Matrix.trace M = 0 ∧ M.conjTranspose = -M

noncomputable def §\phantomsection\label{lean:CGD.Foundations.sigmaX}§CGD.Foundations.sigmaX : Matrix (Fin 2) (Fin 2) ℂ := mkMat 0 1 1 0

noncomputable def §\phantomsection\label{lean:CGD.Foundations.sigmaZ}§CGD.Foundations.sigmaZ : Matrix (Fin 2) (Fin 2) ℂ := mkMat 1 0 0 (-1)

noncomputable def §\phantomsection\label{lean:CGD.Foundations.toSl2c}§CGD.Foundations.toSl2c (M : Matrix (Fin 2) (Fin 2) Complex) : §\hyperref[lean:CGD.Foundations.SL2C]{\detokenize{SL2C}}§ :=
  let tr := M.trace
  let M' := M - (tr / 2) • (1 : Matrix (Fin 2) (Fin 2) Complex)
  ⟨M', by rw[mem_sl_iff, Matrix.trace_sub, Matrix.trace_smul, Matrix.trace_one]; rw[Fintype.card_fin]; simp only[Nat.cast_ofNat]; rw[smul_eq_mul]; field_simp; ring⟩


-- MODULE: CGD.Foundations.Lagrangian.Basic
--------------------------------------------------------------------
def §\phantomsection\label{lean:CGD.Foundations.isSpin4cAlgebra}§CGD.Foundations.isSpin4cAlgebra (M : §\hyperref[lean:CGD.Foundations.ChiralM]{\detokenize{ChiralM}}§) : Prop :=
  ∃ (L R : §\hyperref[lean:CGD.Foundations.SL2C]{\detokenize{SL2C}}§), M = §\hyperref[lean:CGD.Foundations.embedSelfDual]{\detokenize{embedSelfDual}}§ L + §\hyperref[lean:CGD.Foundations.embedAntiSelfDual]{\detokenize{embedAntiSelfDual}}§ R


-- MODULE: CGD.Foundations.Lagrangian.Uniqueness
--------------------------------------------------------------------
/--
The fully antisymmetric Pontryagin topological density is mathematically the unique quadratic, gauge-invariant Lagrangian density that can be constructed without a pre-existing background metric.
-/
@[litlib_track "Topological Lagrangian Uniqueness"]
theorem §\phantomsection\label{lean:CGD.Foundations.topologicalLagrangianUniqueness}§CGD.Foundations.topologicalLagrangianUniqueness
  (L : ((Fin 4 → Fin 4 → §\hyperref[lean:CGD.Foundations.ChiralM]{\detokenize{ChiralM}}§) → Complex))
  (h_inv : ∀ F U, (∀ μ ν, §\hyperref[lean:CGD.Foundations.isSpin4cAlgebra]{\detokenize{isSpin4cAlgebra}}§ (F μ ν)) →
    (∀ μ ν, §\hyperref[lean:CGD.Foundations.isSpin4cAlgebra]{\detokenize{isSpin4cAlgebra}}§ (U * F μ ν * U⁻¹)) →
    L (fun μ ν => U * F μ ν * U⁻¹) = L F)
  (h_topological : ∀ Λ : Matrix (Fin 4) (Fin 4) ℂ,
    (∀ μ ν ρ σ, ∑ α : Fin 4, ∑ β : Fin 4, ∑ γ : Fin 4, ∑ δ : Fin 4,
      Λ α μ * Λ β ν * Λ γ ρ * Λ δ σ * §\hyperref[lean:CGD.Gravity.epsilon4]{\detokenize{CGD.Gravity.epsilon4}}§ α β γ δ = Matrix.det Λ * §\hyperref[lean:CGD.Gravity.epsilon4]{\detokenize{CGD.Gravity.epsilon4}}§ μ ν ρ σ) →
    ∀ F, (∀ μ ν, §\hyperref[lean:CGD.Foundations.isSpin4cAlgebra]{\detokenize{isSpin4cAlgebra}}§ (F μ ν)) →
    (∀ μ ν, §\hyperref[lean:CGD.Foundations.isSpin4cAlgebra]{\detokenize{isSpin4cAlgebra}}§ (∑ α : Fin 4, ∑ β : Fin 4, (Λ μ α * Λ ν β) • F α β)) →
    L (fun μ ν => ∑ α : Fin 4, ∑ β : Fin 4, (Λ μ α * Λ ν β) • F α β) = Matrix.det Λ * L F)
  (hLQuadScale : ∀ (c : ℂ) (F : Fin 4 → Fin 4 → §\hyperref[lean:CGD.Foundations.ChiralM]{\detokenize{ChiralM}}§),
    (∀ μ ν, §\hyperref[lean:CGD.Foundations.isSpin4cAlgebra]{\detokenize{isSpin4cAlgebra}}§ (F μ ν)) → L (fun μ ν => c • F μ ν) = c^2 * L F)
  (hLQuadAdd : ∀ (F G : Fin 4 → Fin 4 → §\hyperref[lean:CGD.Foundations.ChiralM]{\detokenize{ChiralM}}§),
    (∀ μ ν, §\hyperref[lean:CGD.Foundations.isSpin4cAlgebra]{\detokenize{isSpin4cAlgebra}}§ (F μ ν)) → (∀ μ ν, §\hyperref[lean:CGD.Foundations.isSpin4cAlgebra]{\detokenize{isSpin4cAlgebra}}§ (G μ ν)) →
    L (fun μ ν => F μ ν + G μ ν) + L (fun μ ν => F μ ν - G μ ν) = 2 * L F + 2 * L G) :
  ∃ c : ℂ, ∀ F, (∀ μ ν, §\hyperref[lean:CGD.Foundations.isSpin4cAlgebra]{\detokenize{isSpin4cAlgebra}}§ (F μ ν)) → L F = c * ∑ μ : Fin 4, ∑ ν : Fin 4, ∑ ρ : Fin 4, ∑ σ : Fin 4, §\hyperref[lean:CGD.Gravity.epsilon4]{\detokenize{CGD.Gravity.epsilon4}}§ μ ν ρ σ * Matrix.trace (F μ ν * F ρ σ)


-- MODULE: CGD.Foundations.Lagrangian.Variation
--------------------------------------------------------------------
/--
Because the action is the fully antisymmetric Pontryagin density, its
functional variation with respect to compactly supported, smooth gauge field
perturbations evaluates identically to zero.

This theorem rigorously demonstrates that Chiral Gauge Dynamics possesses no
classical Euler-Lagrange equations of motion. Consequently, the dynamic
evolution and stress-energy conservation of the macroscopic universe are not
driven by action minimization, but rather are governed entirely by strict
geometric conservation laws (the Bianchi identities) and topological boundary
conditions.
-/
@[litlib_track "Topological Action Variation"]
theorem §\phantomsection\label{lean:CGD.Foundations.topologicalActionVariationZero}§CGD.Foundations.topologicalActionVariationZero
  (v : ℝ → §\hyperref[lean:CGD.Axioms.PhysicalUniverse]{\detokenize{PhysicalUniverse}}§)
  [dt : §\hyperref[lean:Litlib.Y1965.spivak1965calculus.DivergenceTheoremR4Compact]{\detokenize{Litlib.Y1965.spivak1965calculus.DivergenceTheoremR4Compact}}§ (fun x mu => §\hyperref[lean:CGD.Foundations.variationCurrent]{\detokenize{variationCurrent}}§ v 0 mu x)]
  [lr : §\hyperref[lean:Litlib.Y1976.rudin1976principles.LeibnizIntegralRule]{\detokenize{Litlib.Y1976.rudin1976principles.LeibnizIntegralRule}}§ (fun s x => §\hyperref[lean:CGD.Foundations.lagrangianDensity]{\detokenize{lagrangianDensity}}§ (fun mu nu => §\hyperref[lean:CGD.Foundations.curvature]{\detokenize{curvature}}§ (fun m p => (v s).toUniverse.spin4c_connection m p) mu nu x))]
  (h_valid : §\hyperref[lean:CGD.Foundations.isValidPhysicalVariation]{\detokenize{isValidPhysicalVariation}}§ v) :
  deriv (fun t => §\hyperref[lean:CGD.Foundations.physicalUniverseAction]{\detokenize{physicalUniverseAction}}§ (v t)) 0 = 0


-- MODULE: CGD.Foundations.Lagrangian.Variation.Algebra
--------------------------------------------------------------------
noncomputable def §\phantomsection\label{lean:CGD.Foundations.physicalUniverseAction}§CGD.Foundations.physicalUniverseAction (pu : §\hyperref[lean:CGD.Axioms.PhysicalUniverse]{\detokenize{PhysicalUniverse}}§) : ℂ :=
  universeAction pu.toUniverse

/--
The topological Chern-Simons Variation Current.
K^mu = - 2 * ε^{mu nu rho sigma} Tr( δA_nu F_{rho sigma} )
-/
noncomputable def §\phantomsection\label{lean:CGD.Foundations.variationCurrent}§CGD.Foundations.variationCurrent (v : ℝ → §\hyperref[lean:CGD.Axioms.PhysicalUniverse]{\detokenize{PhysicalUniverse}}§) (t : ℝ) (mu : Fin 4) (x : §\hyperref[lean:CGD.Foundations.SpacetimePoint]{\detokenize{SpacetimePoint}}§) : ℂ :=
  let A := fun s m p => (v s).toUniverse.spin4c_connection m p
  let dA_dt := fun m p => deriv (fun s => A s m p) t
  let F := §\hyperref[lean:CGD.Foundations.curvature]{\detokenize{curvature}}§ (A t)
  (-2 : ℂ) * ∑ nu : Fin 4, ∑ rho : Fin 4, ∑ sigma : Fin 4,
    §\hyperref[lean:CGD.Gravity.epsilon4]{\detokenize{CGD.Gravity.epsilon4}}§ mu nu rho sigma * Matrix.trace (dA_dt nu x * F rho sigma x)


-- MODULE: CGD.Foundations.Spacetime
--------------------------------------------------------------------
/--
The Spacetime Manifold is defined as a map `Fin 4 → ℝ`.
This instantiates the standard `NormedAddCommGroup` and `NormedSpace ℝ` required for rigorous calculus of variations over the manifold.
-/
abbrev §\phantomsection\label{lean:CGD.Foundations.SpacetimePoint}§CGD.Foundations.SpacetimePoint := Fin 4 → ℝ


-- MODULE: CGD.Foundations.TensorCalculus.DifferentialRules
--------------------------------------------------------------------
/-- True geometric definition of the gauge covariant derivative of the field strength tensor:
    D_α F_βγ = ∂_α F_βγ + [A_α, F_βγ] -/
noncomputable def §\phantomsection\label{lean:CGD.Foundations.covariantDeriv}§CGD.Foundations.covariantDeriv (A : Fin 4 → §\hyperref[lean:CGD.Foundations.SpacetimePoint]{\detokenize{SpacetimePoint}}§ → §\hyperref[lean:CGD.Foundations.SL2C]{\detokenize{SL2C}}§) (alpha beta gamma : Fin 4) (x : §\hyperref[lean:CGD.Foundations.SpacetimePoint]{\detokenize{SpacetimePoint}}§) : §\hyperref[lean:CGD.Foundations.SL2C]{\detokenize{SL2C}}§ :=
  let dF := partialDerivSl2c alpha (fun p => §\hyperref[lean:CGD.Foundations.curvatureSl2c]{\detokenize{curvatureSl2c}}§ A beta gamma p) x
  let comm := ⁅A alpha x, §\hyperref[lean:CGD.Foundations.curvatureSl2c]{\detokenize{curvatureSl2c}}§ A beta gamma x⁆
  dF + comm


-- MODULE: CGD.Foundations.Topology
--------------------------------------------------------------------
abbrev §\phantomsection\label{lean:CGD.Foundations.S3}§CGD.Foundations.S3 := { x : Fin 4 → ℝ // (x 0)^2 + (x 1)^2 + (x 2)^2 + (x 3)^2 = 1 }


-- MODULE: CGD.Gravity.DomainSeparation
--------------------------------------------------------------------
/--
A rigorous derivation showing that the scalar volume density μ(x) generates torsion that can be
conformally absorbed, resulting in a physical metric that natively satisfies the Trace-Reversed
Vacuum Einstein Field Equations with a Cosmological Constant (Λ = 1).
-/
@[litlib_track "Macroscopic Cosmological Emergence"]
theorem §\phantomsection\label{lean:CGD.Gravity.macroscopicCosmologicalEmergence}§CGD.Gravity.macroscopicCosmologicalEmergence
  (pu : §\hyperref[lean:CGD.Axioms.PhysicalUniverse]{\detokenize{PhysicalUniverse}}§)
  (urbantke_g : pu.bulk → Fin 4 → Fin 4 → ℝ)
  (urbantke_g_inv : pu.bulk → Fin 4 → Fin 4 → ℝ)
  (urbantke_ricci : pu.bulk → Fin 4 → Fin 4 → ℝ)
  (mu : pu.bulk → ℝ)
  (nabla_mu : (pu.bulk → ℝ) → pu.bulk → Fin 4 → ℝ)
  (nabla_nabla_mu : (pu.bulk → ℝ) → pu.bulk → Fin 4 → Fin 4 → ℝ)
  (g_phys_inv : pu.bulk → Fin 4 → Fin 4 → ℝ)
  (ricci_phys : pu.bulk → Fin 4 → Fin 4 → ℝ)
  (F_ij F_bar_ij : pu.bulk → Fin 3 → Fin 3 → ℂ)
  (plebanski_vacuum : ℂ → (Fin 3 → Fin 3 → ℂ) → (Fin 3 → Fin 3 → ℂ) → Prop)
  (isLeviCivitaRicci : (pu.bulk → (Fin 4 → Fin 4 → ℝ)) → (pu.bulk → (Fin 4 → Fin 4 → ℝ)) → Prop)
  (isLeviCivitaRicciPointwise : (Fin 4 → Fin 4 → ℝ) → (Fin 4 → Fin 4 → ℝ) → Prop)
  [conformal : §\hyperref[lean:Litlib.Y2010.wald2010general.ConformalRicciTransformation]{\detokenize{Litlib.Y2010.wald2010general.ConformalRicciTransformation}}§ 
    pu.bulk urbantke_g (fun p a b => mu p * urbantke_g p a b) urbantke_g_inv 
    urbantke_ricci ricci_phys (fun p => Real.sqrt (mu p)) nabla_mu nabla_nabla_mu isLeviCivitaRicci]
  [eq15 : §\hyperref[lean:Litlib.Y2011.krasnov2011plebanski.Eq15]{\detokenize{Litlib.Y2011.krasnov2011plebanski.Eq15}}§ plebanski_vacuum]
  (pleb_equiv : ∀ p, §\hyperref[lean:Litlib.Y2011.krasnov2011plebanski.PlebanskiToEinsteinEquivalence]{\detokenize{Litlib.Y2011.krasnov2011plebanski.PlebanskiToEinsteinEquivalence}}§ 
    (fun a b => mu p * urbantke_g p a b) (g_phys_inv p) (ricci_phys p) 1 (F_ij p) (F_bar_ij p) plebanski_vacuum isLeviCivitaRicciPointwise)
  (D : (pu.bulk → ℂ) → pu.bulk → ℂ)
  (F_curv : pu.bulk → ℂ)
  (Sigma_urb : pu.bulk → ℂ)
  (h_axiom3 : ∀ p, F_curv p = (mu p : ℂ) * Sigma_urb p)
  (h_bianchi : ∀ p, D F_curv p = 0)
  (h_levi_civita : (∀ p, D (fun x => (mu x : ℂ) * Sigma_urb x) p = 0) → ∀ p, isLeviCivitaRicciPointwise (fun a b => mu p * urbantke_g p a b) (ricci_phys p))
  (h_inv_phys : ∀ p, ∀ a c, (∑ b : Fin 4, (mu p * urbantke_g p a b) * g_phys_inv p b c) = if a = c then 1 else 0)
  (h_pleb_vac_cond : ∀ p, (∑ i : Fin 3, F_ij p i i) = -1 ∧ (∀ i j, F_bar_ij p i j = 0)) :
  ∀ p a c, urbantke_ricci p a c = (mu p * urbantke_g p a c)
      + 2 * nabla_nabla_mu (fun x => Real.log (Real.sqrt (mu x))) p a c 
      + urbantke_g p a c * (∑ d : Fin 4, ∑ e : Fin 4, urbantke_g_inv p d e * nabla_nabla_mu (fun x => Real.log (Real.sqrt (mu x))) p d e) 
      - 2 * nabla_mu (fun x => Real.log (Real.sqrt (mu x))) p a * nabla_mu (fun x => Real.log (Real.sqrt (mu x))) p c 
      + 2 * urbantke_g p a c * (∑ d : Fin 4, ∑ e : Fin 4, urbantke_g_inv p d e * nabla_mu (fun x => Real.log (Real.sqrt (mu x))) p d * nabla_mu (fun x => Real.log (Real.sqrt (mu x))) p e)

/--
In the Capovilla formulation, the Unimodular constraint is governed by the scalar density μ.
This theorem rigorously applies the literature identity to prove that the fundamental volume element
of the emergent spacetime metric (sqrt_g) squared is strictly equal to the Unimodular multiplier μ squared.
-/
@[litlib_track "Macroscopic Unimodular Vacuum Emergence"]
theorem §\phantomsection\label{lean:CGD.Gravity.macroscopicVacuumEmergence}§CGD.Gravity.macroscopicVacuumEmergence
  (pu : §\hyperref[lean:CGD.Axioms.PhysicalUniverse]{\detokenize{PhysicalUniverse}}§)
  (sqrt_g mu eta : §\hyperref[lean:CGD.Foundations.SpacetimePoint]{\detokenize{SpacetimePoint}}§ → ℂ)
  (Psi_3x3 M_3x3 : §\hyperref[lean:CGD.Foundations.SpacetimePoint]{\detokenize{SpacetimePoint}}§ → Matrix (Fin 3) (Fin 3) ℂ)
  (vol_id : §\hyperref[lean:Litlib.Y1991.capovilla1991pure.Theorem_Volume_Element_Identity]{\detokenize{Theorem_Volume_Element_Identity}}§ §\hyperref[lean:CGD.Foundations.SpacetimePoint]{\detokenize{SpacetimePoint}}§ sqrt_g mu eta Psi_3x3 M_3x3) :
  ∀ x ∈ pu.bulk, (sqrt_g x)^2 = (mu x)^2


-- MODULE: CGD.Gravity.ExactSolutions.Abelian
--------------------------------------------------------------------
/--
Provides an exact analytical solution for an Abelian plane wave satisfying the pure CDJ
constraint equation (F ∧ F = 0).
-/
@[litlib_track "Exact Abelian Macroscopic Solution"]
theorem §\phantomsection\label{lean:CGD.Gravity.dynamicExactAbelianSolution}§CGD.Gravity.dynamicExactAbelianSolution (c : ℂ) (hc : c ≠ 0) :
  ∃ (u : §\hyperref[lean:CGD.Axioms.Universe]{\detokenize{Universe}}§),
    §\hyperref[lean:CGD.Gravity.satisfiesPureCdjConstraint]{\detokenize{CGD.Gravity.satisfiesPureCdjConstraint}}§ (fun p m n => §\hyperref[lean:CGD.Gravity.cgdAdjointCurvature]{\detokenize{cgdAdjointCurvature}}§ u m n p) ∧
    (∀ x, §\hyperref[lean:CGD.Foundations.curvatureSl2c]{\detokenize{curvatureSl2c}}§ u.sd_sector 1 2 x = c • §\hyperref[lean:CGD.Foundations.toSl2c]{\detokenize{toSl2c}}§ §\hyperref[lean:CGD.Foundations.sigmaX]{\detokenize{sigmaX}}§) ∧
    (∃ x, §\hyperref[lean:CGD.Foundations.curvatureSl2c]{\detokenize{curvatureSl2c}}§ u.sd_sector 1 2 x ≠ 0)


-- MODULE: CGD.Gravity.ExactSolutions.AlgebraicForms
--------------------------------------------------------------------
/--
Type O (Isotropic / FLRW Cosmological Vacuum)
Algebraic Form: A_0 = 0, A_j(t) = a(t) σ_j
All principal null directions are degenerate.
The scale factor `a(t)` must be strictly real and non-static to ensure a non-degenerate Lorentzian metric.
Acts as the volume-generating Axial Condensate for all other Petrov types.
-/
def §\phantomsection\label{lean:CGD.Gravity.ExactSolutions.IsTypeOForm}§CGD.Gravity.ExactSolutions.IsTypeOForm (pu : §\hyperref[lean:CGD.Axioms.PhysicalUniverse]{\detokenize{PhysicalUniverse}}§) : Prop :=
  ∃ a : ℝ → ℂ, ∀ x : §\hyperref[lean:CGD.Foundations.SpacetimePoint]{\detokenize{SpacetimePoint}}§,
    (a (x 0)).im = 0 ∧
    (fderiv ℝ a (x 0) 1).im = 0 ∧
    a (x 0) ≠ 0 ∧
    fderiv ℝ a (x 0) 1 ≠ 0 ∧
    pu.toUniverse.sd_sector.val 0 x = 0 ∧
    pu.toUniverse.sd_sector.val 1 x = §\hyperref[lean:CGD.Foundations.toSl2c]{\detokenize{toSl2c}}§ (a (x 0) • sigma1.val) ∧
    pu.toUniverse.sd_sector.val 2 x = §\hyperref[lean:CGD.Foundations.toSl2c]{\detokenize{toSl2c}}§ (a (x 0) • sigma2.val) ∧
    pu.toUniverse.sd_sector.val 3 x = §\hyperref[lean:CGD.Foundations.toSl2c]{\detokenize{toSl2c}}§ (a (x 0) • sigma3.val)


-- MODULE: CGD.Gravity.ExactSolutions.CDJEvaluation
--------------------------------------------------------------------
@[litlib_track "CDJ constraint holds"]
lemma §\phantomsection\label{lean:CGD.Gravity.CDJ_constraint_holds}§CGD.Gravity.CDJ_constraint_holds :
  (∑ μ : Fin 4, ∑ ν : Fin 4, ∑ ρ : Fin 4, ∑ σ : Fin 4, §\hyperref[lean:CGD.Gravity.epsilon4]{\detokenize{epsilon4}}§ μ ν ρ σ • (§\hyperref[lean:CGD.Gravity.adj_F]{\detokenize{adj_F}}§ μ ν * §\hyperref[lean:CGD.Gravity.adj_F]{\detokenize{adj_F}}§ ρ σ)) =
  ((∑ μ : Fin 4, ∑ ν : Fin 4, ∑ ρ : Fin 4, ∑ σ : Fin 4, §\hyperref[lean:CGD.Gravity.epsilon4]{\detokenize{epsilon4}}§ μ ν ρ σ • (§\hyperref[lean:CGD.Gravity.adj_F]{\detokenize{adj_F}}§ μ ν * §\hyperref[lean:CGD.Gravity.adj_F]{\detokenize{adj_F}}§ ρ σ)).trace / 3) • 1


-- MODULE: CGD.Gravity.ExactSolutions.Definitions
--------------------------------------------------------------------
/--
A Physical Universe natively satisfies the kinematic reality conditions if and only if
the macroscopic Urbantke metric constructed from its self-dual curvature is strictly
real, strictly non-degenerate, and strictly Lorentzian (signature -+++ or +---)
everywhere within the topological bulk.

This mathematically enforces the historical Reality Conditions without requiring
any non-polynomial differential constraints.
-/
def §\phantomsection\label{lean:CGD.Gravity.ExactSolutions.SatisfiesRealityConditions}§CGD.Gravity.ExactSolutions.SatisfiesRealityConditions (pu : §\hyperref[lean:CGD.Axioms.PhysicalUniverse]{\detokenize{PhysicalUniverse}}§) : Prop :=
  ∀ x ∈ pu.bulk, §\hyperref[lean:CGD.Gravity.isLorentzian]{\detokenize{isLorentzian}}§ (§\hyperref[lean:CGD.Gravity.urbantkeMetric]{\detokenize{urbantkeMetric}}§ (fun μ ν => §\hyperref[lean:CGD.Foundations.curvatureSl2c]{\detokenize{curvatureSl2c}}§ pu.toUniverse.sd_sector.val μ ν x))


-- MODULE: CGD.Gravity.ExactSolutions.FLRW
--------------------------------------------------------------------
/--
LRW (Isotropic) spacetimes intrinsically satisfy the non-degenerate Lorentzian Reality Conditions natively from the SU(2) topology, requiring no exact scalar field assumptions.
-/
@[litlib_track "FLRW connections produce real Lorentzian metrics"]
theorem §\phantomsection\label{lean:CGD.Gravity.ExactSolutions.flrw_satisfies_reality}§CGD.Gravity.ExactSolutions.flrw_satisfies_reality (pu : §\hyperref[lean:CGD.Axioms.PhysicalUniverse]{\detokenize{PhysicalUniverse}}§) (h_typeO : §\hyperref[lean:CGD.Gravity.ExactSolutions.IsTypeOForm]{\detokenize{IsTypeOForm}}§ pu) :
  §\hyperref[lean:CGD.Gravity.ExactSolutions.SatisfiesRealityConditions]{\detokenize{SatisfiesRealityConditions}}§ pu


-- MODULE: CGD.Gravity.ExactSolutions.LorentzianAnsatz
--------------------------------------------------------------------
noncomputable def §\phantomsection\label{lean:CGD.Gravity.F_origin_val}§CGD.Gravity.F_origin_val (mu nu : Fin 4) : Matrix (Fin 2) (Fin 2) ℂ :=
  if mu = 0 ∧ nu = 1 then -§\hyperref[lean:CGD.Foundations.sigmaX]{\detokenize{sigmaX}}§
  else if mu = 0 ∧ nu = 2 then -sigmaY
  else if mu = 0 ∧ nu = 3 then -§\hyperref[lean:CGD.Foundations.sigmaZ]{\detokenize{sigmaZ}}§
  else if mu = 1 ∧ nu = 0 then §\hyperref[lean:CGD.Foundations.sigmaX]{\detokenize{sigmaX}}§
  else if mu = 2 ∧ nu = 0 then sigmaY
  else if mu = 3 ∧ nu = 0 then §\hyperref[lean:CGD.Foundations.sigmaZ]{\detokenize{sigmaZ}}§
  else if mu = 1 ∧ nu = 2 then Complex.I • §\hyperref[lean:CGD.Foundations.sigmaZ]{\detokenize{sigmaZ}}§
  else if mu = 2 ∧ nu = 1 then -Complex.I • §\hyperref[lean:CGD.Foundations.sigmaZ]{\detokenize{sigmaZ}}§
  else if mu = 1 ∧ nu = 3 then -Complex.I • sigmaY
  else if mu = 3 ∧ nu = 1 then Complex.I • sigmaY
  else if mu = 2 ∧ nu = 3 then Complex.I • §\hyperref[lean:CGD.Foundations.sigmaX]{\detokenize{sigmaX}}§
  else if mu = 3 ∧ nu = 2 then -Complex.I • §\hyperref[lean:CGD.Foundations.sigmaX]{\detokenize{sigmaX}}§
  else 0

noncomputable def §\phantomsection\label{lean:CGD.Gravity.adj_F}§CGD.Gravity.adj_F (mu nu : Fin 4) : Matrix (Fin 3) (Fin 3) ℂ :=
  extractAdjoint (§\hyperref[lean:CGD.Gravity.F_origin_val]{\detokenize{F_origin_val}}§ mu nu)

@[litlib_track "Linear-In-Coordinates Lorentzian Matrix Ansatz"]
noncomputable def §\phantomsection\label{lean:CGD.Gravity.exactLorentzianL}§CGD.Gravity.exactLorentzianL (mu : Fin 4) (x : §\hyperref[lean:CGD.Foundations.SpacetimePoint]{\detokenize{SpacetimePoint}}§) : Matrix (Fin 2) (Fin 2) ℂ :=
  if mu = 1 then - (x 0 : ℝ) • §\hyperref[lean:CGD.Foundations.sigmaX]{\detokenize{sigmaX}}§ + (x 3 : ℝ) • ((Complex.I / 2) • sigmaY) - (x 2 : ℝ) • ((Complex.I / 2) • §\hyperref[lean:CGD.Foundations.sigmaZ]{\detokenize{sigmaZ}}§)
  else if mu = 2 then - (x 0 : ℝ) • sigmaY + (x 1 : ℝ) • ((Complex.I / 2) • §\hyperref[lean:CGD.Foundations.sigmaZ]{\detokenize{sigmaZ}}§) - (x 3 : ℝ) • ((Complex.I / 2) • §\hyperref[lean:CGD.Foundations.sigmaX]{\detokenize{sigmaX}}§)
  else if mu = 3 then - (x 0 : ℝ) • §\hyperref[lean:CGD.Foundations.sigmaZ]{\detokenize{sigmaZ}}§ + (x 2 : ℝ) • ((Complex.I / 2) • §\hyperref[lean:CGD.Foundations.sigmaX]{\detokenize{sigmaX}}§) - (x 1 : ℝ) • ((Complex.I / 2) • sigmaY)
  else 0


-- MODULE: CGD.Gravity.ExactSolutions.LorentzianEval
--------------------------------------------------------------------
@[litlib_track "Exact Lorentzian Field Curvature At Origin"]
lemma §\phantomsection\label{lean:CGD.Gravity.F_origin_val_eq_map}§CGD.Gravity.F_origin_val_eq_map (m n : Fin 4) : §\hyperref[lean:CGD.Gravity.F_origin_val]{\detokenize{F_origin_val}}§ m n = §\hyperref[lean:CGD.Gravity.F_val_map]{\detokenize{F_val_map}}§ m n

noncomputable def §\phantomsection\label{lean:CGD.Gravity.F_val_map}§CGD.Gravity.F_val_map (m n : Fin 4) : Matrix (Fin 2) (Fin 2) ℂ :=
  match m, n with
  | 0, 1 => -§\hyperref[lean:CGD.Foundations.sigmaX]{\detokenize{sigmaX}}§ | 0, 2 => -sigmaY | 0, 3 => -§\hyperref[lean:CGD.Foundations.sigmaZ]{\detokenize{sigmaZ}}§
  | 1, 0 => §\hyperref[lean:CGD.Foundations.sigmaX]{\detokenize{sigmaX}}§  | 1, 2 => Complex.I • §\hyperref[lean:CGD.Foundations.sigmaZ]{\detokenize{sigmaZ}}§ | 1, 3 => -Complex.I • sigmaY
  | 2, 0 => sigmaY  | 2, 1 => -Complex.I • §\hyperref[lean:CGD.Foundations.sigmaZ]{\detokenize{sigmaZ}}§ | 2, 3 => Complex.I • §\hyperref[lean:CGD.Foundations.sigmaX]{\detokenize{sigmaX}}§
  | 3, 0 => §\hyperref[lean:CGD.Foundations.sigmaZ]{\detokenize{sigmaZ}}§  | 3, 1 => Complex.I • sigmaY | 3, 2 => -Complex.I • §\hyperref[lean:CGD.Foundations.sigmaX]{\detokenize{sigmaX}}§
  | _, _ => 0


-- MODULE: CGD.Gravity.ExactSolutions.MainTheorem
--------------------------------------------------------------------
/--
This formally constructs an exact analytical non-Abelian $SU(2)$ gauge configuration
that simultaneously satisfies the trace-free Capovilla CDJ constraint at the spatial
origin, and produces a non-degenerate, strictly real, negative-determinant
Lorentzian metric (det g < 0) at x = 0.

This mathematically proves that the foundational vacuum geometry of Chiral Gauge Dynamics
is strictly non-vacuous.
-/
@[litlib_track "Exact Non-Abelian Lorentzian Macroscopic Witness"]
theorem §\phantomsection\label{lean:CGD.Gravity.dynamicExactLorentzianSolution}§CGD.Gravity.dynamicExactLorentzianSolution :
  ∃ (u : §\hyperref[lean:CGD.Axioms.Universe]{\detokenize{Universe}}§) (x : §\hyperref[lean:CGD.Foundations.SpacetimePoint]{\detokenize{SpacetimePoint}}§),
    (∑ μ : Fin 4, ∑ ν : Fin 4, ∑ ρ : Fin 4, ∑ σ : Fin 4,
      §\hyperref[lean:CGD.Gravity.epsilon4]{\detokenize{epsilon4}}§ μ ν ρ σ • (§\hyperref[lean:CGD.Gravity.cgdAdjointCurvature]{\detokenize{cgdAdjointCurvature}}§ u μ ν x * §\hyperref[lean:CGD.Gravity.cgdAdjointCurvature]{\detokenize{cgdAdjointCurvature}}§ u ρ σ x)) =
    ((∑ μ : Fin 4, ∑ ν : Fin 4, ∑ ρ : Fin 4, ∑ σ : Fin 4,
      §\hyperref[lean:CGD.Gravity.epsilon4]{\detokenize{epsilon4}}§ μ ν ρ σ • (§\hyperref[lean:CGD.Gravity.cgdAdjointCurvature]{\detokenize{cgdAdjointCurvature}}§ u μ ν x * §\hyperref[lean:CGD.Gravity.cgdAdjointCurvature]{\detokenize{cgdAdjointCurvature}}§ u ρ σ x)).trace / 3) • 1 ∧
    §\hyperref[lean:CGD.Gravity.isLorentzian]{\detokenize{isLorentzian}}§ (§\hyperref[lean:CGD.Gravity.urbantkeMetric]{\detokenize{urbantkeMetric}}§ (fun m n => §\hyperref[lean:CGD.Foundations.curvatureSl2c]{\detokenize{curvatureSl2c}}§ u.sd_sector m n x))


-- MODULE: CGD.Gravity.ExactSolutions.Soliton
--------------------------------------------------------------------
/--
The Soliton Geometry Theorem.
Proves that a continuous, non-singular Spin(4,C) gauge profile mathematically possesses a
strict non-degenerate Lorentzian emergent metric, while simultaneously satisfying the
Ricci-flat Capovilla CDJ constraint.
-/
@[litlib_track "Exact Analytical Non-Singular Soliton Verification"]
theorem §\phantomsection\label{lean:CGD.Gravity.ExactSolutions.dynamicExactSolitonSolution}§CGD.Gravity.ExactSolutions.dynamicExactSolitonSolution :
  ∃ (u : §\hyperref[lean:CGD.Axioms.Universe]{\detokenize{Universe}}§) (x : §\hyperref[lean:CGD.Foundations.SpacetimePoint]{\detokenize{SpacetimePoint}}§),
    (∑ μ : Fin 4, ∑ ν : Fin 4, ∑ ρ : Fin 4, ∑ σ : Fin 4,
      §\hyperref[lean:CGD.Gravity.epsilon4]{\detokenize{epsilon4}}§ μ ν ρ σ • (§\hyperref[lean:CGD.Gravity.cgdAdjointCurvature]{\detokenize{cgdAdjointCurvature}}§ u μ ν x * §\hyperref[lean:CGD.Gravity.cgdAdjointCurvature]{\detokenize{cgdAdjointCurvature}}§ u ρ σ x)) =
    ((∑ μ : Fin 4, ∑ ν : Fin 4, ∑ ρ : Fin 4, ∑ σ : Fin 4,
      §\hyperref[lean:CGD.Gravity.epsilon4]{\detokenize{epsilon4}}§ μ ν ρ σ • (§\hyperref[lean:CGD.Gravity.cgdAdjointCurvature]{\detokenize{cgdAdjointCurvature}}§ u μ ν x * §\hyperref[lean:CGD.Gravity.cgdAdjointCurvature]{\detokenize{cgdAdjointCurvature}}§ u ρ σ x)).trace / 3) • 1 ∧
    §\hyperref[lean:CGD.Gravity.isLorentzian]{\detokenize{isLorentzian}}§ (§\hyperref[lean:CGD.Gravity.urbantkeMetric]{\detokenize{urbantkeMetric}}§ (fun m n => §\hyperref[lean:CGD.Foundations.curvatureSl2c]{\detokenize{curvatureSl2c}}§ u.sd_sector.val m n x))

@[litlib_track "The exact mathematical gauge field solution"]
noncomputable def §\phantomsection\label{lean:CGD.Gravity.ExactSolutions.solitonAnsatzVal}§CGD.Gravity.ExactSolutions.solitonAnsatzVal (mu : Fin 4) (x : §\hyperref[lean:CGD.Foundations.SpacetimePoint]{\detokenize{SpacetimePoint}}§) : Matrix (Fin 2) (Fin 2) ℂ :=
  let K_c := (§\hyperref[lean:CGD.Gravity.ExactSolutions.solitonK]{\detokenize{solitonK}}§ x : ℂ)
  let t  := (x 0 : ℂ)
  let x1 := (x 1 : ℂ)
  let x2 := (x 2 : ℂ)
  let x3 := (x 3 : ℂ)
  let sX := §\hyperref[lean:CGD.Foundations.sigmaX]{\detokenize{sigmaX}}§
  let sY := solitonSigmaY
  let sZ := §\hyperref[lean:CGD.Foundations.sigmaZ]{\detokenize{sigmaZ}}§
  let isX := Complex.I • sX
  let isY := Complex.I • sY
  let isZ := Complex.I • sZ
  if mu = 1 then -t • sX - (K_c * x3) • isY + (K_c * x2) • isZ
  else if mu = 2 then (K_c * x3) • isX - t • sY - (K_c * x1) • isZ
  else if mu = 3 then -(K_c * x2) • isX + (K_c * x1) • isY - t • sZ
  else 0

@[litlib_track "Smooth scalar profile removing the singularity"]
noncomputable def §\phantomsection\label{lean:CGD.Gravity.ExactSolutions.solitonK}§CGD.Gravity.ExactSolutions.solitonK (x : §\hyperref[lean:CGD.Foundations.SpacetimePoint]{\detokenize{SpacetimePoint}}§) : ℝ :=
  -1 / ((x 1)^2 + (x 2)^2 + (x 3)^2 + 2)


-- MODULE: CGD.Gravity.GeodesicMotion
--------------------------------------------------------------------
/--
Using the rigorous Papapetrou (1951) theorem, we prove that a localized topological defect
(a single-pole singularity in the gauge curvature) must travel along a complex geodesic of the
dynamically emergent Urbantke geometry. This natively utilizes the exact complex covariant
conservation of the emergent Stress-Energy tensor, evaluating strictly inside the macroscopic bulk.
-/
@[litlib_track "Machian Motion of Topological Defects"]
theorem §\phantomsection\label{lean:CGD.Gravity.machianTopologicalDefectMotion}§CGD.Gravity.machianTopologicalDefectMotion
  (pu : §\hyperref[lean:CGD.Axioms.PhysicalUniverse]{\detokenize{PhysicalUniverse}}§)
  (isSmooth : (pu.bulk → ℂ) → Prop)
  (γ : ℂ → pu.bulk)
  (u : ℂ → (Fin 4 → ℂ))
  (du : ℂ → (Fin 4 → ℂ))
  (isSinglePole : (Fin 4 → Fin 4 → pu.bulk → ℂ) → (ℂ → pu.bulk) → Prop)
  [general_bianchi : §\hyperref[lean:Litlib.Y2003.nakahara2003geometry.Theorem_ContractedBianchi]{\detokenize{Litlib.Y2003.nakahara2003geometry.Theorem_ContractedBianchi}}§ pu.bulk (Fin 4) isSmooth (§\hyperref[lean:CGD.Gravity.bulkDeriv]{\detokenize{bulkDeriv}}§ pu)]
  [symm_metric : §\hyperref[lean:Litlib.Y1984.urbantke1984integrability.Eq10_Symmetry]{\detokenize{Litlib.Y1984.urbantke1984integrability.Eq10_Symmetry}}§]
  [papa : §\hyperref[lean:Litlib.Y1951.papapetrou1951spinning.Eq2_12]{\detokenize{Litlib.Y1951.papapetrou1951spinning.Eq2_12}}§ pu.bulk ℂ
    (fun p => Matrix.of (fun m n => §\hyperref[lean:CGD.Gravity.urbantkeMetric]{\detokenize{CGD.Gravity.urbantkeMetric}}§ (fun a b => §\hyperref[lean:CGD.Foundations.curvatureSl2c]{\detokenize{curvatureSl2c}}§ pu.toUniverse.sd_sector a b p.val) m n))
    (fun p => Matrix.of (fun m n => §\hyperref[lean:CGD.Gravity.matrixInv4x4]{\detokenize{CGD.Gravity.matrixInv4x4}}§ (fun a b => §\hyperref[lean:CGD.Gravity.urbantkeMetric]{\detokenize{CGD.Gravity.urbantkeMetric}}§ (fun a b => §\hyperref[lean:CGD.Foundations.curvatureSl2c]{\detokenize{curvatureSl2c}}§ pu.toUniverse.sd_sector a b p.val) a b) m n))
    (fun p rho mu nu => §\hyperref[lean:CGD.Gravity.christoffel]{\detokenize{CGD.Gravity.christoffel}}§ (fun m n p' => §\hyperref[lean:CGD.Gravity.urbantkeMetric]{\detokenize{CGD.Gravity.urbantkeMetric}}§ (fun a b => §\hyperref[lean:CGD.Foundations.curvatureSl2c]{\detokenize{curvatureSl2c}}§ pu.toUniverse.sd_sector a b p') m n) rho mu nu p.val)
    (fun m n p => §\hyperref[lean:CGD.Gravity.emergentStressEnergy]{\detokenize{emergentStressEnergy}}§ (fun a b p' => §\hyperref[lean:CGD.Foundations.curvatureSl2c]{\detokenize{curvatureSl2c}}§ pu.toUniverse.sd_sector a b p') m n p.val)
    (§\hyperref[lean:CGD.Gravity.bulkDeriv]{\detokenize{bulkDeriv}}§ pu) γ u du isSinglePole]
  (h_metric_eq10 : ∀ x : pu.bulk, ∃ (F F_dual : Fin 3 → Fin 4 → Fin 4 → ℂ) (epsilon3 : Fin 3 → Fin 3 → Fin 3 → ℂ),
    (∀ a mu nu, F a mu nu = - F a nu mu) ∧
    (∀ a mu nu, F_dual a mu nu = - F_dual a nu mu) ∧
    (∀ a b c, epsilon3 a b c = - epsilon3 b a c ∧ epsilon3 a b c = - epsilon3 a c b) ∧
    (∀ mu nu, §\hyperref[lean:CGD.Gravity.urbantkeMetric]{\detokenize{CGD.Gravity.urbantkeMetric}}§ (fun a b => §\hyperref[lean:CGD.Foundations.curvatureSl2c]{\detokenize{curvatureSl2c}}§ pu.toUniverse.sd_sector a b x.val) mu nu =
      (-1 / 6 : ℂ) * Finset.sum Finset.univ (fun a => Finset.sum Finset.univ (fun b => Finset.sum Finset.univ (fun c => Finset.sum Finset.univ (fun alpha => Finset.sum Finset.univ (fun beta => epsilon3 a c b * F a mu alpha * F_dual c alpha beta * F b beta nu)))))))
  (h_smooth_g : ∀ i j, isSmooth (fun p => §\hyperref[lean:CGD.Gravity.urbantkeMetric]{\detokenize{CGD.Gravity.urbantkeMetric}}§ (fun a b => §\hyperref[lean:CGD.Foundations.curvatureSl2c]{\detokenize{curvatureSl2c}}§ pu.toUniverse.sd_sector a b p.val) i j))
  (h_smooth_g_inv : ∀ i j, isSmooth (fun p => §\hyperref[lean:CGD.Gravity.matrixInv4x4]{\detokenize{CGD.Gravity.matrixInv4x4}}§ (fun m n => §\hyperref[lean:CGD.Gravity.urbantkeMetric]{\detokenize{CGD.Gravity.urbantkeMetric}}§ (fun a b => §\hyperref[lean:CGD.Foundations.curvatureSl2c]{\detokenize{curvatureSl2c}}§ pu.toUniverse.sd_sector a b p.val) m n) i j))
  (h_smooth_chris : ∀ rho mu nu, isSmooth (fun p => §\hyperref[lean:CGD.Gravity.christoffel]{\detokenize{CGD.Gravity.christoffel}}§ (fun m n p' => §\hyperref[lean:CGD.Gravity.urbantkeMetric]{\detokenize{CGD.Gravity.urbantkeMetric}}§ (fun a b => §\hyperref[lean:CGD.Foundations.curvatureSl2c]{\detokenize{curvatureSl2c}}§ pu.toUniverse.sd_sector a b p') m n) rho mu nu p.val))
  (h_chris_eq : ∀ p : pu.bulk, ∀ rho mu nu, §\hyperref[lean:CGD.Gravity.christoffel]{\detokenize{CGD.Gravity.christoffel}}§ (fun m n p' => §\hyperref[lean:CGD.Gravity.urbantkeMetric]{\detokenize{CGD.Gravity.urbantkeMetric}}§ (fun a b => §\hyperref[lean:CGD.Foundations.curvatureSl2c]{\detokenize{curvatureSl2c}}§ pu.toUniverse.sd_sector a b p') m n) rho mu nu p.val =
    (1/2 : ℂ) * ∑ sigma : Fin 4, §\hyperref[lean:CGD.Gravity.matrixInv4x4]{\detokenize{CGD.Gravity.matrixInv4x4}}§ (fun m n => §\hyperref[lean:CGD.Gravity.urbantkeMetric]{\detokenize{CGD.Gravity.urbantkeMetric}}§ (fun a b => §\hyperref[lean:CGD.Foundations.curvatureSl2c]{\detokenize{curvatureSl2c}}§ pu.toUniverse.sd_sector a b p.val) m n) rho sigma * (
      §\hyperref[lean:CGD.Gravity.bulkDeriv]{\detokenize{bulkDeriv}}§ pu mu (fun p' => §\hyperref[lean:CGD.Gravity.urbantkeMetric]{\detokenize{CGD.Gravity.urbantkeMetric}}§ (fun a b => §\hyperref[lean:CGD.Foundations.curvatureSl2c]{\detokenize{curvatureSl2c}}§ pu.toUniverse.sd_sector a b p'.val) sigma nu) p +
      §\hyperref[lean:CGD.Gravity.bulkDeriv]{\detokenize{bulkDeriv}}§ pu nu (fun p' => §\hyperref[lean:CGD.Gravity.urbantkeMetric]{\detokenize{CGD.Gravity.urbantkeMetric}}§ (fun a b => §\hyperref[lean:CGD.Foundations.curvatureSl2c]{\detokenize{curvatureSl2c}}§ pu.toUniverse.sd_sector a b p'.val) mu sigma) p -
      §\hyperref[lean:CGD.Gravity.bulkDeriv]{\detokenize{bulkDeriv}}§ pu sigma (fun p' => §\hyperref[lean:CGD.Gravity.urbantkeMetric]{\detokenize{CGD.Gravity.urbantkeMetric}}§ (fun a b => §\hyperref[lean:CGD.Foundations.curvatureSl2c]{\detokenize{curvatureSl2c}}§ pu.toUniverse.sd_sector a b p'.val) mu nu) p))
  (h_ricci_eq : ∀ p : pu.bulk, ∀ mu nu, §\hyperref[lean:CGD.Gravity.ricciTensor]{\detokenize{CGD.Gravity.ricciTensor}}§ (fun m n p' => §\hyperref[lean:CGD.Gravity.urbantkeMetric]{\detokenize{CGD.Gravity.urbantkeMetric}}§ (fun a b => §\hyperref[lean:CGD.Foundations.curvatureSl2c]{\detokenize{curvatureSl2c}}§ pu.toUniverse.sd_sector a b p') m n) mu nu p.val =
    ∑ rho : Fin 4, (§\hyperref[lean:CGD.Gravity.bulkDeriv]{\detokenize{bulkDeriv}}§ pu rho (fun p' => §\hyperref[lean:CGD.Gravity.christoffel]{\detokenize{CGD.Gravity.christoffel}}§ (fun m n p'' => §\hyperref[lean:CGD.Gravity.urbantkeMetric]{\detokenize{CGD.Gravity.urbantkeMetric}}§ (fun a b => §\hyperref[lean:CGD.Foundations.curvatureSl2c]{\detokenize{curvatureSl2c}}§ pu.toUniverse.sd_sector a b p'') m n) rho mu nu p'.val) p -
          §\hyperref[lean:CGD.Gravity.bulkDeriv]{\detokenize{bulkDeriv}}§ pu nu (fun p' => §\hyperref[lean:CGD.Gravity.christoffel]{\detokenize{CGD.Gravity.christoffel}}§ (fun m n p'' => §\hyperref[lean:CGD.Gravity.urbantkeMetric]{\detokenize{CGD.Gravity.urbantkeMetric}}§ (fun a b => §\hyperref[lean:CGD.Foundations.curvatureSl2c]{\detokenize{curvatureSl2c}}§ pu.toUniverse.sd_sector a b p'') m n) rho mu rho p'.val) p +
          ∑ lambda : Fin 4, (§\hyperref[lean:CGD.Gravity.christoffel]{\detokenize{CGD.Gravity.christoffel}}§ (fun m n p' => §\hyperref[lean:CGD.Gravity.urbantkeMetric]{\detokenize{CGD.Gravity.urbantkeMetric}}§ (fun a b => §\hyperref[lean:CGD.Foundations.curvatureSl2c]{\detokenize{curvatureSl2c}}§ pu.toUniverse.sd_sector a b p') m n) rho lambda rho p.val * §\hyperref[lean:CGD.Gravity.christoffel]{\detokenize{CGD.Gravity.christoffel}}§ (fun m n p' => §\hyperref[lean:CGD.Gravity.urbantkeMetric]{\detokenize{CGD.Gravity.urbantkeMetric}}§ (fun a b => §\hyperref[lean:CGD.Foundations.curvatureSl2c]{\detokenize{curvatureSl2c}}§ pu.toUniverse.sd_sector a b p') m n) lambda mu nu p.val -
                §\hyperref[lean:CGD.Gravity.christoffel]{\detokenize{CGD.Gravity.christoffel}}§ (fun m n p' => §\hyperref[lean:CGD.Gravity.urbantkeMetric]{\detokenize{CGD.Gravity.urbantkeMetric}}§ (fun a b => §\hyperref[lean:CGD.Foundations.curvatureSl2c]{\detokenize{curvatureSl2c}}§ pu.toUniverse.sd_sector a b p') m n) rho lambda nu p.val * §\hyperref[lean:CGD.Gravity.christoffel]{\detokenize{CGD.Gravity.christoffel}}§ (fun m n p' => §\hyperref[lean:CGD.Gravity.urbantkeMetric]{\detokenize{CGD.Gravity.urbantkeMetric}}§ (fun a b => §\hyperref[lean:CGD.Foundations.curvatureSl2c]{\detokenize{curvatureSl2c}}§ pu.toUniverse.sd_sector a b p') m n) lambda mu rho p.val)))
  (h_single_pole : isSinglePole (fun m n p => §\hyperref[lean:CGD.Gravity.emergentStressEnergy]{\detokenize{emergentStressEnergy}}§ (fun a b p' => §\hyperref[lean:CGD.Foundations.curvatureSl2c]{\detokenize{curvatureSl2c}}§ pu.toUniverse.sd_sector a b p') m n p.val) γ) :
  ∀ s alpha, du s alpha + ∑ mu : Fin 4, ∑ nu : Fin 4, §\hyperref[lean:CGD.Gravity.christoffel]{\detokenize{CGD.Gravity.christoffel}}§ (fun m n p => §\hyperref[lean:CGD.Gravity.urbantkeMetric]{\detokenize{CGD.Gravity.urbantkeMetric}}§ (fun a b => §\hyperref[lean:CGD.Foundations.curvatureSl2c]{\detokenize{curvatureSl2c}}§ pu.toUniverse.sd_sector a b p) m n) alpha mu nu (γ s).val * u s mu * u s nu = 0


-- MODULE: CGD.Gravity.Geometry
--------------------------------------------------------------------
abbrev §\phantomsection\label{lean:CGD.Gravity.SpacetimeIndex}§CGD.Gravity.SpacetimeIndex := Fin 4

noncomputable def §\phantomsection\label{lean:CGD.Gravity.christoffel}§CGD.Gravity.christoffel (g : §\hyperref[lean:CGD.Gravity.SpacetimeIndex]{\detokenize{SpacetimeIndex}}§ → §\hyperref[lean:CGD.Gravity.SpacetimeIndex]{\detokenize{SpacetimeIndex}}§ → §\hyperref[lean:CGD.Foundations.SpacetimePoint]{\detokenize{SpacetimePoint}}§ → Complex) (rho mu nu : §\hyperref[lean:CGD.Gravity.SpacetimeIndex]{\detokenize{SpacetimeIndex}}§) (x : §\hyperref[lean:CGD.Foundations.SpacetimePoint]{\detokenize{SpacetimePoint}}§) : Complex :=
  let g_inv := §\hyperref[lean:CGD.Gravity.matrixInv4x4]{\detokenize{matrixInv4x4}}§ (fun i j => g i j x)
  (1 / 2 : Complex) * ∑ sigma, g_inv rho sigma * (§\hyperref[lean:CGD.Math.partialDeriv]{\detokenize{partialDeriv}}§ mu (fun p => g sigma nu p) x + §\hyperref[lean:CGD.Math.partialDeriv]{\detokenize{partialDeriv}}§ nu (fun p => g mu sigma p) x - §\hyperref[lean:CGD.Math.partialDeriv]{\detokenize{partialDeriv}}§ sigma (fun p => g mu nu p) x)

/-- 4-dimensional Levi-Civita tensor cast to Complex -/
noncomputable def §\phantomsection\label{lean:CGD.Gravity.epsilon4}§CGD.Gravity.epsilon4 (i j k l : Fin 4) : Complex := epsilon4_int i j k l

def §\phantomsection\label{lean:CGD.Gravity.isLorentzian}§CGD.Gravity.isLorentzian (g : Matrix (Fin 4) (Fin 4) Complex) : Prop :=
  (∀ i j, (g i j).im = 0) ∧ (g.det.re < 0) ∧ (g.det.im = 0)

noncomputable def §\phantomsection\label{lean:CGD.Gravity.matrixInv4x4}§CGD.Gravity.matrixInv4x4 (M : Matrix (Fin 4) (Fin 4) Complex) : Matrix (Fin 4) (Fin 4) Complex :=
  (1 / M.det) • M.adjugate

noncomputable def §\phantomsection\label{lean:CGD.Gravity.ricciTensor}§CGD.Gravity.ricciTensor (g : §\hyperref[lean:CGD.Gravity.SpacetimeIndex]{\detokenize{SpacetimeIndex}}§ → §\hyperref[lean:CGD.Gravity.SpacetimeIndex]{\detokenize{SpacetimeIndex}}§ → §\hyperref[lean:CGD.Foundations.SpacetimePoint]{\detokenize{SpacetimePoint}}§ → Complex) (mu nu : §\hyperref[lean:CGD.Gravity.SpacetimeIndex]{\detokenize{SpacetimeIndex}}§) (x : §\hyperref[lean:CGD.Foundations.SpacetimePoint]{\detokenize{SpacetimePoint}}§) : Complex :=
  ∑ rho, (§\hyperref[lean:CGD.Math.partialDeriv]{\detokenize{partialDeriv}}§ rho (fun p => §\hyperref[lean:CGD.Gravity.christoffel]{\detokenize{christoffel}}§ g rho mu nu p) x
        - §\hyperref[lean:CGD.Math.partialDeriv]{\detokenize{partialDeriv}}§ nu (fun p => §\hyperref[lean:CGD.Gravity.christoffel]{\detokenize{christoffel}}§ g rho mu rho p) x
        + ∑ lambda, (§\hyperref[lean:CGD.Gravity.christoffel]{\detokenize{christoffel}}§ g rho lambda rho x * §\hyperref[lean:CGD.Gravity.christoffel]{\detokenize{christoffel}}§ g lambda mu nu x
                   - §\hyperref[lean:CGD.Gravity.christoffel]{\detokenize{christoffel}}§ g rho lambda nu x * §\hyperref[lean:CGD.Gravity.christoffel]{\detokenize{christoffel}}§ g lambda mu rho x))

noncomputable def §\phantomsection\label{lean:CGD.Gravity.urbantkeMetric}§CGD.Gravity.urbantkeMetric (F : Fin 4 -> Fin 4 -> §\hyperref[lean:CGD.Foundations.SL2C]{\detokenize{SL2C}}§) : Matrix (Fin 4) (Fin 4) Complex :=
  fun mu nu =>
    ∑ a : Fin 3, ∑ b : Fin 3, ∑ c : Fin 3,
      let eps_iso := epsilon3 a b c
      let space_term := ∑ alpha : Fin 4, ∑ beta : Fin 4, ∑ gamma : Fin 4, ∑ delta : Fin 4,
          let eps_space := §\hyperref[lean:CGD.Gravity.epsilon4]{\detokenize{epsilon4}}§ alpha beta gamma delta
          let F1 := project F a mu alpha
          let F2 := project F b nu beta
          let F3 := project F c gamma delta
          eps_space * F1 * F2 * F3
      eps_iso * space_term


-- MODULE: CGD.Gravity.MacroscopicVacuum.Basic
--------------------------------------------------------------------
noncomputable def §\phantomsection\label{lean:CGD.Gravity.cgdAdjointCurvature}§CGD.Gravity.cgdAdjointCurvature (u : §\hyperref[lean:CGD.Axioms.Universe]{\detokenize{Universe}}§) (μ ν : Fin 4) (x : §\hyperref[lean:CGD.Foundations.SpacetimePoint]{\detokenize{SpacetimePoint}}§) : Matrix (Fin 3) (Fin 3) ℂ :=
  extractAdjoint (§\hyperref[lean:CGD.Foundations.curvatureSl2c]{\detokenize{curvatureSl2c}}§ u.sd_sector μ ν x).val

/--
The Pure CDJ Vacuum Constraint (Capovilla, Dell, Jacobson 1991).
Enforces that the symmetric tensor Σ^{ab} = ε^{μνρσ} F^a_{μν} F^b_{ρσ} is purely trace-free.
Σ^{ab} - (1/3) δ^{ab} Tr(Σ) = 0
-/
def §\phantomsection\label{lean:CGD.Gravity.satisfiesPureCdjConstraint}§CGD.Gravity.satisfiesPureCdjConstraint (F_adj : §\hyperref[lean:CGD.Foundations.SpacetimePoint]{\detokenize{SpacetimePoint}}§ → Fin 4 → Fin 4 → Matrix (Fin 3) (Fin 3) ℂ) : Prop :=
  ∀ x : §\hyperref[lean:CGD.Foundations.SpacetimePoint]{\detokenize{SpacetimePoint}}§,
    let Sigma := ∑ μ : Fin 4, ∑ ν : Fin 4, ∑ ρ : Fin 4, ∑ σ : Fin 4,
      §\hyperref[lean:CGD.Gravity.epsilon4]{\detokenize{epsilon4}}§ μ ν ρ σ • (F_adj x μ ν * F_adj x ρ σ)
    Sigma = (Matrix.trace Sigma / 3) • (1 : Matrix (Fin 3) (Fin 3) ℂ)


-- MODULE: CGD.Gravity.StressEnergy.Conservation
--------------------------------------------------------------------
/-- Project bulk derivatives into the spacetime coordinate system -/
noncomputable def §\phantomsection\label{lean:CGD.Gravity.bulkDeriv}§CGD.Gravity.bulkDeriv (pu : §\hyperref[lean:CGD.Axioms.PhysicalUniverse]{\detokenize{PhysicalUniverse}}§) (mu : Fin 4) (f : pu.bulk → ℂ) (x : pu.bulk) : ℂ :=
  §\hyperref[lean:CGD.Math.partialDeriv]{\detokenize{partialDeriv}}§ mu (fun p => if h : p ∈ pu.bulk then f ⟨p, h⟩ else 0) x.val

noncomputable def §\phantomsection\label{lean:CGD.Gravity.emergentStressEnergy}§CGD.Gravity.emergentStressEnergy (F : Fin 4 → Fin 4 → §\hyperref[lean:CGD.Foundations.SpacetimePoint]{\detokenize{SpacetimePoint}}§ → §\hyperref[lean:CGD.Foundations.SL2C]{\detokenize{SL2C}}§) (mu nu : Fin 4) (x : §\hyperref[lean:CGD.Foundations.SpacetimePoint]{\detokenize{SpacetimePoint}}§) : ℂ :=
  let g := fun m n p => §\hyperref[lean:CGD.Gravity.urbantkeMetric]{\detokenize{CGD.Gravity.urbantkeMetric}}§ (fun a b => F a b p) m n
  let g_inv := §\hyperref[lean:CGD.Gravity.matrixInv4x4]{\detokenize{CGD.Gravity.matrixInv4x4}}§ (fun m n => g m n x)
  let R_mu_nu := §\hyperref[lean:CGD.Gravity.ricciTensor]{\detokenize{CGD.Gravity.ricciTensor}}§ g mu nu x
  let R_scalar := ∑ alpha : Fin 4, ∑ beta : Fin 4, g_inv alpha beta * §\hyperref[lean:CGD.Gravity.ricciTensor]{\detokenize{CGD.Gravity.ricciTensor}}§ g alpha beta x
  R_mu_nu - (1/2 : ℂ) * g mu nu x * R_scalar

/--
The emergent Stress-Energy tensor (defined as the Einstein tensor of the dynamically emergent Urbantke metric) is covariantly conserved with respect to its own Levi-Civita connection.
This fundamentally restricts the conservation to the macroscopic bulk where det g ≠ 0.
-/
@[litlib_track "Emergent Stress-Energy Conservation"]
theorem §\phantomsection\label{lean:CGD.Gravity.emergentStressEnergyConservation}§CGD.Gravity.emergentStressEnergyConservation
  (pu : §\hyperref[lean:CGD.Axioms.PhysicalUniverse]{\detokenize{PhysicalUniverse}}§)
  (isSmooth : (pu.bulk → ℂ) → Prop)
  [general_bianchi : §\hyperref[lean:Litlib.Y2003.nakahara2003geometry.Theorem_ContractedBianchi]{\detokenize{Litlib.Y2003.nakahara2003geometry.Theorem_ContractedBianchi}}§ pu.bulk (Fin 4) isSmooth (§\hyperref[lean:CGD.Gravity.bulkDeriv]{\detokenize{bulkDeriv}}§ pu)]
  [symm_metric : §\hyperref[lean:Litlib.Y1984.urbantke1984integrability.Eq10_Symmetry]{\detokenize{Litlib.Y1984.urbantke1984integrability.Eq10_Symmetry}}§]
  (h_metric_eq10 : ∀ x : pu.bulk, ∃ (F F_dual : Fin 3 → Fin 4 → Fin 4 → ℂ) (epsilon3 : Fin 3 → Fin 3 → Fin 3 → ℂ),
    (∀ a mu nu, F a mu nu = - F a nu mu) ∧
    (∀ a mu nu, F_dual a mu nu = - F_dual a nu mu) ∧
    (∀ a b c, epsilon3 a b c = - epsilon3 b a c ∧ epsilon3 a b c = - epsilon3 a c b) ∧
    (∀ mu nu, §\hyperref[lean:CGD.Gravity.urbantkeMetric]{\detokenize{CGD.Gravity.urbantkeMetric}}§ (fun a b => §\hyperref[lean:CGD.Foundations.curvatureSl2c]{\detokenize{curvatureSl2c}}§ pu.toUniverse.sd_sector a b x.val) mu nu =
      (-1 / 6 : ℂ) * Finset.sum Finset.univ (fun a => Finset.sum Finset.univ (fun b => Finset.sum Finset.univ (fun c => Finset.sum Finset.univ (fun alpha => Finset.sum Finset.univ (fun beta => epsilon3 a c b * F a mu alpha * F_dual c alpha beta * F b beta nu)))))))
  (h_smooth_g : ∀ i j, isSmooth (fun p => §\hyperref[lean:CGD.Gravity.urbantkeMetric]{\detokenize{CGD.Gravity.urbantkeMetric}}§ (fun a b => §\hyperref[lean:CGD.Foundations.curvatureSl2c]{\detokenize{curvatureSl2c}}§ pu.toUniverse.sd_sector a b p.val) i j))
  (h_smooth_g_inv : ∀ i j, isSmooth (fun p => §\hyperref[lean:CGD.Gravity.matrixInv4x4]{\detokenize{CGD.Gravity.matrixInv4x4}}§ (fun m n => §\hyperref[lean:CGD.Gravity.urbantkeMetric]{\detokenize{CGD.Gravity.urbantkeMetric}}§ (fun a b => §\hyperref[lean:CGD.Foundations.curvatureSl2c]{\detokenize{curvatureSl2c}}§ pu.toUniverse.sd_sector a b p.val) m n) i j))
  (h_smooth_chris : ∀ rho mu nu, isSmooth (fun p => §\hyperref[lean:CGD.Gravity.christoffel]{\detokenize{CGD.Gravity.christoffel}}§ (fun m n p' => §\hyperref[lean:CGD.Gravity.urbantkeMetric]{\detokenize{CGD.Gravity.urbantkeMetric}}§ (fun a b => §\hyperref[lean:CGD.Foundations.curvatureSl2c]{\detokenize{curvatureSl2c}}§ pu.toUniverse.sd_sector a b p') m n) rho mu nu p.val))
  (h_chris_eq : ∀ p : pu.bulk, ∀ rho mu nu, §\hyperref[lean:CGD.Gravity.christoffel]{\detokenize{CGD.Gravity.christoffel}}§ (fun m n p' => §\hyperref[lean:CGD.Gravity.urbantkeMetric]{\detokenize{CGD.Gravity.urbantkeMetric}}§ (fun a b => §\hyperref[lean:CGD.Foundations.curvatureSl2c]{\detokenize{curvatureSl2c}}§ pu.toUniverse.sd_sector a b p') m n) rho mu nu p.val =
    (1/2 : ℂ) * ∑ sigma : Fin 4, §\hyperref[lean:CGD.Gravity.matrixInv4x4]{\detokenize{CGD.Gravity.matrixInv4x4}}§ (fun m n => §\hyperref[lean:CGD.Gravity.urbantkeMetric]{\detokenize{CGD.Gravity.urbantkeMetric}}§ (fun a b => §\hyperref[lean:CGD.Foundations.curvatureSl2c]{\detokenize{curvatureSl2c}}§ pu.toUniverse.sd_sector a b p.val) m n) rho sigma * (
      §\hyperref[lean:CGD.Gravity.bulkDeriv]{\detokenize{bulkDeriv}}§ pu mu (fun p' => §\hyperref[lean:CGD.Gravity.urbantkeMetric]{\detokenize{CGD.Gravity.urbantkeMetric}}§ (fun a b => §\hyperref[lean:CGD.Foundations.curvatureSl2c]{\detokenize{curvatureSl2c}}§ pu.toUniverse.sd_sector a b p'.val) sigma nu) p +
      §\hyperref[lean:CGD.Gravity.bulkDeriv]{\detokenize{bulkDeriv}}§ pu nu (fun p' => §\hyperref[lean:CGD.Gravity.urbantkeMetric]{\detokenize{CGD.Gravity.urbantkeMetric}}§ (fun a b => §\hyperref[lean:CGD.Foundations.curvatureSl2c]{\detokenize{curvatureSl2c}}§ pu.toUniverse.sd_sector a b p'.val) mu sigma) p -
      §\hyperref[lean:CGD.Gravity.bulkDeriv]{\detokenize{bulkDeriv}}§ pu sigma (fun p' => §\hyperref[lean:CGD.Gravity.urbantkeMetric]{\detokenize{CGD.Gravity.urbantkeMetric}}§ (fun a b => §\hyperref[lean:CGD.Foundations.curvatureSl2c]{\detokenize{curvatureSl2c}}§ pu.toUniverse.sd_sector a b p'.val) mu nu) p))
  (h_ricci_eq : ∀ p : pu.bulk, ∀ mu nu, §\hyperref[lean:CGD.Gravity.ricciTensor]{\detokenize{CGD.Gravity.ricciTensor}}§ (fun m n p' => §\hyperref[lean:CGD.Gravity.urbantkeMetric]{\detokenize{CGD.Gravity.urbantkeMetric}}§ (fun a b => §\hyperref[lean:CGD.Foundations.curvatureSl2c]{\detokenize{curvatureSl2c}}§ pu.toUniverse.sd_sector a b p') m n) mu nu p.val =
    ∑ rho : Fin 4, (§\hyperref[lean:CGD.Gravity.bulkDeriv]{\detokenize{bulkDeriv}}§ pu rho (fun p' => §\hyperref[lean:CGD.Gravity.christoffel]{\detokenize{CGD.Gravity.christoffel}}§ (fun m n p'' => §\hyperref[lean:CGD.Gravity.urbantkeMetric]{\detokenize{CGD.Gravity.urbantkeMetric}}§ (fun a b => §\hyperref[lean:CGD.Foundations.curvatureSl2c]{\detokenize{curvatureSl2c}}§ pu.toUniverse.sd_sector a b p'') m n) rho mu nu p'.val) p -
          §\hyperref[lean:CGD.Gravity.bulkDeriv]{\detokenize{bulkDeriv}}§ pu nu (fun p' => §\hyperref[lean:CGD.Gravity.christoffel]{\detokenize{CGD.Gravity.christoffel}}§ (fun m n p'' => §\hyperref[lean:CGD.Gravity.urbantkeMetric]{\detokenize{CGD.Gravity.urbantkeMetric}}§ (fun a b => §\hyperref[lean:CGD.Foundations.curvatureSl2c]{\detokenize{curvatureSl2c}}§ pu.toUniverse.sd_sector a b p'') m n) rho mu rho p'.val) p +
          ∑ lambda : Fin 4, (§\hyperref[lean:CGD.Gravity.christoffel]{\detokenize{CGD.Gravity.christoffel}}§ (fun m n p' => §\hyperref[lean:CGD.Gravity.urbantkeMetric]{\detokenize{CGD.Gravity.urbantkeMetric}}§ (fun a b => §\hyperref[lean:CGD.Foundations.curvatureSl2c]{\detokenize{curvatureSl2c}}§ pu.toUniverse.sd_sector a b p') m n) rho lambda rho p.val * §\hyperref[lean:CGD.Gravity.christoffel]{\detokenize{CGD.Gravity.christoffel}}§ (fun m n p' => §\hyperref[lean:CGD.Gravity.urbantkeMetric]{\detokenize{CGD.Gravity.urbantkeMetric}}§ (fun a b => §\hyperref[lean:CGD.Foundations.curvatureSl2c]{\detokenize{curvatureSl2c}}§ pu.toUniverse.sd_sector a b p') m n) lambda mu nu p.val -
                §\hyperref[lean:CGD.Gravity.christoffel]{\detokenize{CGD.Gravity.christoffel}}§ (fun m n p' => §\hyperref[lean:CGD.Gravity.urbantkeMetric]{\detokenize{CGD.Gravity.urbantkeMetric}}§ (fun a b => §\hyperref[lean:CGD.Foundations.curvatureSl2c]{\detokenize{curvatureSl2c}}§ pu.toUniverse.sd_sector a b p') m n) rho lambda nu p.val * §\hyperref[lean:CGD.Gravity.christoffel]{\detokenize{CGD.Gravity.christoffel}}§ (fun m n p' => §\hyperref[lean:CGD.Gravity.urbantkeMetric]{\detokenize{CGD.Gravity.urbantkeMetric}}§ (fun a b => §\hyperref[lean:CGD.Foundations.curvatureSl2c]{\detokenize{curvatureSl2c}}§ pu.toUniverse.sd_sector a b p') m n) lambda mu rho p.val))) :
  ∀ nu (x : pu.bulk),
    let g_urb := fun m n (p : pu.bulk) => §\hyperref[lean:CGD.Gravity.urbantkeMetric]{\detokenize{CGD.Gravity.urbantkeMetric}}§ (fun a b => §\hyperref[lean:CGD.Foundations.curvatureSl2c]{\detokenize{curvatureSl2c}}§ pu.toUniverse.sd_sector a b p.val) m n
    let g_inv := fun m n (p : pu.bulk) => §\hyperref[lean:CGD.Gravity.matrixInv4x4]{\detokenize{CGD.Gravity.matrixInv4x4}}§ (fun a b => g_urb a b p) m n
    let chris := fun rho mu nu (p : pu.bulk) => §\hyperref[lean:CGD.Gravity.christoffel]{\detokenize{CGD.Gravity.christoffel}}§ (fun m n p' => §\hyperref[lean:CGD.Gravity.urbantkeMetric]{\detokenize{CGD.Gravity.urbantkeMetric}}§ (fun a b => §\hyperref[lean:CGD.Foundations.curvatureSl2c]{\detokenize{curvatureSl2c}}§ pu.toUniverse.sd_sector a b p') m n) rho mu nu p.val
    let T := fun m n (p : pu.bulk) => §\hyperref[lean:CGD.Gravity.emergentStressEnergy]{\detokenize{emergentStressEnergy}}§ (fun a b p' => §\hyperref[lean:CGD.Foundations.curvatureSl2c]{\detokenize{curvatureSl2c}}§ pu.toUniverse.sd_sector a b p') m n p.val
    ∑ mu : Fin 4, ∑ alpha : Fin 4, g_inv mu alpha x * (
      §\hyperref[lean:CGD.Gravity.bulkDeriv]{\detokenize{bulkDeriv}}§ pu alpha (fun p => T mu nu p) x -
      ∑ lambda : Fin 4, (chris lambda alpha mu x * T lambda nu x +
                         chris lambda alpha nu x * T mu lambda x)
    ) = 0

/-- The physical axiom of macroscopic volume strictly implies the bulk subspace is nonempty -/
instance §\phantomsection\label{lean:CGD.Gravity.instPhysicalUniverseBulkNonempty}§CGD.Gravity.instPhysicalUniverseBulkNonempty (pu : §\hyperref[lean:CGD.Axioms.PhysicalUniverse]{\detokenize{PhysicalUniverse}}§) : Nonempty pu.bulk :=
  let ⟨x, hx⟩


-- MODULE: CGD.Math.Calculus
--------------------------------------------------------------------
noncomputable def §\phantomsection\label{lean:CGD.Math.partialDeriv}§CGD.Math.partialDeriv {E : Type*}[NormedAddCommGroup E][NormedSpace ℝ E] (μ : Fin 4) (f : §\hyperref[lean:CGD.Foundations.SpacetimePoint]{\detokenize{SpacetimePoint}}§ → E) : §\hyperref[lean:CGD.Foundations.SpacetimePoint]{\detokenize{SpacetimePoint}}§ → E :=
  fun x => fderiv ℝ f x ((Pi.single μ (1 : ℝ)) : Fin 4 → ℝ)


-- MODULE: CGD.Math.Matrix
--------------------------------------------------------------------
noncomputable instance §\phantomsection\label{lean:CGD.Math.matNormedAddCommGroup}§CGD.Math.matNormedAddCommGroup : NormedAddCommGroup (Matrix (Fin 2) (Fin 2) ℂ) :=
  inferInstanceAs (NormedAddCommGroup (Fin 2 → Fin 2 → ℂ))

noncomputable instance §\phantomsection\label{lean:CGD.Math.matNormedSpaceR}§CGD.Math.matNormedSpaceR : NormedSpace ℝ (Matrix (Fin 2) (Fin 2) ℂ) :=
  inferInstanceAs (NormedSpace ℝ (Fin 2 → Fin 2 → ℂ))


-- MODULE: CGD.Particles.Color
--------------------------------------------------------------------
/--
Defines a gauge field that is constrained to a lower-dimensional Lie subalgebra,
missing at least one of the three internal color generators.
-/
def §\phantomsection\label{lean:CGD.Particles.isColorDeficient}§CGD.Particles.isColorDeficient (F : Fin 4 → Fin 4 → §\hyperref[lean:CGD.Foundations.SL2C]{\detokenize{SL2C}}§) (color : Fin 3) : Prop :=
  ∀ mu nu, project F color mu nu = 0

/--
Demonstrates that if any of the three internal Lie algebra generators (colors) are missing from the gauge field, the scalar triple product natively vanishes, mathematically forcing the macroscopic spacetime volume to zero. A stable spacetime geometry mathematically requires the interaction of exactly all three SU(2) colors.
-/
@[litlib_track "Geometric Degeneracy of Color-Deficient Fields"]
theorem §\phantomsection\label{lean:CGD.Particles.kinematicColorDeficientDegeneracy}§CGD.Particles.kinematicColorDeficientDegeneracy :
  ∀ (F : Fin 4 → Fin 4 → §\hyperref[lean:CGD.Foundations.SL2C]{\detokenize{SL2C}}§) (color : Fin 3),
    §\hyperref[lean:CGD.Particles.isColorDeficient]{\detokenize{isColorDeficient}}§ F color →
    (§\hyperref[lean:CGD.Gravity.urbantkeMetric]{\detokenize{urbantkeMetric}}§ F).det = 0

/--
Demonstrates that a non-zero macroscopic spacetime volume strictly requires non-Abelian fields.
-/
@[litlib_track "Kinematic Multi-Color Requirement"]
theorem §\phantomsection\label{lean:CGD.Particles.kinematicMultiColorRequirement}§CGD.Particles.kinematicMultiColorRequirement :
  ∀ (F : Fin 4 → Fin 4 → §\hyperref[lean:CGD.Foundations.SL2C]{\detokenize{SL2C}}§),
    (§\hyperref[lean:CGD.Gravity.urbantkeMetric]{\detokenize{urbantkeMetric}}§ F).det ≠ 0 →
    ¬ §\hyperref[lean:CGD.Particles.isSingleColor]{\detokenize{isSingleColor}}§ F

/--
Demonstrates that the Urbantke metric determinant fundamentally requires non-commuting Lie algebra generators. For an Abelian (single-color) field, the Lie bracket vanishes, algebraically forcing the macroscopic spacetime volume to zero. Physical spacetime geometries therefore require non-Abelian fields (such as multi-color hadrons) to expand into stable configurations, geometrically manifesting color confinement.
-/
@[litlib_track "Metric Confinement of Abelian Fields"]
theorem §\phantomsection\label{lean:CGD.Particles.kinematicSingleColorDegeneracy}§CGD.Particles.kinematicSingleColorDegeneracy :
  ∀ (F : Fin 4 → Fin 4 → §\hyperref[lean:CGD.Foundations.SL2C]{\detokenize{SL2C}}§),
    §\hyperref[lean:CGD.Particles.isSingleColor]{\detokenize{isSingleColor}}§ F →
    (§\hyperref[lean:CGD.Gravity.urbantkeMetric]{\detokenize{urbantkeMetric}}§ F).det = 0

/--
Demonstrates that a non-zero macroscopic spacetime volume strictly requires non-Abelian fields spanning exactly three active color generators. One or two colors is mathematically insufficient to sustain spacetime volume, natively bounding the minimum unbroken gauge symmetry required for macroscopic existence.
-/
@[litlib_track "Kinematic Three-Color Requirement"]
theorem §\phantomsection\label{lean:CGD.Particles.kinematicThreeColorRequirement}§CGD.Particles.kinematicThreeColorRequirement :
  ∀ (F : Fin 4 → Fin 4 → §\hyperref[lean:CGD.Foundations.SL2C]{\detokenize{SL2C}}§),
    (§\hyperref[lean:CGD.Gravity.urbantkeMetric]{\detokenize{urbantkeMetric}}§ F).det ≠ 0 →
    ¬ §\hyperref[lean:CGD.Particles.isSingleColor]{\detokenize{isSingleColor}}§ F ∧ (∀ color, ¬ §\hyperref[lean:CGD.Particles.isColorDeficient]{\detokenize{isColorDeficient}}§ F color)


-- MODULE: CGD.Particles.DarkMatter
--------------------------------------------------------------------
/--
Dark Matter Decoupling (The "Darkness" Theorem):
A pure Self-Dual defect (Dark Matter) is superposed onto the macroscopic 
vacuum of normal matter. Because the Spin(4,C) Lie algebra perfectly orthogonalizes 
the chiral sectors, the Anti-Self-Dual curvature (which dictates electromagnetism 
and the strong force) evaluates to being mathematically identical to the vacuum state 
without the Dark Matter present. 
-/
@[litlib_track "Dark Matter Anti-Self-Dual Decoupling (Darkness)"]
theorem §\phantomsection\label{lean:CGD.Particles.kinematicDarkMatterDecoupling}§CGD.Particles.kinematicDarkMatterDecoupling (pu : §\hyperref[lean:CGD.Axioms.PhysicalUniverse]{\detokenize{PhysicalUniverse}}§)
  (A_DM : §\hyperref[lean:CGD.Axioms.Sl2cGaugeField]{\detokenize{Sl2cGaugeField}}§) (x : §\hyperref[lean:CGD.Foundations.SpacetimePoint]{\detokenize{SpacetimePoint}}§) (mu nu : Fin 4) :
  let A_tot := fun m p => 
    §\hyperref[lean:CGD.Foundations.embedSelfDual]{\detokenize{embedSelfDual}}§ (A_DM.val m p + pu.toUniverse.sd_sector.val m p) + 
    §\hyperref[lean:CGD.Foundations.embedAntiSelfDual]{\detokenize{embedAntiSelfDual}}§ (pu.toUniverse.asd_sector.val m p);
  (§\hyperref[lean:CGD.Foundations.chiralProject]{\detokenize{chiralProject}}§ (§\hyperref[lean:CGD.Foundations.curvature]{\detokenize{curvature}}§ A_tot mu nu x)).anti_self_dual = 
  (§\hyperref[lean:CGD.Foundations.chiralProject]{\detokenize{chiralProject}}§ (§\hyperref[lean:CGD.Foundations.curvature]{\detokenize{curvature}}§ (fun m p => pu.toUniverse.spin4c_connection m p) mu nu x)).anti_self_dual

/--
Self-Interacting Dark Matter (SIDM):
When two SD-only defects occupy the same spacetime volume, their non-Abelian 
SL(2, C) group structure intrinsically generates a non-linear cross-term in the curvature. 
This strictly forces momentum exchange, proving that pure gravitational defects in CGD 
natively behave as Self-Interacting Dark Matter (SIDM).
-/
@[litlib_track "Dark Matter Self-Interaction (SIDM)"]
theorem §\phantomsection\label{lean:CGD.Particles.kinematicDarkMatterSelfInteraction}§CGD.Particles.kinematicDarkMatterSelfInteraction
  (A_1 A_2 : §\hyperref[lean:CGD.Axioms.Sl2cGaugeField]{\detokenize{Sl2cGaugeField}}§)
  (x : §\hyperref[lean:CGD.Foundations.SpacetimePoint]{\detokenize{SpacetimePoint}}§) (mu nu : Fin 4) :
  let A_tot := fun m p => A_1.val m p + A_2.val m p;
  §\hyperref[lean:CGD.Foundations.curvatureSl2c]{\detokenize{curvatureSl2c}}§ A_tot mu nu x = 
    §\hyperref[lean:CGD.Foundations.curvatureSl2c]{\detokenize{curvatureSl2c}}§ A_1.val mu nu x + 
    §\hyperref[lean:CGD.Foundations.curvatureSl2c]{\detokenize{curvatureSl2c}}§ A_2.val mu nu x + 
    ⁅A_1.val mu x, A_2.val nu x⁆ + ⁅A_2.val mu x, A_1.val nu x⁆


-- MODULE: CGD.Particles.Definitions
--------------------------------------------------------------------
/--
A valid topological projection from the 4D Spacetime gauge field down to a 3D 
spatial boundary. Because the input `Sl2cGaugeField` is natively guaranteed 
to be smooth by the `PhysicalUniverse` axioms, this projection inherently 
guarantees the differentiability and SU(2) bounds of the output matrix field.
-/
structure §\phantomsection\label{lean:CGD.Particles.FermionBoundaryProjection}§CGD.Particles.FermionBoundaryProjection where
  map : §\hyperref[lean:CGD.Axioms.Sl2cGaugeField]{\detokenize{Sl2cGaugeField}}§ → ((Fin 3 → ℝ) → Matrix (Fin 2) (Fin 2) ℂ)
  h_SU2 : ∀ A x, §\hyperref[lean:IsSU2]{\detokenize{IsSU2}}§ (map A x)
  h_smooth : ∀ A i j, Differentiable ℝ (fun x => map A x i j)

/--
Encapsulates the mathematical operation of a 2π adiabatic spatial rotation 
on a topological state, required to evaluate quantum spin-statistics.
-/
structure §\phantomsection\label{lean:CGD.Particles.IsFermionicRotation}§CGD.Particles.IsFermionicRotation
  (path_parity : (ℝ → (Fin 3 → ℝ) → Matrix (Fin 2) (Fin 2) ℂ) → ℕ)
  (U_0 : (Fin 3 → ℝ) → Matrix (Fin 2) (Fin 2) ℂ)
  (U_rot : ℝ → (Fin 3 → ℝ) → Matrix (Fin 2) (Fin 2) ℂ) : Prop where
  h_rot_SU2 : ∀ t x, §\hyperref[lean:IsSU2]{\detokenize{IsSU2}}§ (U_rot t x)
  h_is_2pi_rot : ∀ t x, U_rot t x = U_0 (§\hyperref[lean:rotZ]{\detokenize{rotZ}}§ t x)
  h_fr_parity : path_parity U_rot = 1

/--
Encapsulates the strict mathematical prerequisites for a macroscopic gauge field 
to constitute a topological fermion. By taking a `PhysicalUniverse` and a valid 
projection, it derives its smoothness and SU(2) nature natively from the 
universe's Anti-Self-Dual (Matter) sector, requiring only the assertion of 
the topological boundary conditions (Vacuum at infinity, Baryon Number = 1).
-/
structure §\phantomsection\label{lean:CGD.Particles.IsFermionicState}§CGD.Particles.IsFermionicState
  [MeasureTheory.MeasureSpace (Fin 3 → ℝ)]
  (pu : §\hyperref[lean:CGD.Axioms.PhysicalUniverse]{\detokenize{PhysicalUniverse}}§)
  (proj : §\hyperref[lean:CGD.Particles.FermionBoundaryProjection]{\detokenize{FermionBoundaryProjection}}§) : Prop where
  h_vacuum : §\hyperref[lean:AsymptoticVacuumSU2]{\detokenize{AsymptoticVacuumSU2}}§ (proj.map pu.toUniverse.asd_sector)
  h_integrable : MeasureTheory.Integrable (§\hyperref[lean:BaryonNumberIntegrandSU2]{\detokenize{BaryonNumberIntegrandSU2}}§ (proj.map pu.toUniverse.asd_sector))
  h_degree_one : §\hyperref[lean:baryon_number_SU2]{\detokenize{baryon_number_SU2}}§ 
                   (proj.map pu.toUniverse.asd_sector) 
                   (proj.h_SU2 pu.toUniverse.asd_sector) 
                   (proj.h_smooth pu.toUniverse.asd_sector) 
                   h_vacuum h_integrable = 1

noncomputable def §\phantomsection\label{lean:CGD.Particles.homogeneousChaosAnsatz}§CGD.Particles.homogeneousChaosAnsatz (mu : Fin 4) (x : §\hyperref[lean:CGD.Foundations.SpacetimePoint]{\detokenize{SpacetimePoint}}§) : §\hyperref[lean:CGD.Foundations.SL2C]{\detokenize{SL2C}}§ :=
  if mu = 1 then (Complex.I * (x 1 : ℂ)) • sigma1
  else if mu = 2 then (Complex.I * (x 2 : ℂ)) • sigma2
  else 0

/--
Defines a single-color (Abelian) condensate where all components of the curvature tensor commute, constraining the field to a single U(1) Cartan subalgebra.
-/
def §\phantomsection\label{lean:CGD.Particles.isSingleColor}§CGD.Particles.isSingleColor (F : Fin 4 -> Fin 4 -> §\hyperref[lean:CGD.Foundations.SL2C]{\detokenize{SL2C}}§) : Prop :=
  ∀ mu nu rho sigma, ⁅F mu nu, F rho sigma⁆ = 0

/--
The 4D Topological Soliton (Skyrmion) in the regular gauge.
Topologically localizes the gauge field.
Multiplied by Complex.I to map into the strictly anti-Hermitian su(2) Lie algebra.
Ontologically, this represents the spatial compactification of a 3D Cauchy slice (R3 U {infty} ~= S3)
rather than a 4D Euclidean tunneling event.
-/
noncomputable def §\phantomsection\label{lean:CGD.Particles.topologicalSoliton}§CGD.Particles.topologicalSoliton (mu : Fin 4) (x : §\hyperref[lean:CGD.Foundations.SpacetimePoint]{\detokenize{SpacetimePoint}}§) : §\hyperref[lean:CGD.Foundations.SL2C]{\detokenize{SL2C}}§ :=
  let r2 : ℂ := (x 0 : ℂ)^2 + (x 1 : ℂ)^2 + (x 2 : ℂ)^2 + (x 3 : ℂ)^2
  let D : ℂ := r2 + 1
  if mu = 0 then (Complex.I / D) • ((x 1 : ℂ) • sigma1 + (x 2 : ℂ) • sigma2 + (x 3 : ℂ) • sigma3)
  else if mu = 1 then (Complex.I / D) • (-(x 0 : ℂ) • sigma1 - (x 3 : ℂ) • sigma2 + (x 2 : ℂ) • sigma3)
  else if mu = 2 then (Complex.I / D) • ((x 3 : ℂ) • sigma1 - (x 0 : ℂ) • sigma2 - (x 1 : ℂ) • sigma3)
  else if mu = 3 then (Complex.I / D) • (-(x 2 : ℂ) • sigma1 + (x 1 : ℂ) • sigma2 - (x 0 : ℂ) • sigma3)
  else 0


-- MODULE: CGD.Particles.Electromagnetism
--------------------------------------------------------------------
/--
The true topological Abelian field strength tensor defined as a component
of the full non-Abelian curvature. Incorporates both the curl and the commutator,
enabling the dual divergence to measure the non-trivial topological magnetic current.
F_{μν}^{ij} = ∂_μ A_ν^{ij} - ∂_ν A_μ^{ij} + [A_μ, A_ν]^{ij}
-/
noncomputable def §\phantomsection\label{lean:CGD.Particles.abelianFieldStrength}§CGD.Particles.abelianFieldStrength (pu : §\hyperref[lean:CGD.Axioms.PhysicalUniverse]{\detokenize{PhysicalUniverse}}§) (i j : Fin 4) (μ ν : Fin 4) (x : §\hyperref[lean:CGD.Foundations.SpacetimePoint]{\detokenize{SpacetimePoint}}§) : ℂ :=
  (§\hyperref[lean:CGD.Math.partialDeriv]{\detokenize{partialDeriv}}§ μ (fun p => connectionComponent pu i j ν p) x - §\hyperref[lean:CGD.Math.partialDeriv]{\detokenize{partialDeriv}}§ ν (fun p => connectionComponent pu i j μ p) x) +
  ∑ k : Fin 4, (connectionComponent pu i k μ x * connectionComponent pu k j ν x - connectionComponent pu i k ν x * connectionComponent pu k j μ x)

/--
The emergent U(1) current based on the Duan-Ge (1979) topological decomposition.
Defined as the dual divergence of the Abelian field strength tensor F_ρσ.
J^μ = ε^{μνρσ} ∂_ν F_ρσ
-/
noncomputable def §\phantomsection\label{lean:CGD.Particles.emergentElectricCurrent}§CGD.Particles.emergentElectricCurrent
  (F : Fin 4 → Fin 4 → §\hyperref[lean:CGD.Foundations.SpacetimePoint]{\detokenize{SpacetimePoint}}§ → ℂ)
  (μ : Fin 4) (x : §\hyperref[lean:CGD.Foundations.SpacetimePoint]{\detokenize{SpacetimePoint}}§) : ℂ :=
  ∑ ν : Fin 4, ∑ ρ : Fin 4, ∑ σ : Fin 4,
    §\hyperref[lean:CGD.Gravity.epsilon4]{\detokenize{epsilon4}}§ μ ν ρ σ * §\hyperref[lean:CGD.Math.partialDeriv]{\detokenize{partialDeriv}}§ ν (fun p => F ρ σ p) x

/--
Because the physical Universe is axiomatically a smooth Spin(4,C) connection,
its emergent Abelian projection natively satisfies the n-dimensional Clairaut theorem,
guaranteeing exact topological charge conservation without further assumptions.
-/
@[litlib_track "Kinematic Charge Conservation"]
theorem §\phantomsection\label{lean:CGD.Particles.kinematicChargeConservation}§CGD.Particles.kinematicChargeConservation
  [clairaut : §\hyperref[lean:Litlib.Y1976.rudin1976principles.ClairautTheoremNDimensional]{\detokenize{Litlib.Y1976.rudin1976principles.ClairautTheoremNDimensional}}§]
  (pu : §\hyperref[lean:CGD.Axioms.PhysicalUniverse]{\detokenize{PhysicalUniverse}}§) (i j : Fin 4) :
  ∀ x : §\hyperref[lean:CGD.Foundations.SpacetimePoint]{\detokenize{SpacetimePoint}}§,
    ∑ μ : Fin 4, §\hyperref[lean:CGD.Math.partialDeriv]{\detokenize{partialDeriv}}§ μ (fun p => §\hyperref[lean:CGD.Particles.emergentElectricCurrent]{\detokenize{emergentElectricCurrent}}§ (§\hyperref[lean:CGD.Particles.abelianFieldStrength]{\detokenize{abelianFieldStrength}}§ pu i j) μ p) x = 0

/--
Because the emergent current is purely topological (defined via the Levi-Civita symbol),
its divergence strictly vanishes due to the commutativity of partial derivatives,
justified natively via n-dimensional Clairaut theorem on the Abelian field strength.
-/
@[litlib_track "Topological Charge Conservation"]
theorem §\phantomsection\label{lean:CGD.Particles.topologicalChargeConservation}§CGD.Particles.topologicalChargeConservation
  [clairaut : §\hyperref[lean:Litlib.Y1976.rudin1976principles.ClairautTheoremNDimensional]{\detokenize{Litlib.Y1976.rudin1976principles.ClairautTheoremNDimensional}}§]
  (F : Fin 4 → Fin 4 → §\hyperref[lean:CGD.Foundations.SpacetimePoint]{\detokenize{SpacetimePoint}}§ → ℂ)
  (h_smooth_re : ∀ ρ σ, ContDiffOn ℝ 2 (fun p => (F ρ σ p).re) Set.univ)
  (h_smooth_im : ∀ ρ σ, ContDiffOn ℝ 2 (fun p => (F ρ σ p).im) Set.univ)
  (h_diff_F : ∀ ρ σ x, DifferentiableAt ℝ (fun p => F ρ σ p) x)
  (h_diff : ∀ ν ρ σ x, DifferentiableAt ℝ (fun p => §\hyperref[lean:CGD.Math.partialDeriv]{\detokenize{partialDeriv}}§ ν (fun p' => F ρ σ p') p) x)
  (h_diff_fderiv_re : ∀ ρ σ x, DifferentiableAt ℝ (fderiv ℝ (fun p => (F ρ σ p).re)) x)
  (h_diff_fderiv_im : ∀ ρ σ x, DifferentiableAt ℝ (fderiv ℝ (fun p => (F ρ σ p).im)) x) :
  ∀ x : §\hyperref[lean:CGD.Foundations.SpacetimePoint]{\detokenize{SpacetimePoint}}§,
    ∑ μ : Fin 4, §\hyperref[lean:CGD.Math.partialDeriv]{\detokenize{partialDeriv}}§ μ (fun p => §\hyperref[lean:CGD.Particles.emergentElectricCurrent]{\detokenize{emergentElectricCurrent}}§ F μ p) x = 0


-- MODULE: CGD.Particles.Mass
--------------------------------------------------------------------
/--
The topological rest mass is strictly defined as the absolute value of the Cartan-Maurer 
topological integral evaluated over the asymptotic spatial boundary of the Anti-Self-Dual (Matter) field.

**PHYSICAL ONTOLOGY:**
For a stable Lorentzian particle, `BoundaryManifold` physically represents the spatial compactification 
($S^3$) of a 3D Cauchy slice of spacetime at a fixed time $t$. The particle is therefore modeled strictly 
as a stable Lorentzian Soliton, not a transient Euclidean Instanton.

Because the non-compact boost directions of SL(2, C) form a contractible space ($\mathbb{R}^3$), they possess 
zero topological winding number and strictly drop out of the boundary integral. This guarantees that 
exotic negative-mass boost states are kinematically excluded from the topological invariant, leaving 
only the strictly positive SU(2) rest mass charge.
-/
noncomputable def §\phantomsection\label{lean:CGD.Particles.inertialMass}§CGD.Particles.inertialMass
  {BoundaryManifold : Type*} [TopologicalSpace BoundaryManifold] [Nonempty BoundaryManifold]
  (boundaryMap : (Fin 4 → §\hyperref[lean:CGD.Foundations.SpacetimePoint]{\detokenize{SpacetimePoint}}§ → §\hyperref[lean:CGD.Foundations.SL2C]{\detokenize{SL2C}}§) → BoundaryManifold → §\hyperref[lean:CGD.Foundations.SU2Group]{\detokenize{SU2Group}}§)
  (cartanMaurerIntegral : (BoundaryManifold → §\hyperref[lean:CGD.Foundations.SU2Group]{\detokenize{SU2Group}}§) → ℝ)
  (pu : §\hyperref[lean:CGD.Axioms.PhysicalUniverse]{\detokenize{PhysicalUniverse}}§) : ℝ :=
  |cartanMaurerIntegral (boundaryMap pu.toUniverse.asd_sector.val)|

/-- 
By applying the rigorous Cartan-Maurer topological integral to the SU(2) boundary mapping 
of the matter sector, we prove the topological origin of strictly positive inertial mass. 

Because stable particles must constitute a topological homeomorphism to the gauge boundary 
on a given Cauchy surface, their integrated charge strictly evaluates to a non-zero integer. 
This mathematically verifies that stable topological Solitons in Chiral Gauge Dynamics intrinsically 
possess a non-zero, strictly positive Mass Gap, completely avoiding the vacuum instabilities 
of the non-compact SL(2, C) group.
-/
@[litlib_track "Topological Mass Gap"]
theorem §\phantomsection\label{lean:CGD.Particles.topologicalMassGap}§CGD.Particles.topologicalMassGap
  {BoundaryManifold : Type*} [TopologicalSpace BoundaryManifold] [Nonempty BoundaryManifold]
  (boundaryMap : (Fin 4 → §\hyperref[lean:CGD.Foundations.SpacetimePoint]{\detokenize{SpacetimePoint}}§ → §\hyperref[lean:CGD.Foundations.SL2C]{\detokenize{SL2C}}§) → BoundaryManifold → §\hyperref[lean:CGD.Foundations.SU2Group]{\detokenize{SU2Group}}§)
  (windingNumber : (BoundaryManifold → §\hyperref[lean:CGD.Foundations.SU2Group]{\detokenize{SU2Group}}§) → ℤ)
  (cartanMaurerIntegral : (BoundaryManifold → §\hyperref[lean:CGD.Foundations.SU2Group]{\detokenize{SU2Group}}§) → ℝ)
  [tc : §\hyperref[lean:Litlib.Y2003.nakahara2003geometry.CartanMaurerTopology]{\detokenize{CartanMaurerTopology}}§ (BoundaryManifold → §\hyperref[lean:CGD.Foundations.SU2Group]{\detokenize{SU2Group}}§) Continuous windingNumber cartanMaurerIntegral]
  [belavin : §\hyperref[lean:Litlib.Y1975.belavin1975pseudoparticle.Eq8]{\detokenize{Eq8}}§ BoundaryManifold §\hyperref[lean:CGD.Foundations.SU2Group]{\detokenize{SU2Group}}§ Continuous windingNumber cartanMaurerIntegral]
  (pu : §\hyperref[lean:CGD.Axioms.PhysicalUniverse]{\detokenize{PhysicalUniverse}}§)
  (h_homeo : §\hyperref[lean:Litlib.Y1975.belavin1975pseudoparticle.IsHomeomorphism]{\detokenize{IsHomeomorphism}}§ (boundaryMap pu.toUniverse.asd_sector.val)) :
  §\hyperref[lean:CGD.Particles.inertialMass]{\detokenize{inertialMass}}§ boundaryMap cartanMaurerIntegral pu > 0


-- MODULE: CGD.Particles.Matter
--------------------------------------------------------------------
/--
Proves that the CGD ontology is not an empty vacuum. A topological connection
with a non-zero anti-self-dual curvature natively generates a strictly non-zero
physical Stress-Energy tensor (T_μν ≠ 0), representing the presence of matter.

This relies strictly on the Plebanski expansion, proving that if the physical
energy-momentum tensor vanishes, the anti-self-dual expansion coefficients vanish,
extinguishing the topological matter field entirely.
-/
@[litlib_track "Dynamic Matter Existence"]
theorem §\phantomsection\label{lean:CGD.Particles.dynamicMatterExistence}§CGD.Particles.dynamicMatterExistence
  (pu : §\hyperref[lean:CGD.Axioms.PhysicalUniverse]{\detokenize{PhysicalUniverse}}§)
  (x : §\hyperref[lean:CGD.Foundations.SpacetimePoint]{\detokenize{SpacetimePoint}}§)
  (Sigma Sigma_bar : Fin 3 → Fin 4 → Fin 4 → ℂ)

  -- The true Plebanski expansion coefficients
  (F_ij F_bar_ij T_ij : Fin 3 → Fin 3 → ℂ)

  (Lambda G T_scalar : ℂ)
  (plebanski_matter_eqs : Prop)
  (h_G : G ≠ 0)

  -- Vector representation of the SL(2,C) curvature
  (eval_SL2C : §\hyperref[lean:CGD.Foundations.SL2C]{\detokenize{SL2C}}§ → Fin 3 → ℂ)
  (h_eval_inj : ∀ A, (∀ i, eval_SL2C A i = 0) → A = 0)

  -- Plebanski decomposition of the physical anti-self-dual curvature
  (h_F_asd_decomp : ∀ μ ν i, eval_SL2C (§\hyperref[lean:CGD.Foundations.curvatureSl2c]{\detokenize{curvatureSl2c}}§ pu.toUniverse.asd_sector μ ν x) i = ∑ j, F_bar_ij i j * Sigma_bar j μ ν)

  -- The physics: Eq16 relates the emergent CGD Stress-Energy to the internal T_ij
  (eq16 : §\hyperref[lean:Litlib.Y2011.krasnov2011plebanski.Eq16]{\detokenize{Eq16}}§
    Sigma
    Sigma_bar
    (fun μ ν => §\hyperref[lean:CGD.Gravity.matrixInv4x4]{\detokenize{matrixInv4x4}}§ (fun m n => §\hyperref[lean:CGD.Gravity.urbantkeMetric]{\detokenize{urbantkeMetric}}§ (fun a b => §\hyperref[lean:CGD.Foundations.curvatureSl2c]{\detokenize{curvatureSl2c}}§ pu.toUniverse.sd_sector a b x) m n) μ ν)
    (fun μ ν => §\hyperref[lean:CGD.Gravity.emergentStressEnergy]{\detokenize{emergentStressEnergy}}§ (fun a b p => §\hyperref[lean:CGD.Foundations.curvatureSl2c]{\detokenize{curvatureSl2c}}§ pu.toUniverse.sd_sector a b p) μ ν x)
    T_ij)

  -- The physics: Eq17 directly binds the expansion coefficients to the Stress-Energy
  (eq17 : §\hyperref[lean:Litlib.Y2011.krasnov2011plebanski.Eq17]{\detokenize{Eq17}}§
    Lambda G F_ij F_bar_ij T_scalar T_ij plebanski_matter_eqs)
  (h_matter : plebanski_matter_eqs)

  -- The non-vacuum condition: The physical matter curvature is non-zero
  (h_non_vacuum : ∃ μ ν, §\hyperref[lean:CGD.Foundations.curvatureSl2c]{\detokenize{curvatureSl2c}}§ pu.toUniverse.asd_sector μ ν x ≠ 0) :

  ∃ ρ μ, §\hyperref[lean:CGD.Gravity.emergentStressEnergy]{\detokenize{emergentStressEnergy}}§ (fun a b p => §\hyperref[lean:CGD.Foundations.curvatureSl2c]{\detokenize{curvatureSl2c}}§ pu.toUniverse.sd_sector a b p) ρ μ x ≠ 0


-- MODULE: CGD.Particles.SpinStatistics
--------------------------------------------------------------------
/--
While the local geometric spinor is a commuting Kähler-Dirac field, 
the macroscopic topological defect natively possesses fermionic spin-statistics. 
Via the Finkelstein-Rubinstein theorem (explicitly cited by Witten 1983 for the SU(2) exception),
the 2π adiabatic spatial rotation of a degree-1 SU(2) topological soliton 
traces the non-trivial element of π_4(SU(2)) = Z_2. The resulting quantum 
phase amplitude evaluates strictly to -1.

Because this evaluates directly against the Physical Universe's Anti-Self-Dual (Matter) 
sector via a valid topological projection, the native smoothness (ContDiff) of the 
spacetime gauge field natively fulfills the differentiability requirements without 
external axioms.
-/
@[litlib_track "Macroscopic Topological Fermionization"]
theorem §\phantomsection\label{lean:CGD.Particles.kinematicMacroscopicFermion}§CGD.Particles.kinematicMacroscopicFermion
  [MeasureTheory.MeasureSpace (Fin 3 → ℝ)]
  (pu : §\hyperref[lean:CGD.Axioms.PhysicalUniverse]{\detokenize{PhysicalUniverse}}§)
  (proj : §\hyperref[lean:CGD.Particles.FermionBoundaryProjection]{\detokenize{FermionBoundaryProjection}}§)
  (state : §\hyperref[lean:CGD.Particles.IsFermionicState]{\detokenize{IsFermionicState}}§ pu proj)
  (path_parity : (ℝ → (Fin 3 → ℝ) → Matrix (Fin 2) (Fin 2) ℂ) → ℕ)
  (U_rot : ℝ → (Fin 3 → ℝ) → Matrix (Fin 2) (Fin 2) ℂ)
  (rotation : §\hyperref[lean:CGD.Particles.IsFermionicRotation]{\detokenize{IsFermionicRotation}}§ path_parity (proj.map pu.toUniverse.asd_sector) U_rot)
  (h_fr : §\hyperref[lean:FinkelsteinRubinsteinQuantization]{\detokenize{FinkelsteinRubinsteinQuantization}}§ path_parity 
    (proj.map pu.toUniverse.asd_sector) 
    (proj.h_SU2 pu.toUniverse.asd_sector) 
    (proj.h_smooth pu.toUniverse.asd_sector) 
    state.h_vacuum state.h_integrable state.h_degree_one 
    U_rot rotation.h_rot_SU2 rotation.h_is_2pi_rot rotation.h_fr_parity) :
  §\hyperref[lean:quantum_weight]{\detokenize{quantum_weight}}§ path_parity U_rot = -1


-- MODULE: CGD.Particles.TopologicalStability
--------------------------------------------------------------------
§\phantomsection\label{lean:CGD.Particles.instNonemptyS3}§instance : Nonempty §\hyperref[lean:CGD.Foundations.S3]{\detokenize{S3}}§

§\phantomsection\label{lean:CGD.Particles.instNonemptySU2Group}§instance : Nonempty §\hyperref[lean:CGD.Foundations.SU2Group]{\detokenize{SU2Group}}§

§\phantomsection\label{lean:CGD.Particles.instTopologicalSpaceSU2Group}§noncomputable instance : TopologicalSpace §\hyperref[lean:CGD.Foundations.SU2Group]{\detokenize{SU2Group}}§ := instTopologicalSpaceSubtype

/--
A valid topological homotopy must maintain the integrity of the physical boundary state.
If the continuous mapping to the SU(2) topology breaks, the state has undergone a violent
topological decay.
-/
def §\phantomsection\label{lean:CGD.Particles.isHomotopicConnection}§CGD.Particles.isHomotopicConnection (A0 A1 : Fin 4 → §\hyperref[lean:CGD.Foundations.SpacetimePoint]{\detokenize{SpacetimePoint}}§ → §\hyperref[lean:CGD.Foundations.SL2C]{\detokenize{SL2C}}§) : Prop :=
  ∃ (H : ℝ → Fin 4 → §\hyperref[lean:CGD.Foundations.SpacetimePoint]{\detokenize{SpacetimePoint}}§ → §\hyperref[lean:CGD.Foundations.SL2C]{\detokenize{SL2C}}§),
    (∀ mu x, H 0 mu x = A0 mu x) ∧
    (∀ mu x, H 1 mu x = A1 mu x) ∧
    (∀ mu i j, ContDiff ℝ ⊤ (fun (tx : ℝ × §\hyperref[lean:CGD.Foundations.SpacetimePoint]{\detokenize{SpacetimePoint}}§) => (H tx.1 mu tx.2).val i j)) ∧
    (Continuous (fun t => geometricBoundaryProjection (H t)))

/--
Establishes the absolute topological stability of the instanton configuration. Because the boundary mapping of the instanton constitutes a topological homeomorphism to the SU(2) gauge group, its Cartan-Maurer mapping degree must be strictly non-zero. Consequently, the state is topologically protected and cannot continuously decay into the trivial vacuum.
-/
@[litlib_track "Topological Stability of the Instanton"]
theorem §\phantomsection\label{lean:CGD.Particles.kinematicTopologicalStability}§CGD.Particles.kinematicTopologicalStability
  (windingNumber : (§\hyperref[lean:CGD.Foundations.S3]{\detokenize{S3}}§ → §\hyperref[lean:CGD.Foundations.SU2Group]{\detokenize{SU2Group}}§) → ℤ)
  (cartanMaurerIntegral : (§\hyperref[lean:CGD.Foundations.S3]{\detokenize{S3}}§ → §\hyperref[lean:CGD.Foundations.SU2Group]{\detokenize{SU2Group}}§) → ℝ)
  [tc : §\hyperref[lean:Litlib.Y2003.nakahara2003geometry.CartanMaurerTopology]{\detokenize{CartanMaurerTopology}}§ (§\hyperref[lean:CGD.Foundations.S3]{\detokenize{S3}}§ → §\hyperref[lean:CGD.Foundations.SU2Group]{\detokenize{SU2Group}}§) Continuous windingNumber cartanMaurerIntegral]
  [belavin : §\hyperref[lean:Litlib.Y1975.belavin1975pseudoparticle.Eq8]{\detokenize{Eq8}}§ §\hyperref[lean:CGD.Foundations.S3]{\detokenize{S3}}§ §\hyperref[lean:CGD.Foundations.SU2Group]{\detokenize{SU2Group}}§ Continuous windingNumber cartanMaurerIntegral]
  (integral_zero : cartanMaurerIntegral 1 = 0) :
  ¬ §\hyperref[lean:CGD.Particles.isHomotopicConnection]{\detokenize{isHomotopicConnection}}§ §\hyperref[lean:CGD.Particles.topologicalSoliton]{\detokenize{topologicalSoliton}}§ 0

noncomputable def §\phantomsection\label{lean:CGD.Particles.solitonAsymptoticMap}§CGD.Particles.solitonAsymptoticMap (x : §\hyperref[lean:CGD.Foundations.S3]{\detokenize{S3}}§) : §\hyperref[lean:CGD.Foundations.SU2Group]{\detokenize{SU2Group}}§ := by
  cases x
  rename_i v hv

  let a : ℂ := Complex.mk (v 0) (-(v 3))
  let b : ℂ := Complex.mk (-(v 2)) (-(v 1))

  have h_sum : a * star a + b * star b = 1 := by
    apply Complex.ext
    · change (v 0) * (v 0) - (-(v 3)) * (-(-(v 3))) + ((-(v 2)) * (-(v 2)) - (-(v 1)) * (-(-(v 1)))) = 1
      calc (v 0) * (v 0) - (-(v 3)) * (-(-(v 3))) + ((-(v 2)) * (-(v 2)) - (-(v 1)) * (-(-(v 1))))
        _ = (v 0)^2 + (v 1)^2 + (v 2)^2 + (v 3)^2 := by ring
        _ = 1 := hv
    · change (v 0) * (-(-(v 3))) + (-(v 3)) * (v 0) + ((-(v 2)) * (-(-(v 1))) + (-(v 1)) * (-(v 2))) = 0
      ring

  let M : Matrix (Fin 2) (Fin 2) ℂ := !![a, b; -star b, star a]
  have h_prop : M * M.conjTranspose = 1 ∧ M.det = 1 := su2_matrix_prop a b h_sum

  exact ⟨M, h_prop⟩

@[litlib_track "Topological Soliton Is Homeomorphism"]
theorem §\phantomsection\label{lean:CGD.Particles.soliton_is_homeomorphism}§CGD.Particles.soliton_is_homeomorphism : §\hyperref[lean:Litlib.Y1975.belavin1975pseudoparticle.IsHomeomorphism]{\detokenize{IsHomeomorphism}}§ §\hyperref[lean:CGD.Particles.solitonAsymptoticMap]{\detokenize{solitonAsymptoticMap}}§


-- MODULE: CGD.Phenomenology.AxialCondensate
--------------------------------------------------------------------
/-- The pure Axial Chiral field, defined as the antisymmetric superposition of the left and right sectors. -/
noncomputable def §\phantomsection\label{lean:CGD.Phenomenology.axialField}§CGD.Phenomenology.axialField (u : §\hyperref[lean:CGD.Axioms.Universe]{\detokenize{Universe}}§) (mu : Fin 4) (x : §\hyperref[lean:CGD.Foundations.SpacetimePoint]{\detokenize{SpacetimePoint}}§) : Matrix (Fin 2) (Fin 2) ℂ :=
  (1 / 2 : ℂ) • ((u.sd_sector mu x).val - (u.asd_sector mu x).val)

/--
Proves that the Axial field is strictly an isovector (isospin 1).
It inherits the exact traceless property of the SL(2,C) algebra constituents,
confining it to the adjoint representation of the SU(2) color group.
-/
@[litlib_track "Axial Is Isovector"]
theorem §\phantomsection\label{lean:CGD.Phenomenology.axialIsIsovector}§CGD.Phenomenology.axialIsIsovector (pu : §\hyperref[lean:CGD.Axioms.PhysicalUniverse]{\detokenize{PhysicalUniverse}}§) (mu : Fin 4) (x : §\hyperref[lean:CGD.Foundations.SpacetimePoint]{\detokenize{SpacetimePoint}}§) :
  Matrix.trace (§\hyperref[lean:CGD.Phenomenology.axialField]{\detokenize{axialField}}§ pu.toUniverse mu x) = 0

/--
Proves that the Axial field behaves strictly as a pseudo-vector (parity-odd).
When the geometry is parity-inverted, the Axial field exactly flips its algebraic sign.
-/
@[litlib_track "Axial Is Parity Odd"]
theorem §\phantomsection\label{lean:CGD.Phenomenology.axialIsParityOdd}§CGD.Phenomenology.axialIsParityOdd (pu : §\hyperref[lean:CGD.Axioms.PhysicalUniverse]{\detokenize{PhysicalUniverse}}§) (mu : Fin 4) (x : §\hyperref[lean:CGD.Foundations.SpacetimePoint]{\detokenize{SpacetimePoint}}§) :
  §\hyperref[lean:CGD.Phenomenology.axialField]{\detokenize{axialField}}§ (§\hyperref[lean:CGD.Phenomenology.paritySwap]{\detokenize{paritySwap}}§ pu.toUniverse) mu x = - §\hyperref[lean:CGD.Phenomenology.axialField]{\detokenize{axialField}}§ pu.toUniverse mu x

/--
Generic gauge-covariant derivative action on a fundamental state.
Represented abstractly for a matrix-valued state `Psi` and its bare partial derivative `partial_Psi`.
-/
noncomputable def §\phantomsection\label{lean:CGD.Phenomenology.covariantStateDeriv}§CGD.Phenomenology.covariantStateDeriv (partial_Psi : Matrix (Fin 2) (Fin 2) ℂ) (A : Matrix (Fin 2) (Fin 2) ℂ) (Psi : Matrix (Fin 2) (Fin 2) ℂ) : Matrix (Fin 2) (Fin 2) ℂ :=
  partial_Psi + A * Psi

/--
Proves that the Left-Handed (Self-Dual) gauge connection, which natively couples to
left-handed fermions (such as neutrinos) in the chiral Dirac framework, perfectly
decomposes into a symmetric Vector background and the P-violating Axial condensate.

This establishes that left-handed states must continuously interact with the
axial condensate as a refractive background. Because this interaction occurs purely
at the level of the connection (zero Plebanski curvature), it geometrically derives
kinematic flavor oscillation (the MSW effect) without generating gravitational rest mass.
-/
@[litlib_track "Kinematic Left-Handed Coupling"]
theorem §\phantomsection\label{lean:CGD.Phenomenology.kinematicLeftHandedCoupling}§CGD.Phenomenology.kinematicLeftHandedCoupling (pu : §\hyperref[lean:CGD.Axioms.PhysicalUniverse]{\detokenize{PhysicalUniverse}}§) (mu : Fin 4) (x : §\hyperref[lean:CGD.Foundations.SpacetimePoint]{\detokenize{SpacetimePoint}}§) :
  (pu.toUniverse.sd_sector mu x).val = §\hyperref[lean:CGD.Phenomenology.vectorField]{\detokenize{vectorField}}§ pu.toUniverse mu x + §\hyperref[lean:CGD.Phenomenology.axialField]{\detokenize{axialField}}§ pu.toUniverse mu x

/--
Proves that the covariant derivative of a state under the Left-Handed (Self-Dual)
connection mathematically evaluates to the bare partial derivative, plus the
interaction with the symmetric Vector background, plus the interaction with the
P-violating Axial condensate.

This formalizes the algebraic substitution allowing the paper to definitively claim
that left-handed fermions actively interact with the axial condensate in the
covariant Dirac equation.
-/
@[litlib_track "Kinematic Left-Handed Phase Shift"]
theorem §\phantomsection\label{lean:CGD.Phenomenology.kinematicLeftHandedPhaseShift}§CGD.Phenomenology.kinematicLeftHandedPhaseShift (pu : §\hyperref[lean:CGD.Axioms.PhysicalUniverse]{\detokenize{PhysicalUniverse}}§) (mu : Fin 4) (x : §\hyperref[lean:CGD.Foundations.SpacetimePoint]{\detokenize{SpacetimePoint}}§)
  (Psi partial_Psi : Matrix (Fin 2) (Fin 2) ℂ) :
  §\hyperref[lean:CGD.Phenomenology.covariantStateDeriv]{\detokenize{covariantStateDeriv}}§ partial_Psi (pu.toUniverse.sd_sector mu x).val Psi =
  partial_Psi + §\hyperref[lean:CGD.Phenomenology.vectorField]{\detokenize{vectorField}}§ pu.toUniverse mu x * Psi + §\hyperref[lean:CGD.Phenomenology.axialField]{\detokenize{axialField}}§ pu.toUniverse mu x * Psi

/--
Proves that the Right-Handed (Anti-Self-Dual) gauge connection decomposes into the
symmetric Vector background minus the P-violating Axial condensate. The sign flip
relative to the Left-Handed sector strictly defines the geometric parity asymmetry
of the chiral spacetime vacuum.
-/
@[litlib_track "Kinematic Right-Handed Coupling"]
theorem §\phantomsection\label{lean:CGD.Phenomenology.kinematicRightHandedCoupling}§CGD.Phenomenology.kinematicRightHandedCoupling (pu : §\hyperref[lean:CGD.Axioms.PhysicalUniverse]{\detokenize{PhysicalUniverse}}§) (mu : Fin 4) (x : §\hyperref[lean:CGD.Foundations.SpacetimePoint]{\detokenize{SpacetimePoint}}§) :
  (pu.toUniverse.asd_sector mu x).val = §\hyperref[lean:CGD.Phenomenology.vectorField]{\detokenize{vectorField}}§ pu.toUniverse mu x - §\hyperref[lean:CGD.Phenomenology.axialField]{\detokenize{axialField}}§ pu.toUniverse mu x

/--
Because the macroscopic volume emergent metric strictly forbids global chiral symmetry
(via `macroscopicVolumeImpliesChirality`), this mathematically guarantees that the spacetime
background is natively populated by a strictly non-zero Axial-Vector condensate. Empty
space spontaneously generates the axial field to preserve its volume.
-/
@[litlib_track "Macroscopic Volume Implies Axial Condensate"]
theorem §\phantomsection\label{lean:CGD.Phenomenology.macroscopicVolumeImpliesAxialCondensate}§CGD.Phenomenology.macroscopicVolumeImpliesAxialCondensate
  (pu : §\hyperref[lean:CGD.Axioms.PhysicalUniverse]{\detokenize{PhysicalUniverse}}§) (x : §\hyperref[lean:CGD.Foundations.SpacetimePoint]{\detokenize{SpacetimePoint}}§) (hx : x ∈ pu.bulk)
  (h_vacuum : ∀ μ ν, §\hyperref[lean:CGD.Foundations.curvatureSl2c]{\detokenize{curvatureSl2c}}§ pu.toUniverse.asd_sector.val μ ν x = 0) :
  ∃ y mu, §\hyperref[lean:CGD.Phenomenology.axialField]{\detokenize{axialField}}§ pu.toUniverse mu y ≠ 0

/--
A discrete geometric transformation that completely inverts the orientation of the
spacetime manifold by swapping the Chiral left (self-dual) and right (anti-self-dual) sectors.
-/
noncomputable def §\phantomsection\label{lean:CGD.Phenomenology.paritySwap}§CGD.Phenomenology.paritySwap (u : §\hyperref[lean:CGD.Axioms.Universe]{\detokenize{Universe}}§) : §\hyperref[lean:CGD.Axioms.Universe]{\detokenize{Universe}}§ :=
  universeEquiv.symm (u.asd_sector, u.sd_sector)

/-- The pure Vector Chiral field, defined as the symmetric superposition of the left and right sectors. -/
noncomputable def §\phantomsection\label{lean:CGD.Phenomenology.vectorField}§CGD.Phenomenology.vectorField (u : §\hyperref[lean:CGD.Axioms.Universe]{\detokenize{Universe}}§) (mu : Fin 4) (x : §\hyperref[lean:CGD.Foundations.SpacetimePoint]{\detokenize{SpacetimePoint}}§) : Matrix (Fin 2) (Fin 2) ℂ :=
  (1 / 2 : ℂ) • ((u.sd_sector mu x).val + (u.asd_sector mu x).val)


-- MODULE: CGD.Phenomenology.Chirality
--------------------------------------------------------------------
/--
Proves that a perfectly symmetric, non-chiral universe mathematically destroys itself.

If the Left (Gravity) and Right (Matter) gauge fields are identical in a pure vacuum
(where the Anti-Self-Dual curvature is zero), the Self-Dual (Gravity) curvature is also
extinguished. A zero gravity curvature generates a zero Urbantke metric, which algebraically
forces det(g) = 0. This strictly violates the Macroscopic Volume axiom.
Therefore, empty space must be chiral.
-/
@[litlib_track "Macroscopic volume must have chirality"]
theorem §\phantomsection\label{lean:CGD.Phenomenology.macroscopicVolumeImpliesChirality}§CGD.Phenomenology.macroscopicVolumeImpliesChirality
  (pu : §\hyperref[lean:CGD.Axioms.PhysicalUniverse]{\detokenize{PhysicalUniverse}}§)
  (x : §\hyperref[lean:CGD.Foundations.SpacetimePoint]{\detokenize{SpacetimePoint}}§)
  (hx : x ∈ pu.bulk)
  (h_vacuum : ∀ μ ν, §\hyperref[lean:CGD.Foundations.curvatureSl2c]{\detokenize{curvatureSl2c}}§ pu.toUniverse.asd_sector.val μ ν x = 0) :
  pu.toUniverse.sd_sector.val ≠ pu.toUniverse.asd_sector.val


-- MODULE: CGD.Phenomenology.TmdSignFlips
--------------------------------------------------------------------
/--
The Boer-Mulders effect observable (Quark Transverse Spin vs Quark Transverse Momentum).
Projects the orthogonal SU(2) topology via sigmaX.
Geometrically identical to the Sivers projection.
-/
noncomputable def §\phantomsection\label{lean:CGD.Phenomenology.boerMuldersTransverseKick}§CGD.Phenomenology.boerMuldersTransverseKick (U U_inv : Matrix (Fin 2) (Fin 2) ℂ) : ℂ :=
  Matrix.trace (U * explicitSigmaX * U_inv * explicitSigmaY)

/--
Proves that when evaluating continuous non-Abelian geometry, the Sivers and Boer-Mulders effects yield topologically identical observables.

While perturbative QCD treats them as independent non-perturbative functions, phenomenological
models (like Large-Nc and lattice QCD) observe strong proportionalities. This witness demonstrates how continuous background geometries intrinsically explain this: the macroscopic SU(2) holonomy is strictly blind to the composite vs. bare nature of the initial state, resolving both to the exact same geometric projection.
-/
@[litlib_track "Kinematic Sivers and Boer-Mulders Equivalence Witness"]
theorem §\phantomsection\label{lean:CGD.Phenomenology.kinematicSiversBoerMuldersEquivalence}§CGD.Phenomenology.kinematicSiversBoerMuldersEquivalence (pu : §\hyperref[lean:CGD.Axioms.PhysicalUniverse]{\detokenize{CGD.Axioms.PhysicalUniverse}}§) :
  ∀ (matrixExp : Matrix (Fin 2) (Fin 2) ℂ → Matrix (Fin 2) (Fin 2) ℂ)
    [§\hyperref[lean:Litlib.Y2000.hall2000elementary.DerivativeExponential]{\detokenize{Litlib.Y2000.hall2000elementary.DerivativeExponential}}§ (Fin 2) matrixExp]
    (alpha L : ℝ),
    §\hyperref[lean:CGD.Phenomenology.boerMuldersTransverseKick]{\detokenize{boerMuldersTransverseKick}}§
      (§\hyperref[lean:CGD.Quantum.macroscopicObservable]{\detokenize{CGD.Quantum.macroscopicObservable}}§ (§\hyperref[lean:CGD.Quantum.holonomy]{\detokenize{CGD.Quantum.holonomy}}§ matrixExp) (fun mu p => §\hyperref[lean:CGD.Quantum.rotateYAxis]{\detokenize{CGD.Quantum.rotateYAxis}}§ (fun m p => pu.toUniverse.sd_sector m p) alpha mu p) 1 L)
      (§\hyperref[lean:CGD.Quantum.macroscopicObservable]{\detokenize{CGD.Quantum.macroscopicObservable}}§ (§\hyperref[lean:CGD.Quantum.holonomy]{\detokenize{CGD.Quantum.holonomy}}§ matrixExp) (fun mu p => §\hyperref[lean:CGD.Quantum.rotateYAxis]{\detokenize{CGD.Quantum.rotateYAxis}}§ (fun m p => pu.toUniverse.sd_sector m p) alpha mu p) 1 (-L)) =
    §\hyperref[lean:CGD.Phenomenology.siversTransverseKick]{\detokenize{siversTransverseKick}}§
      (§\hyperref[lean:CGD.Quantum.macroscopicObservable]{\detokenize{CGD.Quantum.macroscopicObservable}}§ (§\hyperref[lean:CGD.Quantum.holonomy]{\detokenize{CGD.Quantum.holonomy}}§ matrixExp) (fun mu p => §\hyperref[lean:CGD.Quantum.rotateYAxis]{\detokenize{CGD.Quantum.rotateYAxis}}§ (fun m p => pu.toUniverse.sd_sector m p) alpha mu p) 1 L)
      (§\hyperref[lean:CGD.Quantum.macroscopicObservable]{\detokenize{CGD.Quantum.macroscopicObservable}}§ (§\hyperref[lean:CGD.Quantum.holonomy]{\detokenize{CGD.Quantum.holonomy}}§ matrixExp) (fun mu p => §\hyperref[lean:CGD.Quantum.rotateYAxis]{\detokenize{CGD.Quantum.rotateYAxis}}§ (fun m p => pu.toUniverse.sd_sector m p) alpha mu p) 1 (-L))

/--
Evaluating the Sivers observable upon the `fluxTubeFrame` ansatz natively forces an exact geometric sign-flip upon path inversion (L -> -L).
While standard perturbative QCD derivations require external factorized gauge links to explain this effect, this theorem acts as an explicit constructive witness demonstrating that the continuous macroscopic gauge geometry intrinsically contains the parity-inverted mechanisms required for the Sivers effect.
-/
@[litlib_track "Kinematic Sivers Sign Flip Witness"]
theorem §\phantomsection\label{lean:CGD.Phenomenology.kinematicSiversSignFlip}§CGD.Phenomenology.kinematicSiversSignFlip (pu : §\hyperref[lean:CGD.Axioms.PhysicalUniverse]{\detokenize{CGD.Axioms.PhysicalUniverse}}§) :
  ∀ (matrixExp : Matrix (Fin 2) (Fin 2) ℂ → Matrix (Fin 2) (Fin 2) ℂ)
    [§\hyperref[lean:Litlib.Y2000.hall2000elementary.DerivativeExponential]{\detokenize{Litlib.Y2000.hall2000elementary.DerivativeExponential}}§ (Fin 2) matrixExp]
    (alpha L : ℝ),
    (∀ t, pu.toUniverse.sd_sector 1 (§\hyperref[lean:CGD.Quantum.straightLinePath]{\detokenize{CGD.Quantum.straightLinePath}}§ t) = §\hyperref[lean:CGD.Quantum.fluxTubeFrame]{\detokenize{CGD.Quantum.fluxTubeFrame}}§ 1 (§\hyperref[lean:CGD.Quantum.straightLinePath]{\detokenize{CGD.Quantum.straightLinePath}}§ t)) →
    §\hyperref[lean:CGD.Phenomenology.siversTransverseKick]{\detokenize{siversTransverseKick}}§
      (§\hyperref[lean:CGD.Quantum.macroscopicObservable]{\detokenize{CGD.Quantum.macroscopicObservable}}§ (§\hyperref[lean:CGD.Quantum.holonomy]{\detokenize{CGD.Quantum.holonomy}}§ matrixExp) (fun mu p => §\hyperref[lean:CGD.Quantum.rotateYAxis]{\detokenize{CGD.Quantum.rotateYAxis}}§ (fun m p => pu.toUniverse.sd_sector m p) alpha mu p) 1 L)
      (§\hyperref[lean:CGD.Quantum.macroscopicObservable]{\detokenize{CGD.Quantum.macroscopicObservable}}§ (§\hyperref[lean:CGD.Quantum.holonomy]{\detokenize{CGD.Quantum.holonomy}}§ matrixExp) (fun mu p => §\hyperref[lean:CGD.Quantum.rotateYAxis]{\detokenize{CGD.Quantum.rotateYAxis}}§ (fun m p => pu.toUniverse.sd_sector m p) alpha mu p) 1 (-L)) =
    - §\hyperref[lean:CGD.Phenomenology.siversTransverseKick]{\detokenize{siversTransverseKick}}§
      (§\hyperref[lean:CGD.Quantum.macroscopicObservable]{\detokenize{CGD.Quantum.macroscopicObservable}}§ (§\hyperref[lean:CGD.Quantum.holonomy]{\detokenize{CGD.Quantum.holonomy}}§ matrixExp) (fun mu p => §\hyperref[lean:CGD.Quantum.rotateYAxis]{\detokenize{CGD.Quantum.rotateYAxis}}§ (fun m p => pu.toUniverse.sd_sector m p) alpha mu p) 1 (-L))
      (§\hyperref[lean:CGD.Quantum.macroscopicObservable]{\detokenize{CGD.Quantum.macroscopicObservable}}§ (§\hyperref[lean:CGD.Quantum.holonomy]{\detokenize{CGD.Quantum.holonomy}}§ matrixExp) (fun mu p => §\hyperref[lean:CGD.Quantum.rotateYAxis]{\detokenize{CGD.Quantum.rotateYAxis}}§ (fun m p => pu.toUniverse.sd_sector m p) alpha mu p) 1 L)

/--
The Topological TMD Geometric Ratio Witness.

By evaluating the `fluxTubeFrame` ansatz, this theorem acts as a constructive witness proving that the Worm-Gear and Sivers observables can natively become geometrically locked by the chiral phase angle
`alpha` of the macroscopic connection. Specifically:
`cos(alpha) * WormGear = - sin(alpha) * Sivers`

Because the observables are defined natively as geometric integrals without any
collisional momentum variables, this geometric lock demonstrates a framework mechanism where their ratio
resolves as a flat kinematic constant (-tan(alpha)), reproducing global supercomputer fits.
-/
@[litlib_track "Topological TMD Geometric Ratio Witness"]
theorem §\phantomsection\label{lean:CGD.Phenomenology.kinematicTmdRatio}§CGD.Phenomenology.kinematicTmdRatio (pu : §\hyperref[lean:CGD.Axioms.PhysicalUniverse]{\detokenize{CGD.Axioms.PhysicalUniverse}}§) :
  ∀ (matrixExp : Matrix (Fin 2) (Fin 2) ℂ → Matrix (Fin 2) (Fin 2) ℂ)
    [§\hyperref[lean:Litlib.Y2000.hall2000elementary.DerivativeExponential]{\detokenize{Litlib.Y2000.hall2000elementary.DerivativeExponential}}§ (Fin 2) matrixExp]
    (alpha L : ℝ),
    (∀ t, pu.toUniverse.sd_sector 1 (§\hyperref[lean:CGD.Quantum.straightLinePath]{\detokenize{CGD.Quantum.straightLinePath}}§ t) = §\hyperref[lean:CGD.Quantum.fluxTubeFrame]{\detokenize{CGD.Quantum.fluxTubeFrame}}§ 1 (§\hyperref[lean:CGD.Quantum.straightLinePath]{\detokenize{CGD.Quantum.straightLinePath}}§ t)) →
    (§\hyperref[lean:CGD.Phenomenology.obs_M_A]{\detokenize{obs_M_A}}§ alpha) * §\hyperref[lean:CGD.Phenomenology.wormGearTransverseKick]{\detokenize{wormGearTransverseKick}}§
      (§\hyperref[lean:CGD.Quantum.macroscopicObservable]{\detokenize{CGD.Quantum.macroscopicObservable}}§ (§\hyperref[lean:CGD.Quantum.holonomy]{\detokenize{CGD.Quantum.holonomy}}§ matrixExp) (fun mu p => §\hyperref[lean:CGD.Quantum.rotateYAxis]{\detokenize{CGD.Quantum.rotateYAxis}}§ (fun m p => pu.toUniverse.sd_sector m p) alpha mu p) 1 L)
      (§\hyperref[lean:CGD.Quantum.macroscopicObservable]{\detokenize{CGD.Quantum.macroscopicObservable}}§ (§\hyperref[lean:CGD.Quantum.holonomy]{\detokenize{CGD.Quantum.holonomy}}§ matrixExp) (fun mu p => §\hyperref[lean:CGD.Quantum.rotateYAxis]{\detokenize{CGD.Quantum.rotateYAxis}}§ (fun m p => pu.toUniverse.sd_sector m p) alpha mu p) 1 (-L)) =
    - (§\hyperref[lean:CGD.Phenomenology.obs_M_B]{\detokenize{obs_M_B}}§ alpha) * §\hyperref[lean:CGD.Phenomenology.siversTransverseKick]{\detokenize{siversTransverseKick}}§
      (§\hyperref[lean:CGD.Quantum.macroscopicObservable]{\detokenize{CGD.Quantum.macroscopicObservable}}§ (§\hyperref[lean:CGD.Quantum.holonomy]{\detokenize{CGD.Quantum.holonomy}}§ matrixExp) (fun mu p => §\hyperref[lean:CGD.Quantum.rotateYAxis]{\detokenize{CGD.Quantum.rotateYAxis}}§ (fun m p => pu.toUniverse.sd_sector m p) alpha mu p) 1 L)
      (§\hyperref[lean:CGD.Quantum.macroscopicObservable]{\detokenize{CGD.Quantum.macroscopicObservable}}§ (§\hyperref[lean:CGD.Quantum.holonomy]{\detokenize{CGD.Quantum.holonomy}}§ matrixExp) (fun mu p => §\hyperref[lean:CGD.Quantum.rotateYAxis]{\detokenize{CGD.Quantum.rotateYAxis}}§ (fun m p => pu.toUniverse.sd_sector m p) alpha mu p) 1 (-L))

/--
Evaluating the Worm-Gear observable upon the `fluxTubeFrame` ansatz establishes that it obeys the exact same geometric sign-flip mechanics as the Sivers effect.
-/
@[litlib_track "Kinematic Worm-Gear Sign Flip Witness"]
theorem §\phantomsection\label{lean:CGD.Phenomenology.kinematicWormGearSignFlip}§CGD.Phenomenology.kinematicWormGearSignFlip (pu : §\hyperref[lean:CGD.Axioms.PhysicalUniverse]{\detokenize{CGD.Axioms.PhysicalUniverse}}§) :
  ∀ (matrixExp : Matrix (Fin 2) (Fin 2) ℂ → Matrix (Fin 2) (Fin 2) ℂ)
    [§\hyperref[lean:Litlib.Y2000.hall2000elementary.DerivativeExponential]{\detokenize{Litlib.Y2000.hall2000elementary.DerivativeExponential}}§ (Fin 2) matrixExp]
    (alpha L : ℝ),
    (∀ t, pu.toUniverse.sd_sector 1 (§\hyperref[lean:CGD.Quantum.straightLinePath]{\detokenize{CGD.Quantum.straightLinePath}}§ t) = §\hyperref[lean:CGD.Quantum.fluxTubeFrame]{\detokenize{CGD.Quantum.fluxTubeFrame}}§ 1 (§\hyperref[lean:CGD.Quantum.straightLinePath]{\detokenize{CGD.Quantum.straightLinePath}}§ t)) →
    §\hyperref[lean:CGD.Phenomenology.wormGearTransverseKick]{\detokenize{wormGearTransverseKick}}§
      (§\hyperref[lean:CGD.Quantum.macroscopicObservable]{\detokenize{CGD.Quantum.macroscopicObservable}}§ (§\hyperref[lean:CGD.Quantum.holonomy]{\detokenize{CGD.Quantum.holonomy}}§ matrixExp) (fun mu p => §\hyperref[lean:CGD.Quantum.rotateYAxis]{\detokenize{CGD.Quantum.rotateYAxis}}§ (fun m p => pu.toUniverse.sd_sector m p) alpha mu p) 1 L)
      (§\hyperref[lean:CGD.Quantum.macroscopicObservable]{\detokenize{CGD.Quantum.macroscopicObservable}}§ (§\hyperref[lean:CGD.Quantum.holonomy]{\detokenize{CGD.Quantum.holonomy}}§ matrixExp) (fun mu p => §\hyperref[lean:CGD.Quantum.rotateYAxis]{\detokenize{CGD.Quantum.rotateYAxis}}§ (fun m p => pu.toUniverse.sd_sector m p) alpha mu p) 1 (-L)) =
    - §\hyperref[lean:CGD.Phenomenology.wormGearTransverseKick]{\detokenize{wormGearTransverseKick}}§
      (§\hyperref[lean:CGD.Quantum.macroscopicObservable]{\detokenize{CGD.Quantum.macroscopicObservable}}§ (§\hyperref[lean:CGD.Quantum.holonomy]{\detokenize{CGD.Quantum.holonomy}}§ matrixExp) (fun mu p => §\hyperref[lean:CGD.Quantum.rotateYAxis]{\detokenize{CGD.Quantum.rotateYAxis}}§ (fun m p => pu.toUniverse.sd_sector m p) alpha mu p) 1 (-L))
      (§\hyperref[lean:CGD.Quantum.macroscopicObservable]{\detokenize{CGD.Quantum.macroscopicObservable}}§ (§\hyperref[lean:CGD.Quantum.holonomy]{\detokenize{CGD.Quantum.holonomy}}§ matrixExp) (fun mu p => §\hyperref[lean:CGD.Quantum.rotateYAxis]{\detokenize{CGD.Quantum.rotateYAxis}}§ (fun m p => pu.toUniverse.sd_sector m p) alpha mu p) 1 L)

noncomputable def §\phantomsection\label{lean:CGD.Phenomenology.obs_M_A}§CGD.Phenomenology.obs_M_A (alpha : ℝ) : ℂ := (Complex.cos (alpha/2))^2 - (Complex.sin (alpha/2))^2

noncomputable def §\phantomsection\label{lean:CGD.Phenomenology.obs_M_B}§CGD.Phenomenology.obs_M_B (alpha : ℝ) : ℂ := 2 * (Complex.cos (alpha/2)) * (Complex.sin (alpha/2))

/--
The Sivers effect observable (Nucleon Transverse Spin vs Quark Transverse Momentum).
Projects the orthogonal SU(2) topology via sigmaX.
-/
noncomputable def §\phantomsection\label{lean:CGD.Phenomenology.siversTransverseKick}§CGD.Phenomenology.siversTransverseKick (U U_inv : Matrix (Fin 2) (Fin 2) ℂ) : ℂ :=
  Matrix.trace (U * explicitSigmaX * U_inv * explicitSigmaY)

/--
The Worm-Gear effect observable (Quark Longitudinal Spin vs Quark Transverse Momentum).
Projects the orthogonal SU(2) topology via sigmaZ.
-/
noncomputable def §\phantomsection\label{lean:CGD.Phenomenology.wormGearTransverseKick}§CGD.Phenomenology.wormGearTransverseKick (U U_inv : Matrix (Fin 2) (Fin 2) ℂ) : ℂ :=
  Matrix.trace (U * explicitSigmaZ * U_inv * explicitSigmaY)


-- MODULE: CGD.Quantum.ActionQuantization
--------------------------------------------------------------------
/--
If a physical connection yields an asymptotic boundary map that is a topological homeomorphism to the gauge group, its Cartan-Maurer topological charge strictly evaluates to an integer quantization bound (±1).

**PHYSICAL ONTOLOGY & METRIC INDEPENDENCE:**
The `BoundaryManifold` is intentionally left as an abstract generic type to mathematically satisfy both 
regimes of Chiral Gauge Dynamics without metric dependence:
1. **The Euclidean Instanton (Tunneling):** If `BoundaryManifold` is the $S^3$ boundary of 4D spacetime 
   ($r \to \infty$ in all directions, including time), this theorem strictly reproduces the Belavin 1975 
   instanton action quantization for transient vacuum transitions.
2. **The Lorentzian Soliton (Stable Matter):** If `BoundaryManifold` is the $S^3$ spatial compactification 
   of a 3D Cauchy surface at a fixed time $t$, this theorem quantizes the invariant topological charge 
   of a stable particle (a Skyrmion/Soliton) persisting through macroscopic Lorentzian time.

Because the Cartan-Maurer degree theorem depends only on the continuous boundary mapping and is strictly blind 
to the bulk metric signature, the exact same geometric quantization mathematically governs both phenomena.
-/
@[litlib_track "Topological Action Quantization"]
theorem §\phantomsection\label{lean:CGD.Quantum.kinematicActionQuantization}§CGD.Quantum.kinematicActionQuantization
  {BoundaryManifold : Type*} [TopologicalSpace BoundaryManifold] [Nonempty BoundaryManifold]
  (boundaryMap : (Fin 4 → §\hyperref[lean:CGD.Foundations.SpacetimePoint]{\detokenize{SpacetimePoint}}§ → §\hyperref[lean:CGD.Foundations.SL2C]{\detokenize{SL2C}}§) → BoundaryManifold → §\hyperref[lean:CGD.Foundations.SU2Group]{\detokenize{SU2Group}}§)
  (windingNumber : (BoundaryManifold → §\hyperref[lean:CGD.Foundations.SU2Group]{\detokenize{SU2Group}}§) → ℤ)
  (cartanMaurerIntegral : (BoundaryManifold → §\hyperref[lean:CGD.Foundations.SU2Group]{\detokenize{SU2Group}}§) → ℝ)
  [tc : §\hyperref[lean:Litlib.Y2003.nakahara2003geometry.CartanMaurerTopology]{\detokenize{CartanMaurerTopology}}§ (BoundaryManifold → §\hyperref[lean:CGD.Foundations.SU2Group]{\detokenize{SU2Group}}§) (Continuous : (BoundaryManifold → §\hyperref[lean:CGD.Foundations.SU2Group]{\detokenize{SU2Group}}§) → Prop) windingNumber cartanMaurerIntegral]
  [belavin : §\hyperref[lean:Litlib.Y1975.belavin1975pseudoparticle.Eq8]{\detokenize{Eq8}}§ BoundaryManifold §\hyperref[lean:CGD.Foundations.SU2Group]{\detokenize{SU2Group}}§ (Continuous : (BoundaryManifold → §\hyperref[lean:CGD.Foundations.SU2Group]{\detokenize{SU2Group}}§) → Prop) windingNumber cartanMaurerIntegral]
  (pu : §\hyperref[lean:CGD.Axioms.PhysicalUniverse]{\detokenize{PhysicalUniverse}}§)
  (h_homeo : §\hyperref[lean:Litlib.Y1975.belavin1975pseudoparticle.IsHomeomorphism]{\detokenize{IsHomeomorphism}}§ (boundaryMap pu.toUniverse.sd_sector.val)) :
  cartanMaurerIntegral (boundaryMap pu.toUniverse.sd_sector.val) = 1 ∨
  cartanMaurerIntegral (boundaryMap pu.toUniverse.sd_sector.val) = -1


-- MODULE: CGD.Quantum.Definitions
--------------------------------------------------------------------
/--
Identifies a non-local topological correlation. Two spatial boundaries are defined
as entangled if they are connected by a continuous degenerate connection (a flux tube).
This definition relies on the exact `fluxTubeFrame` witness to guarantee non-vacuous evaluation.

**PHYSICAL CONSERVATION & THE BIANCHI GATEKEEPER:**
Because the universe lacks dynamical Euler-Lagrange equations, the formation and persistence 
of these degenerate connections are strictly governed by the exact Differential Bianchi Identity 
(`kinematicBianchiIdentity`). The Bianchi constraint ($\nabla \cdot B = 0$ analogue) guarantees 
topological flux conservation, mathematically preventing arbitrary spontaneous flux formation 
in the vacuum ("flux soup"). It enforces that tubes must strictly anchor to paired topological 
defects (particles) with equal and opposite winding numbers, geometrically reproducing 
quantum pair production and bipartite entanglement.
-/
@[litlib_track "Entangled State Witness"]
def §\phantomsection\label{lean:CGD.Quantum.areEntangled}§CGD.Quantum.areEntangled (A : Fin 4 → §\hyperref[lean:CGD.Foundations.SpacetimePoint]{\detokenize{SpacetimePoint}}§ → §\hyperref[lean:CGD.Foundations.SL2C]{\detokenize{SL2C}}§) (x y : §\hyperref[lean:CGD.Foundations.SpacetimePoint]{\detokenize{SpacetimePoint}}§) (theta : ℝ) : Prop :=
  ∃ (γ : ℝ → §\hyperref[lean:CGD.Foundations.SpacetimePoint]{\detokenize{SpacetimePoint}}§) (θ : ℝ → ℝ),
    γ 0 = x ∧ γ 1 = y ∧
    θ 0 = 0 ∧ θ 1 = theta ∧
    ∀ t : ℝ, 0 ≤ t → t ≤ 1 →
      (∀ mu, A mu (γ t) = §\hyperref[lean:CGD.Quantum.rotateYAxis]{\detokenize{rotateYAxis}}§ §\hyperref[lean:CGD.Quantum.fluxTubeFrame]{\detokenize{fluxTubeFrame}}§ (θ t) mu (γ t)) ∧
      (∀ mu nu, partialDerivSl2c nu (A mu) (γ t) = partialDerivSl2c nu (§\hyperref[lean:CGD.Quantum.rotateYAxis]{\detokenize{rotateYAxis}}§ §\hyperref[lean:CGD.Quantum.fluxTubeFrame]{\detokenize{fluxTubeFrame}}§ (θ t) mu) (γ t))

/--
A Constructive Witness for an unbroken 1D string defect along the z-axis.
Rather than a unique dynamic solution, this provides an exact analytical ansatz
(a uniform, constant non-Abelian magnetic field) to rigorously prove the non-vacuous
existence of degenerate topological states without relying on empty sets.
-/
@[litlib_track "Flux Tube Frame Witness"]
noncomputable def §\phantomsection\label{lean:CGD.Quantum.fluxTubeFrame}§CGD.Quantum.fluxTubeFrame (mu : Fin 4) (_x : §\hyperref[lean:CGD.Foundations.SpacetimePoint]{\detokenize{SpacetimePoint}}§) : §\hyperref[lean:CGD.Foundations.SL2C]{\detokenize{SL2C}}§ :=
  §\hyperref[lean:CGD.Foundations.toSl2c]{\detokenize{toSl2c}}§ (if mu = 0 then 0 else if mu = 1 then (Complex.I:ℂ) • sigma3.val else if mu = 2 then (Complex.I:ℂ) • sigma1.val else (Complex.I:ℂ) • sigma2.val)

/--
Identifies a field configuration matching the exact analytical witness of an unbroken flux tube.
-/
def §\phantomsection\label{lean:CGD.Quantum.isFluxTube}§CGD.Quantum.isFluxTube (A : Fin 4 → §\hyperref[lean:CGD.Foundations.SpacetimePoint]{\detokenize{SpacetimePoint}}§ → §\hyperref[lean:CGD.Foundations.SL2C]{\detokenize{SL2C}}§) (x : §\hyperref[lean:CGD.Foundations.SpacetimePoint]{\detokenize{SpacetimePoint}}§) : Prop :=
  (∀ mu, A mu x = §\hyperref[lean:CGD.Quantum.fluxTubeFrame]{\detokenize{fluxTubeFrame}}§ mu x) ∧
  (∀ mu nu, partialDerivSl2c nu (A mu) x = partialDerivSl2c nu (§\hyperref[lean:CGD.Quantum.fluxTubeFrame]{\detokenize{fluxTubeFrame}}§ mu) x)

/--
A multi-fingered gauge rotation applied to a background connection.
-/
noncomputable def §\phantomsection\label{lean:CGD.Quantum.rotateYAxis}§CGD.Quantum.rotateYAxis (A : Fin 4 → §\hyperref[lean:CGD.Foundations.SpacetimePoint]{\detokenize{SpacetimePoint}}§ → §\hyperref[lean:CGD.Foundations.SL2C]{\detokenize{SL2C}}§) (theta : Real) (mu : Fin 4) (x : §\hyperref[lean:CGD.Foundations.SpacetimePoint]{\detokenize{SpacetimePoint}}§) : §\hyperref[lean:CGD.Foundations.SL2C]{\detokenize{SL2C}}§ :=
  let R := Matrix.of ![![Complex.cos (theta/2), -Complex.sin (theta/2)], ![Complex.sin (theta/2), Complex.cos (theta/2)]]
  let R_inv := Matrix.of ![![Complex.cos (theta/2), Complex.sin (theta/2)], ![-Complex.sin (theta/2), Complex.cos (theta/2)]]
  §\hyperref[lean:CGD.Foundations.toSl2c]{\detokenize{toSl2c}}§ (R * (A mu x).val * R_inv)


-- MODULE: CGD.Quantum.Dirac.Definitions
--------------------------------------------------------------------
/-- 
The Geometric Spinor (Kähler-Dirac Mode). 
Constructed strictly from the local internal curvature components F_{ab}.
This guarantees flawless gauge covariance (adjoint representation) 
and removes the need for arbitrary fermion fields.
-/
@[litlib_track "Geometric Spinor Definition"]
noncomputable def §\phantomsection\label{lean:CGD.Quantum.Dirac.geometricSpinor}§CGD.Quantum.Dirac.geometricSpinor (F : Fin 4 → Fin 4 → ℂ) : Matrix (Fin 4) (Fin 4) ℂ :=
  ∑ a : Fin 4, ∑ b : Fin 4, F a b • (§\hyperref[lean:Litlib.Math.Dirac.gammaVec]{\detokenize{gammaVec}}§ a * §\hyperref[lean:Litlib.Math.Dirac.gammaVec]{\detokenize{gammaVec}}§ b)

/-- 
The Geometric Yang-Mills Current. 
Defined as the covariant divergence of the curvature tensor.
-/
noncomputable def §\phantomsection\label{lean:CGD.Quantum.Dirac.yangMillsCurrent}§CGD.Quantum.Dirac.yangMillsCurrent (D_F : Fin 4 → Fin 4 → Fin 4 → ℂ) (b : Fin 4) : ℂ :=
  ∑ c : Fin 4, ∑ a : Fin 4, minkowskiEta c a * D_F c a b


-- MODULE: CGD.Quantum.Dirac.Emergence
--------------------------------------------------------------------
/--
Proves that applying the emergent Dirac operator to the geometric spinor natively 
evaluates to the Yang-Mills geometric current. This mathematically establishes that 
the Dirac equation is not a separate physical postulate, but an intrinsic identity 
of the Spin(4,C) gauge field.
-/
@[litlib_track "Geometric Dirac Emergence"]
theorem §\phantomsection\label{lean:CGD.Quantum.Dirac.geometricDiracEmergence}§CGD.Quantum.Dirac.geometricDiracEmergence (D_F : Fin 4 → Fin 4 → Fin 4 → ℂ)
  (h_anti : ∀ c a b, D_F c a b = - D_F c b a)
  (h_bianchi : ∀ c a b, D_F c a b + D_F a b c + D_F b c a = 0) :
  (∑ c : Fin 4, §\hyperref[lean:Litlib.Math.Dirac.gammaVec]{\detokenize{gammaVec}}§ c * (∑ a : Fin 4, ∑ b : Fin 4, D_F c a b • (§\hyperref[lean:Litlib.Math.Dirac.gammaVec]{\detokenize{gammaVec}}§ a * §\hyperref[lean:Litlib.Math.Dirac.gammaVec]{\detokenize{gammaVec}}§ b))) =
  2 • ∑ b : Fin 4, §\hyperref[lean:CGD.Quantum.Dirac.yangMillsCurrent]{\detokenize{yangMillsCurrent}}§ D_F b • §\hyperref[lean:Litlib.Math.Dirac.gammaVec]{\detokenize{gammaVec}}§ b


-- MODULE: CGD.Quantum.Dirac.Vacuum
--------------------------------------------------------------------
/--
The Vacuum Dirac Equation.
If the macroscopic geometry satisfies the source-free vacuum equations (J = 0), 
the geometric spinor mode mathematically and strictly obeys the massless Dirac equation.
-/
@[litlib_track "Vacuum Dirac Equation"]
theorem §\phantomsection\label{lean:CGD.Quantum.Dirac.vacuumDiracEquation}§CGD.Quantum.Dirac.vacuumDiracEquation (D_F : Fin 4 → Fin 4 → Fin 4 → ℂ)
  (h_anti : ∀ c a b, D_F c a b = - D_F c b a)
  (h_bianchi : ∀ c a b, D_F c a b + D_F a b c + D_F b c a = 0)
  (h_vacuum : ∀ b, §\hyperref[lean:CGD.Quantum.Dirac.yangMillsCurrent]{\detokenize{yangMillsCurrent}}§ D_F b = 0) :
  (∑ c : Fin 4, §\hyperref[lean:Litlib.Math.Dirac.gammaVec]{\detokenize{gammaVec}}§ c * (∑ a : Fin 4, ∑ b : Fin 4, D_F c a b • (§\hyperref[lean:Litlib.Math.Dirac.gammaVec]{\detokenize{gammaVec}}§ a * §\hyperref[lean:Litlib.Math.Dirac.gammaVec]{\detokenize{gammaVec}}§ b))) = 0


-- MODULE: CGD.Quantum.Entanglement.Wormhole
--------------------------------------------------------------------
/--
Evaluates the macroscopic metric of the `fluxTubeFrame` witness to prove that if two
particles are connected by such a 1D uniform gauge connection, the metric determinant
rigorously vanishes. This confirms that non-local, degenerate wormhole topologies
mathematically exist as non-vacuous solutions within the Spin(4,C) geometry.
-/
@[litlib_track "Witness for Degenerate Entanglement Channels"]
theorem §\phantomsection\label{lean:CGD.Quantum.kinematicEntanglementWormhole}§CGD.Quantum.kinematicEntanglementWormhole (pu : §\hyperref[lean:CGD.Axioms.PhysicalUniverse]{\detokenize{PhysicalUniverse}}§) :
  ∀ (x y : §\hyperref[lean:CGD.Foundations.SpacetimePoint]{\detokenize{SpacetimePoint}}§) (theta : ℝ),
    §\hyperref[lean:CGD.Quantum.areEntangled]{\detokenize{areEntangled}}§ pu.toUniverse.sd_sector x y theta →
    (§\hyperref[lean:CGD.Gravity.urbantkeMetric]{\detokenize{urbantkeMetric}}§ (fun m n => §\hyperref[lean:CGD.Foundations.curvatureSl2c]{\detokenize{curvatureSl2c}}§ pu.toUniverse.sd_sector m n x)).det = 0 ∧
    (§\hyperref[lean:CGD.Gravity.urbantkeMetric]{\detokenize{urbantkeMetric}}§ (fun m n => §\hyperref[lean:CGD.Foundations.curvatureSl2c]{\detokenize{curvatureSl2c}}§ pu.toUniverse.sd_sector m n y)).det = 0


-- MODULE: CGD.Quantum.FluxTube
--------------------------------------------------------------------
/--
Proves that the exact `fluxTubeFrame` witness natively satisfies the static, purely magnetic boundary conditions.
-/
-- NOTE: Global topological conservation of this minimal configuration is enforced natively by CGD.Foundations.kinematicBianchiIdentity.
@[litlib_track "The FluxTubeFrame is a valid physical solution"]
theorem §\phantomsection\label{lean:CGD.Quantum.fluxTubeIsMinimal}§CGD.Quantum.fluxTubeIsMinimal : ∀ x, §\hyperref[lean:CGD.Quantum.satisfies1DMinimalEnergyBound]{\detokenize{satisfies1DMinimalEnergyBound}}§ §\hyperref[lean:CGD.Quantum.fluxTubeFrame]{\detokenize{fluxTubeFrame}}§ x

/--
Demonstrates that the macroscopic metric of the `fluxTubeFrame` witness possesses a strictly zero determinant. This proves that a 1D uniform gauge connection natively represents a degenerate, non-macroscopic topological background (det(g) = 0), serving as the formal witness that non-local ER=EPR wormhole channels mathematically exist within the Spin(4,C) geometry.
-/
@[litlib_track "Witness for Degenerate Entanglement Channels"]
theorem §\phantomsection\label{lean:CGD.Quantum.kinematicFluxTubeStability}§CGD.Quantum.kinematicFluxTubeStability (pu : §\hyperref[lean:CGD.Axioms.PhysicalUniverse]{\detokenize{PhysicalUniverse}}§) :
  ∀ (x : §\hyperref[lean:CGD.Foundations.SpacetimePoint]{\detokenize{SpacetimePoint}}§),
    §\hyperref[lean:CGD.Quantum.isFluxTube]{\detokenize{isFluxTube}}§ pu.toUniverse.sd_sector x →
    (§\hyperref[lean:CGD.Gravity.urbantkeMetric]{\detokenize{urbantkeMetric}}§ (fun m n => §\hyperref[lean:CGD.Foundations.curvatureSl2c]{\detokenize{curvatureSl2c}}§ pu.toUniverse.sd_sector m n x)).det = 0

/--
Mathematically enforces a strictly static (∂_t A = 0), purely magnetic (electric curvature F_{0ν} = 0) configuration.
Phenomenologically, this represents the minimal energy ground state of a flux tube. Geometrically,
it simply isolates a static background topology.
-/
def §\phantomsection\label{lean:CGD.Quantum.satisfies1DMinimalEnergyBound}§CGD.Quantum.satisfies1DMinimalEnergyBound (A : Fin 4 → §\hyperref[lean:CGD.Foundations.SpacetimePoint]{\detokenize{SpacetimePoint}}§ → §\hyperref[lean:CGD.Foundations.SL2C]{\detokenize{SL2C}}§) (x : §\hyperref[lean:CGD.Foundations.SpacetimePoint]{\detokenize{SpacetimePoint}}§) : Prop :=
  (∀ nu, §\hyperref[lean:CGD.Foundations.curvatureSl2c]{\detokenize{curvatureSl2c}}§ A 0 nu x = 0) ∧
  (∀ nu, partialDerivSl2c 0 (A nu) x = 0)


-- MODULE: CGD.Quantum.Holonomy.Evaluation
--------------------------------------------------------------------
/--
Evaluates the path-ordered integral of the exact `fluxTubeFrame` spatial gauge connection witness.
This explicitly demonstrates that evaluating an SU(2) holonomy natively maps geometric
parameters directly to the trigonometric coefficients of standard quantum state vectors,
proving that the geometry can natively reproduce quantum observables.
-/
@[litlib_track "Macroscopic Spin State Witness"]
theorem §\phantomsection\label{lean:CGD.Quantum.fluxTubeHolonomyEvaluation}§CGD.Quantum.fluxTubeHolonomyEvaluation
  (matrixExp : Matrix (Fin 2) (Fin 2) ℂ → Matrix (Fin 2) (Fin 2) ℂ)
  [§\hyperref[lean:Litlib.Y2000.hall2000elementary.DerivativeExponential]{\detokenize{DerivativeExponential}}§ (Fin 2) matrixExp]
  (pu : §\hyperref[lean:CGD.Axioms.PhysicalUniverse]{\detokenize{PhysicalUniverse}}§) (alpha L : ℝ)
  (h_field : ∀ t, pu.toUniverse.sd_sector 1 (§\hyperref[lean:CGD.Quantum.straightLinePath]{\detokenize{straightLinePath}}§ t) = §\hyperref[lean:CGD.Quantum.fluxTubeFrame]{\detokenize{fluxTubeFrame}}§ 1 (§\hyperref[lean:CGD.Quantum.straightLinePath]{\detokenize{straightLinePath}}§ t)) :
  §\hyperref[lean:CGD.Quantum.macroscopicObservable]{\detokenize{macroscopicObservable}}§ (§\hyperref[lean:CGD.Quantum.holonomy]{\detokenize{holonomy}}§ matrixExp) (fun mu p => §\hyperref[lean:CGD.Quantum.rotateYAxis]{\detokenize{rotateYAxis}}§ (fun m p => pu.toUniverse.sd_sector m p) alpha mu p) 1 L =
  (Complex.cos (L:ℂ)) • 1 + (Complex.I * Complex.sin (L:ℂ)) • §\hyperref[lean:CGD.Quantum.obs_M]{\detokenize{obs_M}}§ alpha

@[litlib_track "Geometric Holonomy"]
-- NATIVE EVALUATION: The exact path-ordered exponential for a constant field
noncomputable def §\phantomsection\label{lean:CGD.Quantum.holonomy}§CGD.Quantum.holonomy (matrixExp : Matrix (Fin 2) (Fin 2) ℂ → Matrix (Fin 2) (Fin 2) ℂ)
  (A : ℝ → Matrix (Fin 2) (Fin 2) ℂ) (t0 t1 : ℝ) : Matrix (Fin 2) (Fin 2) ℂ :=
  matrixExp ((t1 - t0 : ℂ) • A t0)

@[litlib_track "Macroscopic Observable"]
noncomputable def §\phantomsection\label{lean:CGD.Quantum.macroscopicObservable}§CGD.Quantum.macroscopicObservable
  (§\hyperref[lean:CGD.Quantum.holonomy]{\detokenize{holonomy}}§ : (ℝ → Matrix (Fin 2) (Fin 2) ℂ) → ℝ → ℝ → Matrix (Fin 2) (Fin 2) ℂ)
  (A : Fin 4 → §\hyperref[lean:CGD.Foundations.SpacetimePoint]{\detokenize{SpacetimePoint}}§ → §\hyperref[lean:CGD.Foundations.SL2C]{\detokenize{SL2C}}§)
  (mu : Fin 4) (L : ℝ) : Matrix (Fin 2) (Fin 2) ℂ :=
  §\hyperref[lean:CGD.Quantum.holonomy]{\detokenize{holonomy}}§ (fun t => (A mu (§\hyperref[lean:CGD.Quantum.straightLinePath]{\detokenize{straightLinePath}}§ t)).val) 0 L

noncomputable def §\phantomsection\label{lean:CGD.Quantum.obs_M}§CGD.Quantum.obs_M (alpha : ℝ) : Matrix (Fin 2) (Fin 2) ℂ :=
  let R := Matrix.of ![![Complex.cos (alpha/2), -Complex.sin (alpha/2)], ![Complex.sin (alpha/2), Complex.cos (alpha/2)]]
  let R_inv := Matrix.of ![![Complex.cos (alpha/2), Complex.sin (alpha/2)], ![-Complex.sin (alpha/2), Complex.cos (alpha/2)]]
  R * sigma3.val * R_inv

noncomputable def §\phantomsection\label{lean:CGD.Quantum.straightLinePath}§CGD.Quantum.straightLinePath (t : ℝ) : §\hyperref[lean:CGD.Foundations.SpacetimePoint]{\detokenize{SpacetimePoint}}§ := fun _ => t


-- MODULE: CGD.Quantum.Holonomy.Geometric
--------------------------------------------------------------------
/--
Demonstrates that the self-dual and anti-self-dual spin connection components
natively mathematically decouple under the geometric trace metric because their
underlying 4x4 matrix embeddings occupy perfectly disjoint sub-blocks.
The geometry requires no artificial truncation to separate chiral sectors.
-/
@[litlib_track "Algebraic Chiral Orthogonalization"]
theorem §\phantomsection\label{lean:CGD.Quantum.algebraicChiralOrthogonalization}§CGD.Quantum.algebraicChiralOrthogonalization (pu : §\hyperref[lean:CGD.Axioms.PhysicalUniverse]{\detokenize{PhysicalUniverse}}§) (x : §\hyperref[lean:CGD.Foundations.SpacetimePoint]{\detokenize{SpacetimePoint}}§) (μ ν ρ σ : Fin 4) :
  Matrix.trace (§\hyperref[lean:CGD.Foundations.embedSelfDual]{\detokenize{embedSelfDual}}§ (§\hyperref[lean:CGD.Foundations.curvatureSl2c]{\detokenize{CGD.Foundations.curvatureSl2c}}§ pu.toUniverse.sd_sector μ ν x) *
                §\hyperref[lean:CGD.Foundations.embedAntiSelfDual]{\detokenize{embedAntiSelfDual}}§ (§\hyperref[lean:CGD.Foundations.curvatureSl2c]{\detokenize{CGD.Foundations.curvatureSl2c}}§ pu.toUniverse.asd_sector ρ σ x)) = 0

/--
The exact quantum correlation natively emerges from the Cartan-Killing metric
of the SU(2) group. For unitary SU(2) elements, 1/2 Tr(A B†) perfectly recovers
the cosine of the angle between them.
-/
@[litlib_track "Geometric Bell Correlation"]
noncomputable def §\phantomsection\label{lean:CGD.Quantum.geometricBellCorrelation}§CGD.Quantum.geometricBellCorrelation (A B : §\hyperref[lean:CGD.Foundations.SU2Group]{\detokenize{SU2Group}}§) : ℂ :=
  (1 / 2 : ℂ) * Matrix.trace (A.val * B.val.conjTranspose)

/--
Demonstrates that the geometric path-ordered integration of a purely spatial
SU(2) flux tube connection evaluates strictly to a unitary holonomy within
the SU(2) group manifold.
-/
@[litlib_track "Geometric Holonomy Integration"]
theorem §\phantomsection\label{lean:CGD.Quantum.geometricHolonomyIntegration}§CGD.Quantum.geometricHolonomyIntegration (θ : ℝ) :
  let U := Matrix.of ![![Complex.cos (θ/2), -Complex.sin (θ/2)], ![Complex.sin (θ/2), Complex.cos (θ/2)]];
  U * U.conjTranspose = 1 ∧ Matrix.det U = 1

/--
The CHSH correlation bound is mathematically realized natively by the $2\sqrt{2}$ geometric bounds of the macroscopic SU(2) spatial topology, bypassing the need for abstract Hilbert space operators.
-/
@[litlib_track "Geometric Tsirelson Bound"]
theorem §\phantomsection\label{lean:CGD.Quantum.geometricHolonomyTsirelsonBound}§CGD.Quantum.geometricHolonomyTsirelsonBound
  (A1 A2 B1 B2 : §\hyperref[lean:CGD.Foundations.SU2Group]{\detokenize{SU2Group}}§) :
  let chsh := §\hyperref[lean:CGD.Quantum.geometricBellCorrelation]{\detokenize{geometricBellCorrelation}}§ A1 B1 + §\hyperref[lean:CGD.Quantum.geometricBellCorrelation]{\detokenize{geometricBellCorrelation}}§ A1 B2 +
              §\hyperref[lean:CGD.Quantum.geometricBellCorrelation]{\detokenize{geometricBellCorrelation}}§ A2 B1 - §\hyperref[lean:CGD.Quantum.geometricBellCorrelation]{\detokenize{geometricBellCorrelation}}§ A2 B2;
  (chsh.re)^2 ≤ 8 ∧ chsh.im = 0


-- MODULE: CGD.Quantum.Holonomy.RelationalTime
--------------------------------------------------------------------
/--
Using an exact analytical witness representing the U(1) core of a topological defect
(`Complex.I • sigmaZ`), this theorem demonstrates that computing the geometric holonomy
(spatial parallel transport) is mathematically identical to applying the Schrödinger
time-evolution operator.

This proves that traversal along a defect's internal spatial phase space can physically
manifest as the passage of time, demonstrating that relational geometric phase
mathematically satisfies the role of time evolution.
-/
@[litlib_track "Witness for Relational Time Emergence"]
theorem §\phantomsection\label{lean:CGD.Quantum.Holonomy.relationalTimeEmergence}§CGD.Quantum.Holonomy.relationalTimeEmergence
  (matrixExp : Matrix (Fin 2) (Fin 2) ℂ → Matrix (Fin 2) (Fin 2) ℂ)
  (pu : §\hyperref[lean:CGD.Axioms.PhysicalUniverse]{\detokenize{PhysicalUniverse}}§)
  (t : ℝ)
  (h_core : ∀ s, (pu.toUniverse.sd_sector.val 1 (§\hyperref[lean:CGD.Quantum.straightLinePath]{\detokenize{straightLinePath}}§ s)).val = Complex.I • §\hyperref[lean:CGD.Foundations.sigmaZ]{\detokenize{sigmaZ}}§) :
  §\hyperref[lean:CGD.Quantum.macroscopicObservable]{\detokenize{macroscopicObservable}}§ (§\hyperref[lean:CGD.Quantum.holonomy]{\detokenize{holonomy}}§ matrixExp) pu.toUniverse.sd_sector.val 1 t =
  §\hyperref[lean:CGD.Quantum.Holonomy.unitaryTimeEvolution]{\detokenize{unitaryTimeEvolution}}§ matrixExp (-§\hyperref[lean:CGD.Foundations.sigmaZ]{\detokenize{sigmaZ}}§) t

/-- Standard quantum mechanical time-evolution operator U(t) = exp(-i H t) -/
@[litlib_track "Unitary Time Evolution Operator"]
noncomputable def §\phantomsection\label{lean:CGD.Quantum.Holonomy.unitaryTimeEvolution}§CGD.Quantum.Holonomy.unitaryTimeEvolution
  (matrixExp : Matrix (Fin 2) (Fin 2) ℂ → Matrix (Fin 2) (Fin 2) ℂ)
  (Hamiltonian : Matrix (Fin 2) (Fin 2) ℂ) (t : ℝ) : Matrix (Fin 2) (Fin 2) ℂ :=
  matrixExp ((-Complex.I * (t : ℂ)) • Hamiltonian)


-- MODULE: CGD.Quantum.Measurement.BornRule
--------------------------------------------------------------------
/--
The exact analytical primitive (volume function) of the Bengtsson invariant Hopf metric.
According to Litlib.Y2017.bengtsson2017geometry.Chapter03.Sec05_HopfFibration.Eq3_98,
the invariant metric volume element scales natively as sin(θ).
The strict analytical primitive (antiderivative) of sin(θ) evaluates to -cos(θ).
-/
noncomputable def §\phantomsection\label{lean:CGD.Quantum.Measurement.hopfVolumePrimitive}§CGD.Quantum.Measurement.hopfVolumePrimitive (theta : ℝ) : ℝ :=
  - Real.cos theta

/--
The Kinematic Born Rule Equivalence.
Demonstrates that in Chiral Gauge Dynamics, the phase-space volume fraction
(calculated strictly via the invariant Hopf metric primitive) bounded by the
geometric correlation between the Physical Universe's gauge state and a
macroscopic detector frame mathematically evaluates exactly to the quantum
mechanical Born rule projection.

This derives the fundamental mechanism of quantum probability directly from
the macroscopic continuous geometry of the gauge field, without requiring
wave-function collapse or arbitrary dynamic thresholds.
-/
@[litlib_track "Kinematic Born Rule Equivalence"]
theorem §\phantomsection\label{lean:CGD.Quantum.Measurement.kinematicBornRuleEquivalence}§CGD.Quantum.Measurement.kinematicBornRuleEquivalence
  (pu : §\hyperref[lean:CGD.Axioms.PhysicalUniverse]{\detokenize{PhysicalUniverse}}§)
  (evaluateBoundary : §\hyperref[lean:CGD.Axioms.Sl2cGaugeField]{\detokenize{Sl2cGaugeField}}§ → §\hyperref[lean:CGD.Foundations.SU2Group]{\detokenize{SU2Group}}§)
  (detector_frame : §\hyperref[lean:CGD.Foundations.SU2Group]{\detokenize{SU2Group}}§)
  (theta : ℝ)
  -- The physical angle is strictly defined by the SU(2) trace metric between the universe and the detector
  (h_angle : Real.cos theta = (§\hyperref[lean:CGD.Quantum.geometricBellCorrelation]{\detokenize{geometricBellCorrelation}}§ (evaluateBoundary pu.toUniverse.sd_sector) detector_frame).re) :
  -- The phase-space volume fraction of the universe...
  (§\hyperref[lean:CGD.Quantum.Measurement.hopfVolumePrimitive]{\detokenize{hopfVolumePrimitive}}§ Real.pi - §\hyperref[lean:CGD.Quantum.Measurement.hopfVolumePrimitive]{\detokenize{hopfVolumePrimitive}}§ theta) /
  (§\hyperref[lean:CGD.Quantum.Measurement.hopfVolumePrimitive]{\detokenize{hopfVolumePrimitive}}§ Real.pi - §\hyperref[lean:CGD.Quantum.Measurement.hopfVolumePrimitive]{\detokenize{hopfVolumePrimitive}}§ 0) =
  -- ...is mathematically identical to the linear projection...
  (1 + (§\hyperref[lean:CGD.Quantum.geometricBellCorrelation]{\detokenize{geometricBellCorrelation}}§ (evaluateBoundary pu.toUniverse.sd_sector) detector_frame).re) / 2 ∧
  -- ...which identically equals the quantum mechanical Born rule probability projection
  (1 + (§\hyperref[lean:CGD.Quantum.geometricBellCorrelation]{\detokenize{geometricBellCorrelation}}§ (evaluateBoundary pu.toUniverse.sd_sector) detector_frame).re) / 2 = (Real.cos (theta / 2))^2


-- MODULE: CGD.Quantum.Schroedinger
--------------------------------------------------------------------
noncomputable def §\phantomsection\label{lean:CGD.Quantum.P_minus}§CGD.Quantum.P_minus : Matrix (Fin 4) (Fin 4) Complex :=
  (1 / 2 : Complex) • (1 - §\hyperref[lean:Litlib.Math.Dirac.gamma0]{\detokenize{gamma0}}§)

noncomputable def §\phantomsection\label{lean:CGD.Quantum.P_plus}§CGD.Quantum.P_plus : Matrix (Fin 4) (Fin 4) Complex :=
  (1 / 2 : Complex) • (1 + §\hyperref[lean:Litlib.Math.Dirac.gamma0]{\detokenize{gamma0}}§)

/--
The exact algebraic chiral split of the emergent Dirac equation.
By applying the projection operators, the relativistic Dirac equation natively factors
into a coupled system for the large and small components.
-/
@[litlib_track "Algebraic Dirac Chiral Split"]
theorem §\phantomsection\label{lean:CGD.Quantum.algebraicDiracChiralSplit}§CGD.Quantum.algebraicDiracChiralSplit (dPsi : Fin 4 → §\hyperref[lean:CGD.Foundations.SpacetimePoint]{\detokenize{SpacetimePoint}}§ → Matrix (Fin 4) (Fin 4) Complex)
  (Psi : §\hyperref[lean:CGD.Foundations.SpacetimePoint]{\detokenize{SpacetimePoint}}§ → Matrix (Fin 4) (Fin 4) Complex) (m : Complex) (x : §\hyperref[lean:CGD.Foundations.SpacetimePoint]{\detokenize{SpacetimePoint}}§) :
  §\hyperref[lean:CGD.Quantum.localDiracOp]{\detokenize{localDiracOp}}§ (fun a => dPsi a x) = m • Psi x →
  let D0_mod := §\hyperref[lean:CGD.Quantum.modulatedTemporalDeriv]{\detokenize{modulatedTemporalDeriv}}§ (dPsi 0 x) (Psi x) m
  let D_space := §\hyperref[lean:CGD.Quantum.spatialDiracOp]{\detokenize{spatialDiracOp}}§ (fun mu => dPsi mu x)
  (§\hyperref[lean:CGD.Quantum.P_plus]{\detokenize{P_plus}}§ * D0_mod + §\hyperref[lean:CGD.Quantum.P_plus]{\detokenize{P_plus}}§ * §\hyperref[lean:Litlib.Math.Dirac.gamma0]{\detokenize{gamma0}}§ * D_space = 2 • m • (§\hyperref[lean:CGD.Quantum.P_plus]{\detokenize{P_plus}}§ * Psi x)) ∧
  (§\hyperref[lean:CGD.Quantum.P_minus]{\detokenize{P_minus}}§ * D0_mod + §\hyperref[lean:CGD.Quantum.P_minus]{\detokenize{P_minus}}§ * §\hyperref[lean:Litlib.Math.Dirac.gamma0]{\detokenize{gamma0}}§ * D_space = 0)

/--
An algebraic decomposition of the emergent Dirac equation.

By applying the standard chiral projection operators (P_plus, P_minus), the
relativistic equation natively factors into coupled relations for the large and
small components. This verifies that the standard algebraic structures required
for non-relativistic limits (which produce the 1/2m Schrödinger/Pauli Hamiltonian)
are natively supported by the macroscopic gauge-covariant geometry.
-/
@[litlib_track "Exact Schroedinger Reduction"]
theorem §\phantomsection\label{lean:CGD.Quantum.exactSchroedingerReduction}§CGD.Quantum.exactSchroedingerReduction (dPsi : Fin 4 → §\hyperref[lean:CGD.Foundations.SpacetimePoint]{\detokenize{SpacetimePoint}}§ → Matrix (Fin 4) (Fin 4) Complex)
  (Psi : §\hyperref[lean:CGD.Foundations.SpacetimePoint]{\detokenize{SpacetimePoint}}§ → Matrix (Fin 4) (Fin 4) Complex) (m : Complex) (x : §\hyperref[lean:CGD.Foundations.SpacetimePoint]{\detokenize{SpacetimePoint}}§) :
  §\hyperref[lean:CGD.Quantum.localDiracOp]{\detokenize{localDiracOp}}§ (fun a => dPsi a x) = m • Psi x →
  let D0_mod := §\hyperref[lean:CGD.Quantum.modulatedTemporalDeriv]{\detokenize{modulatedTemporalDeriv}}§ (dPsi 0 x) (Psi x) m
  let Psi_small := §\hyperref[lean:CGD.Quantum.P_plus]{\detokenize{P_plus}}§ * Psi x
  (2 • m • Psi_small = §\hyperref[lean:CGD.Quantum.P_plus]{\detokenize{P_plus}}§ * D0_mod + §\hyperref[lean:Litlib.Math.Dirac.gammaVec]{\detokenize{gammaVec}}§ 1 * (§\hyperref[lean:CGD.Quantum.P_minus]{\detokenize{P_minus}}§ * dPsi 1 x) + §\hyperref[lean:Litlib.Math.Dirac.gammaVec]{\detokenize{gammaVec}}§ 2 * (§\hyperref[lean:CGD.Quantum.P_minus]{\detokenize{P_minus}}§ * dPsi 2 x) + §\hyperref[lean:Litlib.Math.Dirac.gammaVec]{\detokenize{gammaVec}}§ 3 * (§\hyperref[lean:CGD.Quantum.P_minus]{\detokenize{P_minus}}§ * dPsi 3 x)) ∧
  (§\hyperref[lean:CGD.Quantum.P_minus]{\detokenize{P_minus}}§ * D0_mod = §\hyperref[lean:Litlib.Math.Dirac.gammaVec]{\detokenize{gammaVec}}§ 1 * (§\hyperref[lean:CGD.Quantum.P_plus]{\detokenize{P_plus}}§ * dPsi 1 x) + §\hyperref[lean:Litlib.Math.Dirac.gammaVec]{\detokenize{gammaVec}}§ 2 * (§\hyperref[lean:CGD.Quantum.P_plus]{\detokenize{P_plus}}§ * dPsi 2 x) + §\hyperref[lean:Litlib.Math.Dirac.gammaVec]{\detokenize{gammaVec}}§ 3 * (§\hyperref[lean:CGD.Quantum.P_plus]{\detokenize{P_plus}}§ * dPsi 3 x))

/--
The core Dirac operator evaluated in the local tangent space frame (using tangent space index 'a'
instead of coordinate index 'mu'). In curved space, dP represents the tetrad-contracted derivative.
-/
noncomputable def §\phantomsection\label{lean:CGD.Quantum.localDiracOp}§CGD.Quantum.localDiracOp (dP : Fin 4 → Matrix (Fin 4) (Fin 4) Complex) : Matrix (Fin 4) (Fin 4) Complex :=
  ∑ a, §\hyperref[lean:Litlib.Math.Dirac.gammaVec]{\detokenize{gammaVec}}§ a * dP a

noncomputable def §\phantomsection\label{lean:CGD.Quantum.modulatedTemporalDeriv}§CGD.Quantum.modulatedTemporalDeriv (dPsi0 Psi : Matrix (Fin 4) (Fin 4) Complex) (m : Complex) : Matrix (Fin 4) (Fin 4) Complex :=
  dPsi0 + m • Psi

noncomputable def §\phantomsection\label{lean:CGD.Quantum.spatialDiracOp}§CGD.Quantum.spatialDiracOp (dPsi : Fin 4 → Matrix (Fin 4) (Fin 4) Complex) : Matrix (Fin 4) (Fin 4) Complex :=
  §\hyperref[lean:Litlib.Math.Dirac.gammaVec]{\detokenize{gammaVec}}§ 1 * dPsi 1 + §\hyperref[lean:Litlib.Math.Dirac.gammaVec]{\detokenize{gammaVec}}§ 2 * dPsi 2 + §\hyperref[lean:Litlib.Math.Dirac.gammaVec]{\detokenize{gammaVec}}§ 3 * dPsi 3


-- MODULE: CGD.Quantum.Vacuum
--------------------------------------------------------------------
/--
In pure connection emergent gravity, spacetime volume is not a fixed background; it is dynamically generated by the gauge field. Therefore, the existence of macroscopic space (det g ≠ 0) mathematically forbids the trivial vacuum (A = 0). Spacetime itself requires a non-zero gauge condensate.
-/
@[litlib_track "Spacetime Vacuum Condensate"]
theorem §\phantomsection\label{lean:CGD.Quantum.kinematicVacuumCondensate}§CGD.Quantum.kinematicVacuumCondensate (pu : §\hyperref[lean:CGD.Axioms.PhysicalUniverse]{\detokenize{PhysicalUniverse}}§) (x : §\hyperref[lean:CGD.Foundations.SpacetimePoint]{\detokenize{SpacetimePoint}}§) (hx : x ∈ pu.bulk) :
  pu.toUniverse.sd_sector.val ≠ (fun _ _ => 0)


-- MODULE: CGD.Quantum.YangMills
--------------------------------------------------------------------
/--
Evaluates the chaotic non-linear self-interaction of the homogenous gauge field ansatz.
-/
@[litlib_track "Yang-Mills Chaos Bound"]
theorem §\phantomsection\label{lean:CGD.Quantum.kinematicYangMillsChaos}§CGD.Quantum.kinematicYangMillsChaos :
  ∀ (x : §\hyperref[lean:CGD.Foundations.SpacetimePoint]{\detokenize{SpacetimePoint}}§),
    Matrix.trace (⁅§\hyperref[lean:CGD.Particles.homogeneousChaosAnsatz]{\detokenize{homogeneousChaosAnsatz}}§ 1 x, §\hyperref[lean:CGD.Particles.homogeneousChaosAnsatz]{\detokenize{homogeneousChaosAnsatz}}§ 2 x⁆.val *
                  ⁅§\hyperref[lean:CGD.Particles.homogeneousChaosAnsatz]{\detokenize{homogeneousChaosAnsatz}}§ 1 x, §\hyperref[lean:CGD.Particles.homogeneousChaosAnsatz]{\detokenize{homogeneousChaosAnsatz}}§ 2 x⁆.val) =
    -8 * (x 1 : ℂ)^2 * (x 2 : ℂ)^2


-- MODULE: Litlib.Math.Dirac
--------------------------------------------------------------------
/--
Physical Interpretation:
The temporal Dirac Gamma matrix $\gamma^0$ in the Weyl/Chiral basis. Flips parity.
-/
noncomputable def §\phantomsection\label{lean:Litlib.Math.Dirac.gamma0}§Litlib.Math.Dirac.gamma0 : Matrix (Fin 4) (Fin 4) Complex :=
  let m : Matrix (Fin 2 ⊕ Fin 2) (Fin 2 ⊕ Fin 2) Complex :=
    Matrix.fromBlocks 0 1 1 0
  Matrix.reindex chiralIso chiralIso m

/--
Physical Interpretation:
The covariant 4-vector of Gamma matrices $\gamma^\mu$.
-/
noncomputable def §\phantomsection\label{lean:Litlib.Math.Dirac.gammaVec}§Litlib.Math.Dirac.gammaVec (mu : Fin 4) : Matrix (Fin 4) (Fin 4) Complex :=
  match mu.val with
  | 0 => §\hyperref[lean:Litlib.Math.Dirac.gamma0]{\detokenize{gamma0}}§
  | 1 => gammaSpatial 0
  | 2 => gammaSpatial 1
  | 3 => gammaSpatial 2
  | _ => 0


-- MODULE: Litlib.Y1951.papapetrou1951spinning.Signature
--------------------------------------------------------------------
/--
Geometric Non-Degeneracy Constraint: The macroscopic metric density determinant must be strictly non-zero.
Velocity Non-Degeneracy Constraint: The 4-velocity must not be identically zero.
Algebraic Domain Constraint: Parameterized over a generic Field `F` with Characteristic Zero (`CharZero`).
Physical Domain Binding: Papapetrou (1951) Eq (2.12): If a covariant tensor distribution T represents a single-pole 
test particle, and it is covariantly conserved with respect to the background connection 
(expressed equivalently via the inverse metric contraction), then its worldline must exactly satisfy the geodesic equation.
-/
class §\phantomsection\label{lean:Litlib.Y1951.papapetrou1951spinning.Eq2_12}§Litlib.Y1951.papapetrou1951spinning.Eq2_12§\textsuperscript{\pdftooltip{\cite{papapetrou1951spinning}}{Reference: papapetrou1951spinning}}§
    (M : Type*) [TopologicalSpace M]
    (F : Type*) [Field F] [CharZero F]
    (Metric : M → Matrix (Fin 4) (Fin 4) F)
    (invMetric : M → Matrix (Fin 4) (Fin 4) F)
    (Christoffel : M → (Fin 4 → Fin 4 → Fin 4 → F))
    (T : Fin 4 → Fin 4 → M → F)
    (§\hyperref[lean:CGD.Math.partialDeriv]{\detokenize{partialDeriv}}§ : Fin 4 → (M → F) → M → F)
    (worldline : F → M)
    (u_up : F → (Fin 4 → F))
    (du_up_ds : F → (Fin 4 → F))
    (isSinglePole : (Fin 4 → Fin 4 → M → F) → (F → M) → Prop)
    where
  metric_nondegenerate : ∀ p, (Metric p).det ≠ 0
  invMetric_prop : ∀ p, Metric p * invMetric p = 1
  u_norm_nonzero : ∀ s, ∑ μ, ∑ ν, Metric (worldline s) μ ν * u_up s μ * u_up s ν ≠ 0
  single_pole_eom : 
    (∀ x b, ∑ a : Fin 4, ∑ c : Fin 4, invMetric x a c * (
      §\hyperref[lean:CGD.Math.partialDeriv]{\detokenize{partialDeriv}}§ c (fun p => T a b p) x -
      ∑ d : Fin 4, (Christoffel x d c a * T d b x + Christoffel x d c b * T a d x)
    ) = 0) →
    isSinglePole T worldline →
    ∀ s α, du_up_ds s α + ∑ μ, ∑ ν, Christoffel (worldline s) α μ ν * u_up s μ * u_up s ν = 0


-- MODULE: Litlib.Y1956.utiyama1956invariant.Signature
--------------------------------------------------------------------
class §\phantomsection\label{lean:Litlib.Y1956.utiyama1956invariant.AppendixI_BilinearForm}§Litlib.Y1956.utiyama1956invariant.AppendixI_BilinearForm§\textsuperscript{\pdftooltip{\cite{utiyama1956invariant}}{Reference: utiyama1956invariant}}§ where
  /-- 
  Invariant Bilinear Form (Appendix I):
  Constructs the uniquely non-degenerate invariant metric (the Killing form) 
  for the Lie algebra generators. The `isLieAlgebra` constraint physically restricts 
  the domain to semi-simple representations, mathematically preventing pathological 
  central extensions and non-compact ghosts from trivializing the gauge structure.
  -/
  spans 
    (M : Type*) [Ring M] [Algebra ℂ M]
    (Trace : M → ℂ)
    (isLieAlgebra : M → Prop)
    (hNonDegenerate : ∃ x, isLieAlgebra x ∧ x ≠ 0) :
    ∀ (B : M → M → ℂ),
    (∀ c x y, isLieAlgebra x → isLieAlgebra y → B (c • x) y = c * B x y) →
    (∀ x1 x2 y, isLieAlgebra x1 → isLieAlgebra x2 → isLieAlgebra y → B (x1 + x2) y = B x1 y + B x2 y) →
    (∀ x y1 y2, isLieAlgebra x → isLieAlgebra y1 → isLieAlgebra y2 → B x (y1 + y2) = B x y1 + B x y2) →
    (∀ x y (U : Mˣ), isLieAlgebra x → isLieAlgebra y → 
      isLieAlgebra ((U : M) * x * (↑U⁻¹ : M)) → 
      isLieAlgebra ((U : M) * y * (↑U⁻¹ : M)) → 
      B ((U : M) * x * (↑U⁻¹ : M)) ((U : M) * y * (↑U⁻¹ : M)) = B x y) →
    ∃ (k : ℂ), ∀ x y, isLieAlgebra x → isLieAlgebra y → B x y = k * Trace (x * y)

class §\phantomsection\label{lean:Litlib.Y1956.utiyama1956invariant.AppendixI_Expansion}§Litlib.Y1956.utiyama1956invariant.AppendixI_Expansion§\textsuperscript{\pdftooltip{\cite{utiyama1956invariant}}{Reference: utiyama1956invariant}}§ where
  /--
  Utiyama Expansion Theorem (Appendix I):
  A gauge-invariant, quadratic (renormalizable) Lagrangian constructed from the 
  field strength tensor uniquely decomposes into a linear combination of the 
  traces of the field strength squared. This strictly binds the dynamical terms 
  to the invariant metric of the gauge group, establishing the algebraic necessity 
  of the Yang-Mills action.
  -/
  yieldsTraceExpansion 
    (M : Type*) [Ring M] [Algebra ℂ M]
    (Trace : M → ℂ)
    (isLieAlgebra : M → Prop)
    (hNonDegenerate : ∃ x, isLieAlgebra x ∧ x ≠ 0)
    (L : (Fin 4 → Fin 4 → M) → ℂ)
    (hTraceSpans : ∀ (B : M → M → ℂ),
      (∀ c x y, isLieAlgebra x → isLieAlgebra y → B (c • x) y = c * B x y) →
      (∀ x1 x2 y, isLieAlgebra x1 → isLieAlgebra x2 → isLieAlgebra y → B (x1 + x2) y = B x1 y + B x2 y) →
      (∀ x y1 y2, isLieAlgebra x → isLieAlgebra y1 → isLieAlgebra y2 → B x (y1 + y2) = B x y1 + B x y2) →
      (∀ x y (U : Mˣ), isLieAlgebra x → isLieAlgebra y → 
        isLieAlgebra ((U : M) * x * (↑U⁻¹ : M)) → 
        isLieAlgebra ((U : M) * y * (↑U⁻¹ : M)) → 
        B ((U : M) * x * (↑U⁻¹ : M)) ((U : M) * y * (↑U⁻¹ : M)) = B x y) →
      ∃ (k : ℂ), ∀ x y, isLieAlgebra x → isLieAlgebra y → B x y = k * Trace (x * y))
    (hLQuadScale : ∀ (c : ℂ) (F : Fin 4 → Fin 4 → M), 
      (∀ μ ν, isLieAlgebra (F μ ν)) → L (fun μ ν => c • F μ ν) = c^2 * L F)
    (hLQuadAdd : ∀ (F G : Fin 4 → Fin 4 → M), 
      (∀ μ ν, isLieAlgebra (F μ ν)) → (∀ μ ν, isLieAlgebra (G μ ν)) → 
      L (fun μ ν => F μ ν + G μ ν) + L (fun μ ν => F μ ν - G μ ν) = 2 * L F + 2 * L G)
    (hLGauge : ∀ (F : Fin 4 → Fin 4 → M) (U : Mˣ), 
      (∀ μ ν, isLieAlgebra (F μ ν)) → 
      (∀ μ ν, isLieAlgebra ((U : M) * F μ ν * (↑U⁻¹ : M))) → 
      L (fun μ ν => (U : M) * F μ ν * (↑U⁻¹ : M)) = L F) :
    ∃ (T : Fin 4 → Fin 4 → Fin 4 → Fin 4 → ℂ), 
      ∀ F, (∀ μ ν, isLieAlgebra (F μ ν)) → 
      L F = ∑ μ : Fin 4, ∑ ν : Fin 4, ∑ ρ : Fin 4, ∑ σ : Fin 4, T μ ν ρ σ * Trace (F μ ν * F ρ σ)


-- MODULE: Litlib.Y1965.spivak1965calculus.Chapter05.IntegrationOnChains
--------------------------------------------------------------------
class §\phantomsection\label{lean:Litlib.Y1965.spivak1965calculus.DivergenceTheoremR4Compact}§Litlib.Y1965.spivak1965calculus.DivergenceTheoremR4Compact§\textsuperscript{\pdftooltip{\cite{spivak1965calculus}}{Reference: spivak1965calculus}}§
    (f : (Fin 4 → ℝ) → (Fin 4 → ℂ))
    where
  /--
  Compactly Supported Divergence Free Constraint:
  By Stokes' Theorem, the integral of an exact differential form over a manifold 
  without boundary evaluates to zero if the form has compact support. 
  Here, we enforce continuous differentiability and explicitly require the closure 
  of the support to be compact. This rigorously ensures the Lebesgue integral of 
  the divergence over Euclidean space is mathematically well-posed and evaluates 
  to zero, precluding boundary term singularities.
  -/
  hf_smooth : ContDiff ℝ 1 f
  hf_compact : IsCompact (closure {x | f x ≠ 0})
  
  integral_div_zero :
    ∫ x : Fin 4 → ℝ, (∑ i : Fin 4, deriv (fun t => f (Function.update x i t) i) (x i)) ∂MeasureTheory.volume = 0


-- MODULE: Litlib.Y1975.belavin1975pseudoparticle.Signature
--------------------------------------------------------------------
class §\phantomsection\label{lean:Litlib.Y1975.belavin1975pseudoparticle.Eq8}§Litlib.Y1975.belavin1975pseudoparticle.Eq8§\textsuperscript{\pdftooltip{\cite{belavin1975pseudoparticle}}{Reference: belavin1975pseudoparticle}}§
  (BoundaryManifold Group : Type*) [TopologicalSpace BoundaryManifold] [TopologicalSpace Group]
  [Nonempty BoundaryManifold] [Nonempty Group]
  (isSmooth : (BoundaryManifold → Group) → Prop)
  (windingNumber : (BoundaryManifold → Group) → ℤ)
  (cartanMaurerIntegral : (BoundaryManifold → Group) → ℝ)
  [§\hyperref[lean:Litlib.Y2003.nakahara2003geometry.CartanMaurerTopology]{\detokenize{Litlib.Y2003.nakahara2003geometry.CartanMaurerTopology}}§ (BoundaryManifold → Group) isSmooth windingNumber cartanMaurerIntegral] where
  
  /-- 
  Topological Non-Triviality Constraint: Ensure the existence of at least one 
  map with a unit winding number to verify the principal bundle possesses 
  a non-trivial topological structure. 
  -/
  h_exists_degree_one : ∃ (g : BoundaryManifold → Group), windingNumber g = 1

  /-- 
  Physical Interpretation: Belavin 1975, Page 86: "Hence q is the number of 
  times the SU(2) is covered under this mapping." Any homeomorphism between 
  topologically equivalent mapping spaces preserves the absolute degree of the map.
  -/
  degree_of_homeomorph :
    ∀ (f : BoundaryManifold → Group),
    §\hyperref[lean:Litlib.Y1975.belavin1975pseudoparticle.IsHomeomorphism]{\detokenize{IsHomeomorphism}}§ f →
    windingNumber f = 1 ∨ windingNumber f = -1

/-- 
Topological Well-Posedness: A mathematically strict definition of a 
topological homeomorphism—a bijection between two topological spaces 
that is continuous and has a continuous inverse. 
-/
structure §\phantomsection\label{lean:Litlib.Y1975.belavin1975pseudoparticle.IsHomeomorphism}§Litlib.Y1975.belavin1975pseudoparticle.IsHomeomorphism§\textsuperscript{\pdftooltip{\cite{belavin1975pseudoparticle}}{Reference: belavin1975pseudoparticle}}§ {X Y : Type*} [TopologicalSpace X] [TopologicalSpace Y] (f : X → Y) : Prop where
  bij : Function.Bijective f
  cont : Continuous f
  inv_cont : ∃ (g : Y → X), Function.LeftInverse g f ∧ Function.RightInverse g f ∧ Continuous g


-- MODULE: Litlib.Y1976.rudin1976principles.Chapter09.Sec08_DerivativesOfHigherOrder
--------------------------------------------------------------------
class §\phantomsection\label{lean:Litlib.Y1976.rudin1976principles.ClairautTheoremNDimensional}§Litlib.Y1976.rudin1976principles.ClairautTheoremNDimensional§\textsuperscript{\pdftooltip{\cite{rudin1976principles}}{Reference: rudin1976principles}}§ where
  /--
  Exercise 29 (page 243): The n-dimensional, k-th order generalization of Clairaut's theorem.
  
  Rigorous Formalization Constraint: We employ Mathlib's `ContDiffOn` to assert that 
  the function is k-times continuously differentiable. Instead of defining ad-hoc 
  index-based partial derivatives (which can obscure topological boundaries), we utilize 
  Mathlib's `iteratedFDeriv`, evaluating it on a tuple of vectors `v`. The theorem 
  establishes that permuting the order of these vectors (via `fun i => v (σ i)`) 
  leaves the result invariant, serving as the rigorous, coordinate-free equivalent 
  of permuting partial derivative indices.
  -/
  symmetry_of_higher_partials
    (n m : ℕ)
    (E : Set (Fin n → ℝ))
    (hOpen : IsOpen E)
    (f : (Fin n → ℝ) → ℝ)
    (hf : ContDiffOn ℝ m f E)
    (x : Fin n → ℝ) (hx : x ∈ E)
    (v : Fin m → (Fin n → ℝ))
    (σ : Equiv.Perm (Fin m)) :
    iteratedFDeriv ℝ m f x (fun i => v (σ i)) = iteratedFDeriv ℝ m f x v


-- MODULE: Litlib.Y1976.rudin1976principles.Chapter11.LebesgueIntegral
--------------------------------------------------------------------
class §\phantomsection\label{lean:Litlib.Y1976.rudin1976principles.LeibnizIntegralRule}§Litlib.Y1976.rudin1976principles.LeibnizIntegralRule§\textsuperscript{\pdftooltip{\cite{rudin1976principles}}{Reference: rudin1976principles}}§
    (f : ℝ → (Fin 4 → ℝ) → ℂ)
    where
  /--
  Leibniz Integral Rule via Dominated Convergence (Theorem 11.28):
  
  Topological Gatekeeping (Override Applied): A naive formalization requires 
  only pointwise differentiability, which allows pathological functions to generate 
  locally unbounded derivatives that Mathlib's default integral evaluates to zero 
  (Garbage-In Exploit). To prevent this, we mathematically enforce `ContDiff ℝ 1` 
  (joint continuous differentiability) on `f` over the product space, which 
  guarantees the local boundedness of `∂f/∂t` necessary for the rigorous application 
  of the Dominated Convergence Theorem. The variation is compactly supported 
  to prevent escape to spatial infinity.
  -/
  hf_smooth : ContDiff ℝ 1 (fun p : ℝ × (Fin 4 → ℝ) => f p.1 p.2)
  hf_integrable : ∀ t, MeasureTheory.Integrable (f t) MeasureTheory.volume
  hf_variation_compact : ∃ K : Set (Fin 4 → ℝ), IsCompact K ∧ ∀ t x, x ∉ K → f t x = f 0 x
  
  leibniz_commute :
    ∀ t, deriv (fun s => ∫ x, f s x ∂MeasureTheory.volume) t = 
         ∫ x, deriv (fun s => f s x) t ∂MeasureTheory.volume


-- MODULE: Litlib.Y1983.witten1983current.Signature
--------------------------------------------------------------------
def §\phantomsection\label{lean:AsymptoticVacuumSU2}§AsymptoticVacuumSU2 (U : (Fin 3 → ℝ) → Matrix (Fin 2) (Fin 2) ℂ) : Prop :=
  ∀ ε : ℝ, ε > 0 → ∃ R : ℝ, R > 0 ∧ ∀ x : Fin 3 → ℝ, 
    (x 0)^2 + (x 1)^2 + (x 2)^2 > R^2 → 
    ∀ i j : Fin 2, Complex.normSq ((U x) i j - (1 : Matrix (Fin 2) (Fin 2) ℂ) i j) < ε^2

noncomputable def §\phantomsection\label{lean:BaryonNumberIntegrandSU2}§BaryonNumberIntegrandSU2 (U : (Fin 3 → ℝ) → Matrix (Fin 2) (Fin 2) ℂ) (x : Fin 3 → ℝ) : ℂ :=
  ∑ i : Fin 3, ∑ j : Fin 3, ∑ k : Fin 3,
    (eps3 i j k) *
    Matrix.trace (
      (U x)⁻¹ * (partialDerivSU2 i U x) *
      ((U x)⁻¹ * (partialDerivSU2 j U x)) *
      ((U x)⁻¹ * (partialDerivSU2 k U x))
    )

/--
Finkelstein-Rubinstein Spin-Statistics for SU(2) (Page 436):
While the Wess-Zumino term vanishes for SU(2), the 4th homotopy group is non-trivial: π_4(SU(2)) = Z_2.
An adiabatic 2π spatial rotation of a degree-1 SU(2) soliton traces the non-trivial loop in configuration space.
Weighting this non-trivial topological history with a factor of -1 mathematically quantizes the soliton as a fermion.
-/
class §\phantomsection\label{lean:FinkelsteinRubinsteinQuantization}§FinkelsteinRubinsteinQuantization
  [MeasureSpace (Fin 3 → ℝ)]
  (path_parity : (ℝ → (Fin 3 → ℝ) → Matrix (Fin 2) (Fin 2) ℂ) → ℕ)
  (U_0 : (Fin 3 → ℝ) → Matrix (Fin 2) (Fin 2) ℂ)
  (_h_SU2 : ∀ x, §\hyperref[lean:IsSU2]{\detokenize{IsSU2}}§ (U_0 x))
  (_h_smooth : ∀ i j : Fin 2, Differentiable ℝ (fun x => U_0 x i j))
  (_h_vacuum : §\hyperref[lean:AsymptoticVacuumSU2]{\detokenize{AsymptoticVacuumSU2}}§ U_0)
  (_h_integrable : Integrable (§\hyperref[lean:BaryonNumberIntegrandSU2]{\detokenize{BaryonNumberIntegrandSU2}}§ U_0))
  (_h_degree_one : §\hyperref[lean:baryon_number_SU2]{\detokenize{baryon_number_SU2}}§ U_0 _h_SU2 _h_smooth _h_vacuum _h_integrable = 1)
  (U_rot : ℝ → (Fin 3 → ℝ) → Matrix (Fin 2) (Fin 2) ℂ)
  (_h_rot_SU2 : ∀ t x, §\hyperref[lean:IsSU2]{\detokenize{IsSU2}}§ (U_rot t x))
  (_h_is_2pi_rot : ∀ t x, U_rot t x = U_0 (§\hyperref[lean:rotZ]{\detokenize{rotZ}}§ t x))
  (_h_fr_parity : path_parity U_rot = 1) : Prop where
  is_fermion : §\hyperref[lean:quantum_weight]{\detokenize{quantum_weight}}§ path_parity U_rot = -1

def §\phantomsection\label{lean:IsSU2}§IsSU2 (A : Matrix (Fin 2) (Fin 2) ℂ) : Prop :=
  Matrix.det A = 1 ∧ star A * A = 1

noncomputable def §\phantomsection\label{lean:baryon_number_SU2}§baryon_number_SU2 
  [MeasureSpace (Fin 3 → ℝ)]
  (U : (Fin 3 → ℝ) → Matrix (Fin 2) (Fin 2) ℂ)
  (_h_SU2 : ∀ x, §\hyperref[lean:IsSU2]{\detokenize{IsSU2}}§ (U x)) 
  (_h_smooth : ∀ i j : Fin 2, Differentiable ℝ (fun x => U x i j))
  (_h_vacuum : §\hyperref[lean:AsymptoticVacuumSU2]{\detokenize{AsymptoticVacuumSU2}}§ U)
  (_h_integrable : Integrable (§\hyperref[lean:BaryonNumberIntegrandSU2]{\detokenize{BaryonNumberIntegrandSU2}}§ U)) : ℂ :=
  (1 / (24 * Real.pi ^ 2)) * 
  ∫ (x : Fin 3 → ℝ), §\hyperref[lean:BaryonNumberIntegrandSU2]{\detokenize{BaryonNumberIntegrandSU2}}§ U x

/-- 
Evaluates the quantum phase amplitude derived from the topological parity 
(the mapping class in π_4) of a configuration path.
-/
noncomputable def §\phantomsection\label{lean:quantum_weight}§quantum_weight (path_parity : (ℝ → (Fin 3 → ℝ) → Matrix (Fin 2) (Fin 2) ℂ) → ℕ)
  (U_path : ℝ → (Fin 3 → ℝ) → Matrix (Fin 2) (Fin 2) ℂ) : ℂ :=
  (-1 : ℂ) ^ (path_parity U_path)

/-- 
Explicit SO(3) spatial rotation matrix around the Z-axis, parameterized by an angle t. 
This continuous coordinate transformation is strictly required to evaluate the 
spatial rotation of a configuration space field.
-/
noncomputable def §\phantomsection\label{lean:rotZ}§rotZ (t : ℝ) (x : Fin 3 → ℝ) : Fin 3 → ℝ :=
  fun i =>
    if i.val = 0 then (Real.cos t) * x 0 - (Real.sin t) * x 1
    else if i.val = 1 then (Real.sin t) * x 0 + (Real.cos t) * x 1
    else x 2


-- MODULE: Litlib.Y1984.urbantke1984integrability.Signature
--------------------------------------------------------------------
class §\phantomsection\label{lean:Litlib.Y1984.urbantke1984integrability.Eq10_Symmetry}§Litlib.Y1984.urbantke1984integrability.Eq10_Symmetry§\textsuperscript{\pdftooltip{\cite{urbantke1984integrability}}{Reference: urbantke1984integrability}}§ where
  /--
  Urbantke Metric Symmetry Identity:
  Demonstrates that the Urbantke quasimetric g_μν, when constructed from an 
  antisymmetric curvature tensor F and its dual, is inherently symmetric (g_μν = g_νμ).
  This algebraic property is a foundational requirement for g_μν to act as a 
  valid pseudo-Riemannian conformal metric.
  -/
  symmetricQuasimetric
    (F : Fin 3 → Fin 4 → Fin 4 → ℂ)
    (F_dual : Fin 3 → Fin 4 → Fin 4 → ℂ)
    (epsilon3 : Fin 3 → Fin 3 → Fin 3 → ℂ)
    (g : Fin 4 → Fin 4 → ℂ)
    (hF_anti : ∀ a μ ν, F a μ ν = - F a ν μ)
    (hF_dual_anti : ∀ a μ ν, F_dual a μ ν = - F_dual a ν μ)
    (hepsilon3_anti : ∀ a b c, epsilon3 a b c = - epsilon3 b a c ∧ epsilon3 a b c = - epsilon3 a c b)
    (hg_def : ∀ μ ν, g μ ν = (-1 / 6 : ℂ) * Finset.sum Finset.univ (fun a => Finset.sum Finset.univ (fun b => Finset.sum Finset.univ (fun c => 
      Finset.sum Finset.univ (fun α => Finset.sum Finset.univ (fun β => 
        epsilon3 a c b * F a μ α * F_dual c α β * F b β ν)))))) :
    ∀ μ ν, g μ ν = g ν μ


-- MODULE: Litlib.Y1991.capovilla1991pure.Signature
--------------------------------------------------------------------
class §\phantomsection\label{lean:Litlib.Y1991.capovilla1991pure.Theorem_Volume_Element_Identity}§Litlib.Y1991.capovilla1991pure.Theorem_Volume_Element_Identity§\textsuperscript{\pdftooltip{\cite{capovilla1991pure}}{Reference: capovilla1991pure}}§
    (Spacetime : Type*) [TopologicalSpace Spacetime]
    (sqrt_g : Spacetime → ℂ)
    (mu : Spacetime → ℂ)
    (eta : Spacetime → ℂ)
    (Psi_3x3 : Spacetime → Matrix (Fin 3) (Fin 3) ℂ)
    (M_3x3 : Spacetime → Matrix (Fin 3) (Fin 3) ℂ) where
  /-- Field Non-Degeneracy Constraint: Precludes division by zero and ensures metric non-degeneracy. -/
  h_sqrt_g_nondeg : ∀ x, sqrt_g x ≠ 0
  h_mu_nondeg : ∀ x, mu x ≠ 0
  h_M_nondeg : ∀ x, Matrix.det (M_3x3 x) ≠ 0
  h_Psi_nondeg : ∀ x, Matrix.det (Psi_3x3 x) ≠ 0
  
  -- Eq 2.16
  eq2_16_holds : ∀ x, (eta x)^2 = mu x / Matrix.det (M_3x3 x)
  -- Eq 2.21
  eq2_21_holds : ∀ x, eta x = (sqrt_g x * Matrix.det (Psi_3x3 x))⁻¹
  -- Eq 2.11 taking determinant: det(Psi) = mu^(-3/2) * det(M^(1/2)) => det(Psi)^2 = mu^(-3) * det(M)
  eq2_11_det_holds : ∀ x, (Matrix.det (Psi_3x3 x))^2 = (mu x)⁻¹ * (mu x)⁻¹ * (mu x)⁻¹ * Matrix.det (M_3x3 x)
  
  /-- 
  Physical Volume Density Evaluation:
  The text on page 67 implicitly defines and utilizes this identity, showing that the 
  Lagrange multiplier μ enforcing tracelessness in the matrix action natively encodes 
  the emergent square root metric determinant \sqrt{g}. Computed algebraically 
  from comparing the evaluation of η across equations 2.16 and 2.21.
  -/
  volume_element_identity : ∀ x, (sqrt_g x)^2 = (mu x)^2


-- MODULE: Litlib.Y2000.hall2000elementary.Signature
--------------------------------------------------------------------
class §\phantomsection\label{lean:Litlib.Y2000.hall2000elementary.DerivativeExponential}§Litlib.Y2000.hall2000elementary.DerivativeExponential§\textsuperscript{\pdftooltip{\cite{hall2000elementary}}{Reference: hall2000elementary}}§ 
    (n : Type*) [Fintype n] [DecidableEq n]
    [NormedAddCommGroup (Matrix n n ℂ)] [NormedSpace ℝ (Matrix n n ℂ)]
    (exp : Matrix n n ℂ → Matrix n n ℂ) where
  hIsExp : ∀ X, Tendsto (fun m : ℕ => ∑ k ∈ Finset.range m, (1 / (Nat.factorial k : ℂ)) • X^k) atTop (nhds  (exp X))
  smooth_exp : 
    ∀ X : Matrix n n ℂ, ContDiff ℝ ⊤ (fun (t : ℝ) => exp ((t : ℂ) • X))
    
  /--
  # Proposition: Derivative of the Exponential Mapping
  **Literature:** Hall 2000, Proposition 3.4, Page 29.
  **Physical Interpretation:** The matrix exponential maps a tangent vector (Lie algebra generator) to a continuous path in the Lie group. The temporal derivative of this path evaluated at the identity uniquely recovers the generator $X$. This ensures the rigorous binding between macroscopic group transformations and infinitesimal generators.
  **Mathematical Boundaries:** Enforces that the mapping from $\mathbb{R}$ to the space of matrices is smooth (`ContDiff`). The derivative applies strictly under the scalar action of the real line.
  -/
  deriv_exp : 
    ∀ (X : Matrix n n ℂ) (t : ℝ), 
      deriv (fun (s : ℝ) => exp ((s : ℂ) • X)) t = X * exp ((t : ℂ) • X) ∧ 
      deriv (fun (s : ℝ) => exp ((s : ℂ) • X)) t = exp ((t : ℂ) • X) * X
      
  deriv_exp_zero : 
    ∀ X : Matrix n n ℂ, deriv (fun (s : ℝ) => exp ((s : ℂ) • X)) 0 = X


-- MODULE: Litlib.Y2003.nakahara2003geometry.Chapter07.Sec04_LeviCivita
--------------------------------------------------------------------
class §\phantomsection\label{lean:Litlib.Y2003.nakahara2003geometry.Theorem_ContractedBianchi}§Litlib.Y2003.nakahara2003geometry.Theorem_ContractedBianchi§\textsuperscript{\pdftooltip{\cite{nakahara2003geometry}}{Reference: nakahara2003geometry}}§ 
    (Point Index : Type _) [Fintype Index] [DecidableEq Index] [Nonempty Index] [Nonempty Point]
    (isSmooth : (Point → ℂ) → Prop)
    (§\hyperref[lean:CGD.Math.partialDeriv]{\detokenize{partialDeriv}}§ : Index → (Point → ℂ) → Point → ℂ) where
  
  -- The fundamental rules of calculus required for the geometry
  derivCommute : ∀ mu nu f x, isSmooth f → isSmooth (§\hyperref[lean:CGD.Math.partialDeriv]{\detokenize{partialDeriv}}§ mu f) → isSmooth (§\hyperref[lean:CGD.Math.partialDeriv]{\detokenize{partialDeriv}}§ nu f) →
    §\hyperref[lean:CGD.Math.partialDeriv]{\detokenize{partialDeriv}}§ mu (fun p => §\hyperref[lean:CGD.Math.partialDeriv]{\detokenize{partialDeriv}}§ nu f p) x = §\hyperref[lean:CGD.Math.partialDeriv]{\detokenize{partialDeriv}}§ nu (fun p => §\hyperref[lean:CGD.Math.partialDeriv]{\detokenize{partialDeriv}}§ mu f p) x
  derivLeibniz : ∀ mu f1 f2 x, isSmooth f1 → isSmooth f2 →
    §\hyperref[lean:CGD.Math.partialDeriv]{\detokenize{partialDeriv}}§ mu (fun p => f1 p * f2 p) x = §\hyperref[lean:CGD.Math.partialDeriv]{\detokenize{partialDeriv}}§ mu f1 x * f2 x + f1 x * §\hyperref[lean:CGD.Math.partialDeriv]{\detokenize{partialDeriv}}§ mu f2 x

  -- The General Theorem: If a metric satisfies the Levi-Civita premises, it satisfies Bianchi.
  all_metrics_satisfy_bianchi :
    ∀ (g g_inv : Index → Index → Point → ℂ)
      (§\hyperref[lean:CGD.Gravity.christoffel]{\detokenize{christoffel}}§ : Index → Index → Index → Point → ℂ)
      (ricci : Index → Index → Point → ℂ)
      (scalarCurv : Point → ℂ)
      (G : Index → Index → Point → ℂ),
      -- PREMISES (Algebraic):
      (∀ x i j, g i j x = g j i x) →
      (∀ x i j, g_inv i j x = g_inv j i x) →
      (∀ x i j, (∑ k, g i k x * g_inv k j x) = if i = j then 1 else 0) →
      (∀ x rho mu nu, §\hyperref[lean:CGD.Gravity.christoffel]{\detokenize{christoffel}}§ rho mu nu x = (1/2 : ℂ) * ∑ sigma, g_inv rho sigma x * (
        §\hyperref[lean:CGD.Math.partialDeriv]{\detokenize{partialDeriv}}§ mu (fun p => g sigma nu p) x + §\hyperref[lean:CGD.Math.partialDeriv]{\detokenize{partialDeriv}}§ nu (fun p => g mu sigma p) x - §\hyperref[lean:CGD.Math.partialDeriv]{\detokenize{partialDeriv}}§ sigma (fun p => g mu nu p) x)) →
      (∀ x mu nu, ricci mu nu x = ∑ rho, (
        §\hyperref[lean:CGD.Math.partialDeriv]{\detokenize{partialDeriv}}§ rho (fun p => §\hyperref[lean:CGD.Gravity.christoffel]{\detokenize{christoffel}}§ rho mu nu p) x - §\hyperref[lean:CGD.Math.partialDeriv]{\detokenize{partialDeriv}}§ nu (fun p => §\hyperref[lean:CGD.Gravity.christoffel]{\detokenize{christoffel}}§ rho mu rho p) x + 
        ∑ lambda, (§\hyperref[lean:CGD.Gravity.christoffel]{\detokenize{christoffel}}§ rho lambda rho x * §\hyperref[lean:CGD.Gravity.christoffel]{\detokenize{christoffel}}§ lambda mu nu x - §\hyperref[lean:CGD.Gravity.christoffel]{\detokenize{christoffel}}§ rho lambda nu x * §\hyperref[lean:CGD.Gravity.christoffel]{\detokenize{christoffel}}§ lambda mu rho x))) →
      (∀ x, scalarCurv x = ∑ alpha, ∑ beta, g_inv alpha beta x * ricci alpha beta x) →
      (∀ x mu nu, G mu nu x = ricci mu nu x - (1/2 : ℂ) * g mu nu x * scalarCurv x) →
      -- PREMISES (Differential Smoothness):
      (∀ i j, isSmooth (fun p => g i j p)) →
      (∀ i j, isSmooth (fun p => g_inv i j p)) →
      (∀ rho mu nu, isSmooth (fun p => §\hyperref[lean:CGD.Gravity.christoffel]{\detokenize{christoffel}}§ rho mu nu p)) →
      -- CONCLUSION:
      ∀ (nu : Index) (x : Point),
        ∑ mu : Index, ∑ alpha : Index, g_inv mu alpha x * (
          §\hyperref[lean:CGD.Math.partialDeriv]{\detokenize{partialDeriv}}§ alpha (fun p => G mu nu p) x -
          ∑ lambda : Index, (§\hyperref[lean:CGD.Gravity.christoffel]{\detokenize{christoffel}}§ lambda alpha mu x * G lambda nu x + §\hyperref[lean:CGD.Gravity.christoffel]{\detokenize{christoffel}}§ lambda alpha nu x * G mu lambda x)
        ) = 0


-- MODULE: Litlib.Y2003.nakahara2003geometry.Chapter10.Sec05_GaugeTheories
--------------------------------------------------------------------
class §\phantomsection\label{lean:Litlib.Y2003.nakahara2003geometry.CartanMaurerTopology}§Litlib.Y2003.nakahara2003geometry.CartanMaurerTopology§\textsuperscript{\pdftooltip{\cite{nakahara2003geometry}}{Reference: nakahara2003geometry}}§ 
    (GroupMap : Type _) [TopologicalSpace GroupMap] [Nonempty GroupMap]
    (isSmooth : GroupMap → Prop)
    (windingNumber : GroupMap → ℤ)
    (cartanMaurerIntegral : GroupMap → ℝ) where
  degreeTheorem :
    ∀ g : GroupMap, isSmooth g → cartanMaurerIntegral g = (windingNumber g : ℝ)
  homotopyInvariance
    (H : ℝ → GroupMap)
    (hCont : Continuous H) :
    ∀ t1 t2 : ℝ, windingNumber (H t1) = windingNumber (H t2)


-- MODULE: Litlib.Y2010.wald2010general.AppendixD.Conformal
--------------------------------------------------------------------
/-- Conformal Ricci Transformation (Wald Eq D.8 for n=4).
This explicitly maps the Ricci tensor of a conformally transformed metric to the 
original Ricci tensor using the fully expanded covariant derivatives of the conformal 
factor. Explicit index expansion (Fin 4) is utilized to prevent opaque function 
exploits and ensure dimensional rigor. -/
class §\phantomsection\label{lean:Litlib.Y2010.wald2010general.ConformalRicciTransformation}§Litlib.Y2010.wald2010general.ConformalRicciTransformation§\textsuperscript{\pdftooltip{\cite{wald2010general}}{Reference: wald2010general}}§
    (M : Type _)
    (g g_tilde g_inv : M → Fin 4 → Fin 4 → ℝ)
    (Ricci Ricci_tilde : M → Fin 4 → Fin 4 → ℝ)
    (Omega : M → ℝ)
    (nabla : (M → ℝ) → M → Fin 4 → ℝ)
    (nabla_nabla : (M → ℝ) → M → Fin 4 → Fin 4 → ℝ)
    (isLeviCivitaRicci : (M → (Fin 4 → Fin 4 → ℝ)) → (M → (Fin 4 → Fin 4 → ℝ)) → Prop)
    where
  /-- Topological Gatekeeping (Conformal Factor): 
  The conformal factor `Omega` must be strictly positive everywhere on the manifold. 
  This mathematically prevents evaluation of undefined limits for the natural 
  logarithm `ln(Omega)` present in the connection difference tensor. -/
  omega_pos : ∀ p : M, 0 < Omega p

  /-- Inverse Metric Constraint: `g_inv` is explicitly verified as the strictly 
  valid matrix inverse of the original metric `g` via Kronecker delta contraction. 
  This provides the explicit algebraic structure required for downstream theorem proving. -/
  valid_inverse : ∀ (p : M) (a c : Fin 4), 
    (∑ b : Fin 4, g p a b * g_inv p b c) = if a = c then 1 else 0

  /-- Geometric Connection Constraint: 
  Both Ricci tensors must be explicitly bound to the torsion-free Levi-Civita 
  connections of their respective metrics. -/
  levi_civita_bound : isLeviCivitaRicci g Ricci
  levi_civita_bound_tilde : isLeviCivitaRicci g_tilde Ricci_tilde

  /-- Conformal Ricci Transformation (Explicit 4D Expansion):
  The exact algebraic relationship between the transformed Ricci tensor and the 
  original Ricci tensor, expanded for n=4 dimensions. The derivative operators 
  are explicitly applied to the natural logarithm of the conformal factor to 
  prevent decoupled type signature exploits. -/
  eq_D8 : ∀ (p : M) (a c : Fin 4), 
    Ricci_tilde p a c = Ricci p a c 
      - 2 * nabla_nabla (fun x => Real.log (Omega x)) p a c 
      - g p a c * (∑ d : Fin 4, ∑ e : Fin 4, g_inv p d e * nabla_nabla (fun x => Real.log (Omega x)) p d e) 
      + 2 * nabla (fun x => Real.log (Omega x)) p a * nabla (fun x => Real.log (Omega x)) p c 
      - 2 * g p a c * (∑ d : Fin 4, ∑ e : Fin 4, g_inv p d e * nabla (fun x => Real.log (Omega x)) p d * nabla (fun x => Real.log (Omega x)) p e)


-- MODULE: Litlib.Y2011.krasnov2011plebanski.Signature
--------------------------------------------------------------------
/--
Physical Interpretation: Formulates the vacuum Einstein equations in the Plebański formalism. The trace of the self-dual curvature matrix is strictly proportional to the cosmological constant, and its anti-self-dual part strictly vanishes.
Mathematical Boundaries: This replaces the traditional differential, tensorial Ricci-flatness conditions with a purely algebraic constraint on the curvature components in the self-dual basis.
-/
class §\phantomsection\label{lean:Litlib.Y2011.krasnov2011plebanski.Eq15}§Litlib.Y2011.krasnov2011plebanski.Eq15§\textsuperscript{\pdftooltip{\cite{krasnov2011plebanski}}{Reference: krasnov2011plebanski}}§
    (plebanski_vacuum : ℂ → (Fin 3 → Fin 3 → ℂ) → (Fin 3 → Fin 3 → ℂ) → Prop)
    where
  plebanski_vacuum_iff : ∀ (Lambda : ℂ) (F_ij F_bar_ij : Fin 3 → Fin 3 → ℂ),
    plebanski_vacuum Lambda F_ij F_bar_ij ↔ 
    ((∑ i : Fin 3, F_ij i i) = -Lambda ∧ (∀ i j, F_bar_ij i j = 0))

/--
Physical Interpretation: Defines the projection of the trace-free macroscopic stress-energy tensor onto the mixed self-dual/anti-self-dual basis. This term serves as the source coupling matter to the gravitational field in the Plebański formalism.
Mathematical Boundaries: The construction explicitly requires raising the indices of the anti-self-dual basis forms (`Sigma_bar`) using the inverse background metric (`g_inv`), strictly binding the definition to domains where the metric is non-degenerate.
-/
class §\phantomsection\label{lean:Litlib.Y2011.krasnov2011plebanski.Eq16}§Litlib.Y2011.krasnov2011plebanski.Eq16§\textsuperscript{\pdftooltip{\cite{krasnov2011plebanski}}{Reference: krasnov2011plebanski}}§
    (Sigma : Fin 3 → Fin 4 → Fin 4 → ℂ)
    (Sigma_bar : Fin 3 → Fin 4 → Fin 4 → ℂ)
    (g_inv : Fin 4 → Fin 4 → ℂ)
    (T_tilde : Fin 4 → Fin 4 → ℂ)
    (T_ij : Fin 3 → Fin 3 → ℂ)
    where
  T_ij_def : ∀ (i j : Fin 3),
    T_ij i j = 
      ∑ mu : Fin 4, ∑ nu : Fin 4, ∑ rho : Fin 4, ∑ alpha : Fin 4, ∑ beta : Fin 4,
        T_tilde rho mu * 
        Sigma i nu rho * 
        g_inv mu alpha * 
        g_inv nu beta * 
        Sigma_bar j alpha beta

/--
Physical Interpretation: The full non-vacuum Einstein equations coupled to macroscopic matter in the Plebański formulation. The trace of the self-dual curvature is determined by the cosmological constant and the trace of the stress-energy tensor. The anti-self-dual part is proportional to the trace-free stress-energy tensor.
Mathematical Boundaries: By coupling the curvature to the stress-energy components `T` and `T_ij`, these equations are strictly bound to domains where macroscopic matter fields are well-defined.
-/
class §\phantomsection\label{lean:Litlib.Y2011.krasnov2011plebanski.Eq17}§Litlib.Y2011.krasnov2011plebanski.Eq17§\textsuperscript{\pdftooltip{\cite{krasnov2011plebanski}}{Reference: krasnov2011plebanski}}§
    (Lambda : ℂ)
    (G : ℂ)
    (F_ij : Fin 3 → Fin 3 → ℂ)
    (F_bar_ij : Fin 3 → Fin 3 → ℂ)
    (T : ℂ)
    (T_ij : Fin 3 → Fin 3 → ℂ)
    (plebanski_matter_eqs : Prop)
    where
  einstein_eqs_iff : plebanski_matter_eqs ↔ 
    ((∑ i : Fin 3, F_ij i i) = -Lambda - 2 * (Real.pi : ℂ) * G * T ∧ 
     (∀ i j, F_bar_ij i j = -2 * (Real.pi : ℂ) * G * T_ij i j))

/--
Physical Interpretation: Establishes the bridge between the Plebański formulation and the standard metric formulation of general relativity (Page 6). It mathematically proves that for a Levi-Civita (torsion-free) connection, the algebraic self-dual curvature conditions (vanishing anti-self-dual part and trace proportional to the cosmological constant) are rigorously equivalent to the tensorial trace-reversed vacuum Einstein field equations ($R_{ab} = \Lambda g_{ab}$).
Mathematical Boundaries: The equivalence strictly requires a non-degenerate spacetime metric `g` with a well-defined inverse (`g_inv`) to prevent dimensional collapse. This constraint, alongside the explicit requirement that the connection is Levi-Civita (`isLeviCivitaRicci`), securely binds the theorem to standard Riemannian geometries where the isomorphism between the self-dual curvature and the Riemann curvature physically holds.
-/
class §\phantomsection\label{lean:Litlib.Y2011.krasnov2011plebanski.PlebanskiToEinsteinEquivalence}§Litlib.Y2011.krasnov2011plebanski.PlebanskiToEinsteinEquivalence§\textsuperscript{\pdftooltip{\cite{krasnov2011plebanski}}{Reference: krasnov2011plebanski}}§
    (g : Fin 4 → Fin 4 → ℝ)
    (g_inv : Fin 4 → Fin 4 → ℝ)
    (Ricci : Fin 4 → Fin 4 → ℝ)
    (Lambda : ℝ)
    (F_ij F_bar_ij : Fin 3 → Fin 3 → ℂ)
    (plebanski_vacuum : ℂ → (Fin 3 → Fin 3 → ℂ) → (Fin 3 → Fin 3 → ℂ) → Prop)
    (isLeviCivitaRicci : (Fin 4 → Fin 4 → ℝ) → (Fin 4 → Fin 4 → ℝ) → Prop)
    where
  equivalence_iff : 
    (∀ mu nu, (∑ alpha : Fin 4, g mu alpha * g_inv alpha nu) = if mu = nu then 1 else 0) → 
    isLeviCivitaRicci g Ricci → 
    (plebanski_vacuum (Lambda : ℂ) F_ij F_bar_ij ↔ (∀ a b, Ricci a b = Lambda * g a b))

\end{minted}
